% Traduction et production de textes liés à la citoyenneté et d'interviews divers — Allemand 1ère A — Cours signé OnBuch+
% Programme officiel MINESEC. Compiler avec : tectonic ALA27-traduction-et-production-de-textes-lies-a-la-citoyennete-et-.tex
\newcommand{\DOCMATIERE}{Allemand}
\newcommand{\DOCNIVEAU}{1ère A}
\newcommand{\DOCMODULE}{Chapitre 5 : Citoyenneté}
\newcommand{\DOCLECON}{ALA27}
\newcommand{\DOCTITRE}{Traduction et production de textes liés à la citoyenneté et d'interviews divers}
% OnBuch+ — préambule « cours signés » ENRICHI (compilé avec tectonic / XeLaTeX)
% Système visuel « Pop » d'OnBuch+ : fond crème (jamais blanc), bordures encre
% épaisses, ombres « dures » décalées, titres Archivo Black en capitales.
% Version MATHÉMATIQUES : graphes (pgfplots/tikz), boîtes pédagogiques
% (définition, propriété, méthode, attention, le savais-tu, exercice résolu,
% exercice, TP, bilan), page de garde et signature OnBuch+ ancrée en bas.
%
% Macros attendues dans main.tex AVANT \input{preamble} :
%   \DOCMATIERE  \DOCNIVEAU  \DOCMODULE  \DOCLECON (numéro)  \DOCTITRE
\documentclass[11pt]{article}

\usepackage{fontspec}
\usepackage[ngerman,french]{babel} % français = langue principale ; l'allemand via \all{..} / textebox
\usepackage[a4paper,margin=2cm,top=3.4cm,bottom=2.8cm,headheight=1.4cm,footskip=1.3cm]{geometry}
\usepackage{amsmath,amssymb,amsfonts}
\usepackage{mathtools} % \xmapsto, \coloneqq... souvent utilisés par les modèles
\usepackage{cancel} % simplifications barrées (bilans de forces)
% \abs{z} (module, valeur absolue) et \norm{u} (norme), utilisés par les modèles
\providecommand{\abs}{}\let\abs\relax\DeclarePairedDelimiter{\abs}{\lvert}{\rvert}
\providecommand{\norm}{}\let\norm\relax\DeclarePairedDelimiter{\norm}{\lVert}{\rVert}
\usepackage[table]{xcolor} % [table] : fournit aussi \rowcolors (alternance de lignes)
\usepackage{pagecolor}
\usepackage{tikz}
\usetikzlibrary{arrows.meta,positioning,calc,shapes.geometric,shapes.misc,shapes.arrows,decorations.pathmorphing,decorations.pathreplacing,decorations.markings,patterns,shadows,mindmap,trees,fit,backgrounds,matrix,intersections,angles,quotes}
\usepackage{pgfplots}
\usepackage[version=4]{mhchem} % équations nucléaires \ce{^{235}_{92}U}
\usepackage[european,siunitx]{circuitikz} % schémas électriques
\pgfplotsset{compat=1.18}
\usepgfplotslibrary{fillbetween,groupplots} % aires entre courbes (calcul intégral)
\usepackage{enumitem}
\usepackage{fancyhdr}
% [most] charge toutes les bibliothèques tcolorbox (listings, minted, raster,
% documentation...) et gonfle inutilement la mémoire TeX sur un document long
% (~300 boîtes) : on ne charge que ce qui est réellement utilisé.
\usepackage[skins,breakable]{tcolorbox}
\usepackage{graphicx}
\usepackage{colortbl}
\usepackage{booktabs}
\usepackage{tabularx}
\usepackage{diagbox} % case d'en-tête diagonale des tableaux à double entrée
% Colonnes extensibles centrée / gauche / droite (C, L, R), souvent utilisées par les modèles
\newcolumntype{C}{>{\centering\arraybackslash}X}
\newcolumntype{L}{>{\raggedright\arraybackslash}X}
\newcolumntype{R}{>{\raggedleft\arraybackslash}X}
\usepackage{array}
\usepackage{multicol}
\usepackage{multirow}
\usepackage{siunitx}
% parse-numbers=false : accepte aussi \SI{2,5 \times 10^{-3}}{...}
\sisetup{locale=FR,output-decimal-marker={,},group-minimum-digits=4,parse-numbers=false}
\DeclareSIUnit\ohm{\ensuremath{\symup{\Omega}}} % U+2126 absent de la police
\DeclareSIUnit\division{div} % réglages d'oscilloscope (ms/div, V/div)
\DeclareSIUnit\tr{tr} % tours (vitesse de rotation, ex. \SI{1500}{\tr\per\minute})
\DeclareSIUnit\molar{M} % concentration molaire (ex. \SI{3}{\molar}, \SI{1}{\micro\molar})
\DeclareSIUnit\an{an} % année (le modèle confond parfois avec \year, un compteur TeX) — ex. \SI{5}{\kilo\meter\per\an}
\DeclareSIUnit\FCFA{FCFA} % franc CFA (ex. \SI{500}{\FCFA\per\kilo\gram})
\DeclareSIUnit\degreeBrix{\textdegree Brix} % teneur en sucre d'un sirop/jus (réfractométrie)
\DeclareSIUnit\month{mois} % le modèle utilise parfois \month, un compteur TeX, comme unité de durée
\usepackage{qrcode}
% \degree : commande fréquemment utilisée par les modèles pour un angle ou une
% température en dehors de \SI{}{} (non fournie par les paquets chargés).
\providecommand{\degree}{^{\circ}}
% \qed (fin de démonstration), utilisé par les modèles sans amsthm
\providecommand{\qed}{\hfill$\square$}
% \barre : double barre des valeurs interdites dans un tableau de variations (tkz-tab)
\providecommand{\barre}{\ensuremath{\|}}
% environnement proof (amsthm non chargé) utilisé par les modèles
\ifdefined\proof\else\newenvironment{proof}[1][Démonstration]{\par\noindent\textit{#1.}\ }{\qed\par}\fi
\usepackage{lastpage}
\usepackage{hyperref}
\hypersetup{hidelinks}

% ============================================================
% Palette « Pop » OnBuch+ — mêmes hex que la classe P (plus_ui.dart)
% ============================================================
\definecolor{poporange}{HTML}{F59321}
\definecolor{popdark}{HTML}{E07A0C}
\definecolor{popcream}{HTML}{FFF6E8}
\definecolor{popink}{HTML}{1C1714}
\definecolor{poppurple}{HTML}{7A5AE0}
\definecolor{popgreen}{HTML}{2E9E6A}
\definecolor{poppink}{HTML}{E0457B}
\definecolor{popblue}{HTML}{2F7DE1}
\definecolor{popgold}{HTML}{C98A12}
\definecolor{popmuted}{HTML}{8A7F76}
\definecolor{popline}{HTML}{EADFD2}
\definecolor{popgray}{HTML}{8A7F76}
\colorlet{popgrey}{popgray}
% Alias tolérants (le modèle nomme parfois les couleurs différemment)
\colorlet{popwhite}{white}
\colorlet{popblack}{popink}
\colorlet{popred}{poppink}
\colorlet{popyellow}{popgold}
% Teintes claires (fonds de tableaux/encadrés) — le modèle nomme parfois
% les couleurs avec un suffixe L (ex. popblueL) sans les avoir définies.
\colorlet{poporangeL}{poporange!20}
\colorlet{popdarkL}{popdark!20}
\colorlet{poppurpleL}{poppurple!20}
\colorlet{popgreenL}{popgreen!20}
\colorlet{poppinkL}{poppink!20}
\colorlet{popblueL}{popblue!20}
\colorlet{popgoldL}{popgold!20}
\colorlet{popredL}{popred!20}
\colorlet{popgrayL}{popgray!20}
% Teintes claires pour fonds de tableaux / zones
\colorlet{popblueL}{popblue!10!white}
\colorlet{poporangeL}{poporange!14!white}
\colorlet{popgreenL}{popgreen!12!white}
\colorlet{poppurpleL}{poppurple!12!white}
\colorlet{poppinkL}{poppink!12!white}
\colorlet{popgoldL}{popgold!14!white}

\pagecolor{popcream}

% ============================================================
% Typographie
% ============================================================
\defaultfontfeatures{Ligatures=TeX}
\setmainfont{PlusJakartaSans-Medium.ttf}[Path=../../../fonts/,BoldFont=PlusJakartaSans-Bold.ttf]
% Maths en sans-serif (Fira Math) avec chiffres et lettres droites de Plus
% Jakarta, pour que les formules se fondent dans le texte.
\usepackage[math-style=ISO,bold-style=ISO]{unicode-math}
\setmathfont{FiraMath-Regular.otf}[Path=../../../fonts/]
\setmathfont{PlusJakartaSans-Medium.ttf}[Path=../../../fonts/,range={up/num,up/{Latin,latin},\mathup}]
\usepackage{icomma} % virgule décimale sans espace en mode math
% Caractères mathématiques Unicode absents des polices de texte (Plus Jakarta,
% Archivo Black) : rendus par leur équivalent mathématique, en texte comme en
% formule — sinon ils s'affichent en carré vide (ex. « Arithmétique dans ℤ »).
\usepackage{newunicodechar}
\newunicodechar{ℝ}{\ensuremath{\mathbb{R}}}
\newunicodechar{ℤ}{\ensuremath{\mathbb{Z}}}
\newunicodechar{ℂ}{\ensuremath{\mathbb{C}}}
\newunicodechar{ℕ}{\ensuremath{\mathbb{N}}}
\newunicodechar{ℚ}{\ensuremath{\mathbb{Q}}}
\newunicodechar{𝔻}{\ensuremath{\mathbb{D}}}
\newunicodechar{α}{\ensuremath{\alpha}}
\newunicodechar{β}{\ensuremath{\beta}}
\newunicodechar{γ}{\ensuremath{\gamma}}
\newunicodechar{δ}{\ensuremath{\delta}}
\newunicodechar{ε}{\ensuremath{\varepsilon}}
\newunicodechar{θ}{\ensuremath{\theta}}
\newunicodechar{λ}{\ensuremath{\lambda}}
\newunicodechar{μ}{\ensuremath{\mu}}
\newunicodechar{σ}{\ensuremath{\sigma}}
\newunicodechar{τ}{\ensuremath{\tau}}
\newunicodechar{φ}{\ensuremath{\varphi}}
\newunicodechar{ψ}{\ensuremath{\psi}}
\newunicodechar{ω}{\ensuremath{\omega}}
\newunicodechar{Σ}{\ensuremath{\Sigma}}
% majuscules produites par \MakeUppercase dans les titres (α-aminés -> Α)
\newunicodechar{Α}{\ensuremath{\alpha}}
\newunicodechar{Β}{\ensuremath{\beta}}
\newunicodechar{Γ}{\ensuremath{\Gamma}}
\newunicodechar{Δ}{\ensuremath{\Delta}}
\newunicodechar{Ε}{\ensuremath{\varepsilon}}
\newunicodechar{Θ}{\ensuremath{\Theta}}
\newunicodechar{Λ}{\ensuremath{\Lambda}}
\newunicodechar{Μ}{\ensuremath{\mu}}
\newunicodechar{Τ}{\ensuremath{\tau}}
\newunicodechar{Φ}{\ensuremath{\Phi}}
\newunicodechar{Ψ}{\ensuremath{\Psi}}
\newunicodechar{Ω}{\ensuremath{\Omega}}
\newunicodechar{↦}{\ensuremath{\mapsto}}
\newunicodechar{∈}{\ensuremath{\in}}
\newunicodechar{∉}{\ensuremath{\notin}}
\newunicodechar{∧}{\ensuremath{\wedge}}
\newunicodechar{⇔}{\ensuremath{\Leftrightarrow}}
\newunicodechar{⟺}{\ensuremath{\Longleftrightarrow}}
\newunicodechar{⇒}{\ensuremath{\Rightarrow}}
\newunicodechar{⟹}{\ensuremath{\Longrightarrow}}
\newunicodechar{∘}{\ensuremath{\circ}}
\newunicodechar{∪}{\ensuremath{\cup}}
\newunicodechar{∩}{\ensuremath{\cap}}
\newunicodechar{≡}{\ensuremath{\equiv}}
\newunicodechar{⊂}{\ensuremath{\subset}}
\newunicodechar{⊥}{\ensuremath{\perp}}
\newunicodechar{∃}{\ensuremath{\exists}}
\newunicodechar{∀}{\ensuremath{\forall}}
\newunicodechar{∛}{\ensuremath{\sqrt[3]{\,}}}
\newunicodechar{∜}{\ensuremath{\sqrt[4]{\,}}}
\newunicodechar{⃗}{\ensuremath{{}^{\to}}}
\newunicodechar{ⁿ}{\ifmmode^{n}\else\textsuperscript{\ensuremath{n}}\fi}
\newunicodechar{ˣ}{\ifmmode^{x}\else\textsuperscript{\ensuremath{x}}\fi}
\newunicodechar{ᵇ}{\ifmmode^{b}\else\textsuperscript{\ensuremath{b}}\fi}
\newunicodechar{ʳ}{\ifmmode^{r}\else\textsuperscript{\ensuremath{r}}\fi}
\newunicodechar{ᵘ}{\ifmmode^{u}\else\textsuperscript{\ensuremath{u}}\fi}
\newunicodechar{ʸ}{\ifmmode^{y}\else\textsuperscript{\ensuremath{y}}\fi}
\newunicodechar{ᵃ}{\ifmmode^{a}\else\textsuperscript{\ensuremath{a}}\fi}
\newunicodechar{ᵉ}{\ifmmode^{e}\else\textsuperscript{\ensuremath{e}}\fi}
\newunicodechar{ᵏ}{\ifmmode^{k}\else\textsuperscript{\ensuremath{k}}\fi}
\newunicodechar{ᵖ}{\ifmmode^{p}\else\textsuperscript{\ensuremath{p}}\fi}
\newunicodechar{ᴺ}{\ifmmode^{N}\else\textsuperscript{\ensuremath{N}}\fi}
\newunicodechar{ᶜ}{\ifmmode^{c}\else\textsuperscript{\ensuremath{c}}\fi}
\newunicodechar{ʰ}{\ifmmode^{h}\else\textsuperscript{\ensuremath{h}}\fi}
\newunicodechar{ᵠ}{\ifmmode^{q}\else\textsuperscript{\ensuremath{q}}\fi}
\newunicodechar{ₙ}{\ifmmode_{n}\else\textsubscript{\ensuremath{n}}\fi}
\newunicodechar{ₐ}{\ifmmode_{a}\else\textsubscript{\ensuremath{a}}\fi}
\newunicodechar{ᵢ}{\ifmmode_{i}\else\textsubscript{\ensuremath{i}}\fi}
\newunicodechar{ᵣ}{\ifmmode_{r}\else\textsubscript{\ensuremath{r}}\fi}
\newunicodechar{ₖ}{\ifmmode_{k}\else\textsubscript{\ensuremath{k}}\fi}
\newunicodechar{ᵦ}{\ifmmode_{\beta}\else\textsubscript{\ensuremath{\beta}}\fi}
\newunicodechar{ₓ}{\ifmmode_{x}\else\textsubscript{\ensuremath{x}}\fi}
\newfontfamily\popdisplay{ArchivoBlack-Regular.ttf}[Path=../../../fonts/]
\newfontfamily\poplogo{SpaceGrotesk-Bold.ttf}[Path=../../../fonts/]
\newfontfamily\popsemi{PlusJakartaSans-SemiBold.ttf}[Path=../../../fonts/]
\newfontfamily\popxbold{PlusJakartaSans-ExtraBold.ttf}[Path=../../../fonts/]
\newfontfamily\popxboldls{PlusJakartaSans-ExtraBold.ttf}[Path=../../../fonts/,LetterSpace=3.0]

\setlist[enumerate]{leftmargin=1.7em,itemsep=5pt,topsep=4pt}
\setlist[itemize]{leftmargin=1.7em,itemsep=4pt,topsep=4pt,label={\color{poporange}\textbullet}}
\setlength{\parindent}{0pt}
\setlength{\parskip}{5pt}
\renewcommand{\arraystretch}{1.25}

% pgfplots : style Pop par défaut
\pgfplotsset{
  every axis/.append style={axis line style={popink,line width=1pt},tick style={popink},
    grid style={popline,line width=0.5pt},label style={font=\small},tick label style={font=\footnotesize}},
  popcurve/.style={poporange,line width=1.8pt,no markers,smooth},
}
% Même style disponible dans un \draw TikZ simple (hors pgfplots)
\tikzset{popcurve/.style={poporange,line width=1.8pt}}
\tikzset{popgrid/.style={popline,very thin}}
\colorlet{popcurve}{poporange} % le nom du style est parfois utilisé comme couleur

% ============================================================
% Pill / squiggle
% ============================================================
\newcommand{\poppill}[3]{%
  \tikz[baseline=-0.55ex]\node[draw=popink,line width=1.2pt,rounded corners=2.6pt,%
    fill=#2,inner xsep=6.5pt,inner ysep=3pt]{%
    \popxboldls\fontsize{7}{7}\selectfont\color{#3}\MakeUppercase{#1}};%
}
\newcommand{\popsquiggle}[1]{%
  \tikz[baseline=-0.4ex]{
    \draw[#1,line width=2.6pt,line cap=round]
      plot [smooth] coordinates {(0,0.09) (0.34,0.01) (0.68,0.21) (1.02,0.01) (1.36,0.10)};
  }%
}
% Mot-clé surligné (marqueur orange sous le texte)
\newcommand{\cle}[1]{{\popxbold\color{popdark}#1}}

% ============================================================
% En-tête / pied
% ============================================================
\pagestyle{fancy}
\fancyhf{}
\renewcommand{\headrulewidth}{0pt}
\renewcommand{\footrulewidth}{0pt}
\lhead{\raisebox{-0.15ex}{\poplogo\fontsize{13}{13}\selectfont\color{popink}OnBuch\color{poporange}+}}
\rhead{\poppill{\DOCMATIERE\ \textbullet\ \DOCNIVEAU\ \textbullet\ Leçon \DOCLECON}{white}{popink}}
\lfoot{\popxbold\fontsize{7.6}{7.6}\selectfont\color{popmuted}\MakeUppercase{Cours signé OnBuch+}\ \textbullet\ {\popsemi\DOCTITRE}}
\rfoot{\tikz[baseline=-0.6ex]\node[rounded corners=8.5pt,fill=popink,minimum height=17pt,minimum width=17pt,inner xsep=5pt,inner sep=0]{\popxbold\fontsize{8}{8}\selectfont\color{popcream}\thepage\,/\,\pageref*{LastPage}};}

% ============================================================
% Titres
% ============================================================
\newcommand{\chaptitre}[2]{%
  \begin{tcolorbox}[enhanced,colback=poporange,colframe=popink,boxrule=2.6pt,arc=11pt,
      drop shadow=popink,left=15pt,right=15pt,top=12pt,bottom=11pt,before skip=3pt,after skip=18pt]
    \poppill{#2}{popink}{popcream}\par\vspace{9pt}
    {\raggedright\hyphenpenalty=10000\exhyphenpenalty=10000\popdisplay\fontsize{22}{24}\selectfont\color{popcream}\MakeUppercase{#1}\par}
  \end{tcolorbox}
}
% Section numérotée : pastille orange avec le numéro + titre Archivo
\newcounter{popsec}
\newcommand{\coursec}[1]{%
  \stepcounter{popsec}\setcounter{popsub}{0}%
  \par\vspace{16pt}\noindent
  \begin{minipage}{\linewidth}
  \tikz[baseline=-0.7ex]\node[draw=popink,line width=1.6pt,fill=poporange,rounded corners=4pt,minimum size=22pt,inner sep=0]{\popdisplay\fontsize{12}{12}\selectfont\color{popcream}\thepopsec};%
  \hspace{8pt}{\popdisplay\fontsize{15}{17}\selectfont\color{popink}\MakeUppercase{#1}}\par
  \vspace{4pt}\noindent{\color{poporange}\rule{36pt}{3.6pt}}
  \end{minipage}\par\nopagebreak\vspace{8pt}
}
\newcounter{popsub}
\newcommand{\courssub}[1]{%
  \stepcounter{popsub}%
  \par\vspace{9pt}\noindent{\popxbold\fontsize{11.5}{13}\selectfont\color{popink}{\color{poporange}\thepopsec.\thepopsub}\hspace{6pt}#1}\par\nopagebreak\vspace{4pt}
}

% ============================================================
% Boîtes pédagogiques — même recette : fond blanc, bordure épaisse colorée,
% ombre dure encre, badge plein. Titre optionnel [#1].
% ============================================================
\tcbset{popbase/.style={breakable,enhanced,colback=white,boxrule=2.3pt,arc=9pt,
  shadow={3pt}{-3pt}{0pt}{popink},left=14pt,right=14pt,top=11pt,bottom=11pt,before skip=11pt,after skip=15pt,
  fontupper=\popsemi\fontsize{10.6}{14.5}\selectfont\color{popink},
  fontlower=\popsemi\fontsize{10.6}{14.5}\selectfont\color{popink}}}
% Coche dessinée (le glyphe ✓ n'existe pas dans les polices du document)
\newcommand{\popcheck}[1]{\tikz[baseline=-0.5ex]\draw[#1,line width=1.6pt,line cap=round,line join=round](-0.13,0.0)--(-0.04,-0.09)--(0.13,0.1);}
\providecommand{\coche}{\popcheck{popgreen}}
\providecommand{\resultat}[1]{\par\smallskip\noindent\fcolorbox{popgreen}{popgreenL}{\parbox{\dimexpr\linewidth-2\fboxsep-2\fboxrule\relax}{\textbf{Résultat.} #1}}\par\smallskip}
\newcommand{\popboxhead}[3]{\poppill{#1}{#2}{white}\ifx\relax#3\relax\else\hspace{6pt}{\popxbold\fontsize{10.6}{12}\selectfont\color{popink}#3}\fi\par\vspace{7pt}}

\newtcolorbox{aretenir}[1][]{popbase,colframe=popgreen,before upper={\popboxhead{À retenir}{popgreen}{#1}}}
% Texte en allemand (document de lecture, dialogue, extrait) : cadre
% d'étude. Un titre optionnel (source, auteur) s'affiche dans l'en-tête.
\newtcolorbox{textebox}[1][]{popbase,colframe=popblue,before upper={\popboxhead{Text}{popblue}{#1}\begin{otherlanguage}{ngerman}},after upper={\end{otherlanguage}}}
% Marques ✓ / ✗ (forme correcte / fautive) souvent utilisées par les modèles
\providecommand{\cross}{\textcolor{poppink}{\ensuremath{\times}}}
\providecommand{\xmark}{\cross}\providecommand{\crossmark}{\cross}\providecommand{\cmark}{\ensuremath{\checkmark}}
% Ligne de réponse / justification à compléter (inventée par les modèles)
\providecommand{\justification}[1]{\par\noindent\textit{Justification~:} \dotfill\par}
% \textipa (package tipa non chargé) : la police ne couvre pas l'API -> simple passage
\providecommand{\textipa}[1]{#1}
% Soulignement (ulem non chargé) : \uline, \ul
\providecommand{\uline}[1]{\underline{#1}}\providecommand{\ul}[1]{\underline{#1}}
% \glossaire{\item ...} inventée par les modèles : liste de glossaire
\providecommand{\glossaire}[1]{\begin{itemize}#1\end{itemize}}
% \alert (beamer) inventée par les modèles : texte mis en évidence
\providecommand{\alert}[1]{\textbf{\textcolor{poppink}{#1}}}
% variantes ASCII d'ordinaux écrites en macro
\providecommand{\iere}{\textsuperscript{re}}\providecommand{\ieme}{\textsuperscript{e}}\providecommand{\ere}{\textsuperscript{re}}\providecommand{\eme}{\textsuperscript{e}}\providecommand{\er}{\textsuperscript{er}}\providecommand{\ie}{\textsuperscript{e}}
% Ordinaux français écrits en macro par les modèles : 1\ière, 3\ième
\newcommand{\ière}{\textsuperscript{re}}\newcommand{\ième}{\textsuperscript{e}}
% \exemple (inventée par les modèles) : étiquette « Exemple : »
\providecommand{\exemple}{\textbf{Exemple~:}\ }
% \square utilisable aussi hors mode math (cases à cocher)
\AtBeginDocument{\let\squareMath\square\renewcommand{\square}{\ensuremath{\squareMath}}}
% Mot ou phrase en allemand à l'intérieur d'un paragraphe français
\newcommand{\all}[1]{\ifmmode\text{\textit{#1}}\else\begingroup\selectlanguage{ngerman}\itshape #1\endgroup\fi}
\let\esp\all % alias toléré
\newtcolorbox{exemplebox}[1][]{popbase,colframe=poporange,before upper={\popboxhead{Exemple}{poporange}{#1}}}
\newtcolorbox{definition}[1][]{popbase,colframe=popblue,before upper={\popboxhead{Définition}{popblue}{#1}}}
\newtcolorbox{propriete}[1][]{popbase,colframe=poppurple,before upper={\popboxhead{Propriété}{poppurple}{#1}}}
\newtcolorbox{methode}[1][]{popbase,colframe=popgold,before upper={\popboxhead{Méthode}{popgold}{#1}}}
\newtcolorbox{attention}[1][]{popbase,colframe=poppink,before upper={\popboxhead{Attention}{poppink}{#1}}}
\newtcolorbox{experience}[1][]{popbase,colframe=popblue,colback=popblueL,before upper={\popboxhead{Expérience}{popblue}{#1}}}
% « Le savais-tu ? » : carte violette pleine (contexte, culture, Cameroun)
\newtcolorbox{savaistu}[1][]{popbase,colback=poppurple,colframe=popink,
  fontupper=\popsemi\fontsize{10.4}{14.2}\selectfont\color{white},
  before upper={\poppill{Le savais-tu\,?}{white}{poppurple}\ifx\relax#1\relax\else\hspace{6pt}{\popxbold\color{white}#1}\fi\par\vspace{7pt}}}
% Exercice résolu : énoncé en haut, \tcblower puis la solution sur fond crème
\newcounter{popexo}\newcounter{popexr}
\newtcolorbox{exoresolu}[1][]{popbase,colframe=poporange,colbacklower=popcream,
  segmentation style={popink,line width=1.2pt,dashed},
  before upper={\stepcounter{popexr}\popboxhead{Exercice résolu \thepopexr}{poporange}{#1}},
  before lower={\poppill{Solution}{popink}{popcream}\par\vspace{6pt}}}
% Exercice (non résolu en place) — la correction est dans la partie Corrigés
\newtcolorbox{exercice}[1][]{popbase,colframe=popink,
  before upper={\stepcounter{popexo}\popboxhead{Exercice \thepopexo}{popink}{#1}}}
% Corrigé
\newtcolorbox{corrige}[1][]{popbase,colframe=popgreen,colback=popgreenL,
  before upper={\popboxhead{Corrigé}{popgreen}{#1}}}
% Objectifs / prérequis / bilan
\newtcolorbox{objectifs}{popbase,colframe=popink,colback=poporangeL,
  before upper={\popboxhead{Objectifs de la leçon}{popink}{}}}
\newtcolorbox{prerequis}{popbase,colframe=popink,colback=popblueL,
  before upper={\popboxhead{Prérequis}{popblue}{}}}
\newtcolorbox{bilan}{popbase,colframe=popink,colback=white,boxrule=2.8pt,
  before upper={\popboxhead{Fiche bilan}{poporange}{L'essentiel à connaître pour l'examen}}}
% Figure encadrée avec légende
\newtcolorbox{popfigure}[1][]{enhanced,breakable=false,colback=white,colframe=popline,boxrule=1.4pt,arc=8pt,
  left=10pt,right=10pt,top=10pt,bottom=8pt,before skip=10pt,after skip=12pt,halign=center,
  lower separated=false,halign lower=center,
  fontlower=\popsemi\fontsize{9}{11.5}\selectfont\color{popmuted},
  #1}
\newcommand{\legende}[1]{\par\vspace{6pt}{\popsemi\fontsize{9}{11.5}\selectfont\color{popmuted}#1}}

% ============================================================
% Signature OnBuch+
% ============================================================
\newcommand{\poplogoinline}[1]{{\poplogo\fontsize{#1}{#1}\selectfont\color{popink}OnBuch\color{poporange}+}}
\newcommand{\signatureonbuch}{%
  \begin{tcolorbox}[enhanced,colback=popink,colframe=popink,boxrule=2.2pt,arc=10pt,
      drop shadow=poporange,left=14pt,right=14pt,top=10pt,bottom=10pt,before skip=0pt,after skip=0pt]
    \tikz[baseline=-0.7ex]\node[circle,fill=popgreen,draw=popcream,line width=1.3pt,minimum size=22pt,inner sep=0]{\popcheck{white}};%
    \hspace{9pt}\begin{minipage}[c]{0.62\linewidth}
      {\popxbold\fontsize{12}{13}\selectfont\color{popcream}Signé\ \textbullet\ {\poplogo OnBuch\color{poporange}+}}\par\vspace{2pt}
      {\raggedright\popsemi\fontsize{8}{10.5}\selectfont\color{popline}\MakeUppercase{Équipe pédagogique OnBuch+}\\\MakeUppercase{Conforme au programme officiel MINESEC}\par}
    \end{minipage}\hfill
    {\poplogo\fontsize{22}{22}\selectfont\color{popcream}OnBuch\color{poporange}+}
  \end{tcolorbox}%
}

% ============================================================
% Page de garde
% #1 titre · #2 matière · #3 niveau · #4 module · #5 n° leçon · #6 durée ·
% #7 description / objectifs (texte ou itemize)
% ============================================================
\newcommand{\pagegarde}[7]{%
  \thispagestyle{empty}
  \newgeometry{a4paper,margin=1.7cm,top=1.6cm,bottom=1.4cm}%
  \begin{tikzpicture}[remember picture,overlay]
    \fill[poporange!18] ([xshift=-3.2cm,yshift=-3.6cm]current page.north east) circle (4.6cm);
    \fill[poppurple!14] ([xshift=2.4cm,yshift=7.6cm]current page.south west) circle (3.8cm);
  \end{tikzpicture}%
  {\poplogo\fontsize{30}{30}\selectfont\color{popink}OnBuch\color{poporange}+}\hfill
  \raisebox{0.6ex}{\poppill{Cours signé \textbullet\ Premium}{popink}{popcream}}\par
  \vspace{1pt}\popsquiggle{poppurple}\par
  \vspace{8pt}
  \begin{tcolorbox}[enhanced,colback=poporange,colframe=popink,boxrule=2.8pt,arc=13pt,
      drop shadow=popink,left=18pt,right=18pt,top=12pt,bottom=13pt,after skip=10pt]
    \poppill{#2 \textbullet\ #3}{popink}{popcream}\hspace{5pt}\poppill{#4}{white}{popink}\par\vspace{8pt}
    {\popdisplay\fontsize{14}{14}\selectfont\color{popink}LEÇON #5}\par\vspace{4pt}
    {\raggedright\hyphenpenalty=10000\exhyphenpenalty=10000\popdisplay\fontsize{25}{27}\selectfont\color{popcream}\MakeUppercase{#1}\par}
    \vspace{8pt}
    {\popxbold\fontsize{9.5}{9.5}\selectfont\color{popink}\MakeUppercase{Durée conseillée : #6}}
  \end{tcolorbox}
  \begin{tcolorbox}[enhanced,colback=white,colframe=popink,boxrule=2.2pt,arc=10pt,
      drop shadow=popink,left=16pt,right=16pt,top=10pt,bottom=10pt,after skip=8pt]
    \poppill{Dans cette leçon}{poppurple}{white}\par\vspace{5pt}
    \popsemi\fontsize{9.3}{12.6}\selectfont\color{popink}#7
  \end{tcolorbox}
  \noindent\begin{minipage}[t]{0.487\linewidth}
    \begin{tcolorbox}[enhanced,colback=popgreen,colframe=popink,boxrule=2.2pt,arc=10pt,
        drop shadow=popink,left=11pt,right=11pt,top=10pt,bottom=10pt,height=3.1cm,valign=center]
      \tcbox[colback=white,colframe=popink,boxrule=1.3pt,arc=5pt,boxsep=0pt,nobeforeafter,
        left=3pt,right=3pt,top=3pt,bottom=3pt,valign=center]{\qrcode[height=1.9cm]{https://play.google.com/store/apps/details?id=com.luvvix.onbuch&hl=fr}}%
      \hspace{8pt}\begin{minipage}[c]{0.5\linewidth}
        {\popdisplay\fontsize{11}{12.5}\selectfont\color{white}\MakeUppercase{Télécharge OnBuch}}\par\vspace{4pt}
        {\raggedright\popsemi\fontsize{8.4}{11}\selectfont\color{white}Gratuit \textbullet\ Google Play\\Tuteur IA, annales, concours, résultats\par}
      \end{minipage}
    \end{tcolorbox}
  \end{minipage}\hfill
  \begin{minipage}[t]{0.487\linewidth}
    \begin{tcolorbox}[enhanced,colback=poppurple,colframe=popink,boxrule=2.2pt,arc=10pt,
        drop shadow=popink,left=11pt,right=11pt,top=10pt,bottom=10pt,height=3.1cm,valign=center]
      \tikz[baseline=-0.7ex]\node[circle,fill=white,draw=popink,line width=1.3pt,minimum size=1.6cm,inner sep=0]{\popdisplay\fontsize{24}{24}\selectfont\color{poppurple}+};%
      \hspace{8pt}\begin{minipage}[c]{0.6\linewidth}
        {\popdisplay\fontsize{11}{12.5}\selectfont\color{white}\MakeUppercase{Passe à OnBuch+}}\par\vspace{4pt}
        {\raggedright\popsemi\fontsize{8.4}{11}\selectfont\color{white}Tous les cours signés, Tuteur IA illimité, fiches PDF — onglet OnBuch+ de l'app\par}
      \end{minipage}
    \end{tcolorbox}
  \end{minipage}
  \vspace{8pt}
  \popcontenu
  \vfill
  \signatureonbuch
  \restoregeometry
  \newpage
}

% Rangée « Ce cours contient » (page de garde)
\newcommand{\poptile}[3]{% #1 couleur #2 gros texte #3 légende
  \begin{tcolorbox}[enhanced,colback=#1,colframe=popink,boxrule=2pt,arc=9pt,shadow={2.5pt}{-2.5pt}{0pt}{popink},
      left=6pt,right=6pt,top=9pt,bottom=9pt,halign=center,valign=center,height=1.75cm,width=\linewidth,nobeforeafter]
    {\popdisplay\fontsize{12.5}{13}\selectfont\color{white}\MakeUppercase{#2}}\par\vspace{3pt}
    {\centering\popsemi\fontsize{7.8}{9.5}\selectfont\color{white}#3\par}
  \end{tcolorbox}}
\newcommand{\popcontenu}{%
  \par\noindent{\popxboldls\fontsize{7.5}{7.5}\selectfont\color{popmuted}CE COURS SIGNÉ CONTIENT}\par\vspace{7pt}
  \noindent\begin{minipage}[t]{0.235\linewidth}\poptile{popblue}{Cours}{complet, illustré, pas à pas}\end{minipage}\hfill
  \begin{minipage}[t]{0.235\linewidth}\poptile{poporange}{Méthodes}{et exercices résolus}\end{minipage}\hfill
  \begin{minipage}[t]{0.235\linewidth}\poptile{poppink}{Exercices}{type examen + corrigés}\end{minipage}\hfill
  \begin{minipage}[t]{0.235\linewidth}\poptile{popgreen}{Fiche bilan}{l'essentiel à retenir}\end{minipage}\par}

% Sommaire
\newtcolorbox{sommaire}{enhanced,colback=white,colframe=popink,boxrule=2.2pt,arc=10pt,
  drop shadow=popink,left=18pt,right=18pt,top=13pt,bottom=13pt,before skip=6pt,after skip=20pt,
  before upper={\poppill{Sommaire}{poppurple}{white}\par\vspace{9pt}\popsemi\fontsize{10.6}{14.5}\selectfont\color{popink}}}

% Signature de fin de cours — poussée tout en bas de la dernière page
\newcommand{\findecours}{%
  \par\vfill\nopagebreak
  \begin{center}\popsquiggle{poporange}\end{center}\vspace{4pt}
  \signatureonbuch
}
\AtBeginDocument{\def\llbracket{\mathopen{[\mkern-3.5mu[}}\def\rrbracket{\mathclose{]\mkern-3.5mu]}}} % intervalles d'entiers (glyphe ⟦ absent de la police)
% Glyphes absents de Fira Math : redéfinis à partir de glyphes disponibles
\AtBeginDocument{%
  \def\setminus{\mathbin{\backslash}}\let\smallsetminus\setminus
  \def\vdots{\mathord{\vbox{\baselineskip3.2pt\lineskiplimit0pt\kern4pt\hbox{.}\hbox{.}\hbox{.}}}}%
  \def\ddots{\mathinner{\mkern1mu\raise6.5pt\hbox{.}\mkern2mu\raise3.6pt\hbox{.}\mkern2mu\raise0.7pt\hbox{.}\mkern1mu}}%
  \def\triangle{\mathord{\scalebox{0.9}{$\symup{\Delta}$}}}%
  \def\nabla{\mathord{\raisebox{\depth}{\scalebox{1}[-1]{$\symup{\Delta}$}}}}%
  \def\checkmark{\popcheck{popdark}}%
  \def\star{\mathbin{\ast}}\def\bigstar{\mathord{\ast}}%
  \def\diamond{\mathbin{\scalebox{0.6}{$\Diamond$}}}%
  \def\bigcup{\mathop{\mathchoice{\raisebox{-0.3em}{\scalebox{1.7}{$\cup$}}}{\scalebox{1.25}{$\cup$}}{\cup}{\cup}}\displaylimits}%
  \def\bigcap{\mathop{\mathchoice{\raisebox{-0.3em}{\scalebox{1.7}{$\cap$}}}{\scalebox{1.25}{$\cap$}}{\cap}{\cap}}\displaylimits}%
  \def\lesseqqgtr{\mathrel{\lessgtr}}\def\bigoplus{\mathop{\oplus}\displaylimits}%
}
\providecommand{\ssi}{si et seulement si} % abréviation fréquente
\AtBeginDocument{\providecommand{\wideparen}{\overparen}} % arc de cercle

%% FIGFIX-BEGIN v4 — correctifs globaux des figures (docs/onbuch-generation/tools/figfix)
\usepackage{adjustbox}\usepackage{etoolbox}
% toute figure TikZ/pgfplots plus large que la ligne est réduite à la largeur disponible
\BeforeBeginEnvironment{tikzpicture}{\begin{adjustbox}{max width=\linewidth}}
\AfterEndEnvironment{tikzpicture}{\end{adjustbox}}
% légendes pgfplots lisibles par-dessus les courbes
\pgfplotsset{every axis/.append style={legend style={fill=white,fill opacity=0.92,text opacity=1,draw=black!15,inner sep=3pt}}}
% étiquettes d'arêtes (syntaxe edge["..."]) avec fond blanc
\tikzset{every edge quotes/.append style={fill=white,inner sep=1.2pt}}
%% FIGFIX-END
\begin{document}

\pagegarde{Traduction et production de textes liés à la citoyenneté et d'interviews divers}{Allemand}{1ère A}{Chapitre 5 : Citoyenneté}{ALA27}{2 h}{Cette leçon mixte entraîne les élèves à produire et à traduire des textes courts (interview, lettre, petit discours) sur le thème de la citoyenneté. Elle réinvestit l'ensemble de la grammaire du programme de Première à travers le thème et la version.
\vspace{4pt}
\begin{multicols}{2}\small
\begin{itemize}
\item À la fin de cette leçon, je sais définir la citoyenneté et utiliser le lexique allemand de la vie civique
\item je sais lire et comprendre un texte allemand sur l'engagement citoyen des jeunes
\item je sais rédiger un petit discours électoral structuré en allemand
\item je sais construire et mener un interview (questions, relances, remerciements)
\item je sais rédiger une lettre formelle à une autorité (Maire, Bürgermeister)
\item je sais traduire de courtes phrases et un texte du français vers l'allemand (thème)
\end{itemize}\end{multicols}}

\begin{sommaire}
\begin{multicols}{2}
\begin{enumerate}
\item Texte support : Jugendliche engagieren sich
\item Lexique thématique de la citoyenneté
\item Produire : interview, lettre, petit discours
\item Révision grammaticale pour la traduction
\item Entraînement à la traduction : thème et version
\item Commentaire modèle et production guidée
\item Activité d'intégration
\item Méthodes et exercices résolus
\item Exercices
\item Corrigés des exercices
\item Fiche bilan
\end{enumerate}
\end{multicols}
\end{sommaire}

\begin{objectifs}
\begin{itemize}
\item À la fin de cette leçon, je sais définir la citoyenneté et utiliser le lexique allemand de la vie civique
\item je sais lire et comprendre un texte allemand sur l'engagement citoyen des jeunes
\item je sais rédiger un petit discours électoral structuré en allemand
\item je sais construire et mener un interview (questions, relances, remerciements)
\item je sais rédiger une lettre formelle à une autorité (Maire, Bürgermeister)
\item je sais traduire de courtes phrases et un texte du français vers l'allemand (thème)
\item je sais traduire un texte court de l'allemand vers le français (version)
\item je sais réinvestir les structures grammaticales du programme (ordre des mots, cas, temps, subordonnées)
\item je sais relire et corriger ma production (orthographe, majuscules, place du verbe, cas)
\end{itemize}
\end{objectifs}

\begin{prerequis}
\begin{itemize}
\item L'ordre des mots en allemand : verbe en position 2, verbe conjugué et infinitif/participe en fin de phrase
\item Les cas (nominatif, accusatif, datif, génitif) et les déclinaisons des articles et adjectifs
\item Les temps principaux : Präsens, Perfekt, Präteritum, Futur I
\item Les propositions subordonnées (weil, dass, wenn, obwohl) et les pronoms relatifs
\item Le vocabulaire de base de la vie quotidienne et scolaire acquis en Seconde
\end{itemize}
\end{prerequis}

\newpage

\coursec{Texte support : Jugendliche engagieren sich}

\courssub{Situation d'accroche et problématique}

Imaginez la scène : à Yaoundé, dans votre lycée, c'est la semaine des élections des délégués de classe. Certains élèves rédigent des discours, d'autres préparent des affiches, et un petit groupe interviewe les candidats pour le journal du lycée. Pendant ce temps, à Munich, des jeunes de seize ans se préparent à voter aux élections européennes, et d'autres siègent dans un \all{Jugendparlament}, un parlement des jeunes, où ils défendent leurs idées devant les élus de la ville.

La citoyenneté n'est donc pas seulement une notion abstraite : c'est une pratique quotidienne, au Cameroun comme en Allemagne. Mais comment présenter ses idées citoyennes en allemand ? Comment mener un interview avec un correspondant de Munich ou de Vienne ? Comment traduire un texte sur les droits et les devoirs du citoyen ? Cette leçon vous donne les outils pour écrire, traduire et parler de la citoyenneté en allemand.

\textbf{Problématique :} \emph{Qu'est-ce qu'un bon citoyen, et comment exprimer l'engagement citoyen en allemand ?}

\textbf{Objectifs de la leçon :} à la fin de cette leçon, vous serez capable de (1) comprendre et commenter un texte sur la citoyenneté, (2) rédiger un interview, une lettre ou un petit discours, (3) traduire des phrases et un texte court du français vers l'allemand et de l'allemand vers le français.

\courssub{Lecture du texte : Jugendliche engagieren sich}

Lisez le texte suivant à voix haute, puis en silence. Soulignez les mots que vous ne comprenez pas et essayez de deviner leur sens grâce au contexte avant de consulter le glossaire.

\begin{textebox}[Jugendliche engagieren sich]
In Deutschland können Jugendliche ab 16 Jahren bei Europawahlen wählen. Viele engagieren sich auch in Jugendparlamenten oder im Freiwilligendienst. Sie helfen alten Menschen, schützen die Umwelt oder arbeiten in Flüchtlingsheimen. Lisa aus München zum Beispiel arbeitet jeden Samstag in einem Tierheim, und Paul aus Wien sammelt Müll im Park. „Ich engagiere mich, weil ich die Gesellschaft verändern will“, sagt Lisa. „Man muss nicht nur reden, man muss auch handeln.“

Auch in Kamerun zeigen junge Leute Zivilcourage: Sie organisieren Saubermachaktionen in ihren Vierteln und helfen bei Wahlen. In Yaoundé und Douala putzen Schüler ihre Schulen und ihre Straßen. Ein guter Bürger kennt seine Rechte, aber auch seine Pflichten: Er respektiert die Gesetze, zahlt Steuern und hilft seiner Gemeinschaft. Er informiert sich über die Politik und denkt an die anderen, nicht nur an sich selbst.

Viele Jugendliche engagieren sich, weil sie die Gesellschaft verändern wollen. Sie lernen dabei auch viel: Sie sprechen vor Gruppen, sie arbeiten im Team und sie verstehen die Demokratie besser. Engagement macht die Demokratie stark, sagen die Experten. Und die Jugendlichen selbst sagen: „Wir sind die Zukunft, aber wir wollen auch heute schon etwas bewegen.“
\end{textebox}

\begin{definition}[Glossaire : le vocabulaire de la citoyenneté]
\begin{itemize}
\item \cle{die Staatsbürgerschaft} (\emph{die Staatsbürgerschaften}) = la nationalité, la citoyenneté ; \cle{die Bürgerschaft} = la citoyenneté, l'ensemble des citoyens
\item \cle{der Bürger} (\emph{die Bürger}) = le citoyen ; \cle{die Bürgerin} (\emph{die Bürgerinnen}) = la citoyenne
\item \cle{sich engagieren} = s'engager (verbe réfléchi : \all{ich engagiere mich, du engagierst dich})
\item \cle{die Wahl} (\emph{die Wahlen}) = l'élection, le choix ; \cle{wählen} = choisir, élire, voter
\item \cle{die Pflicht} (\emph{die Pflichten}) = le devoir ; \cle{das Recht} (\emph{die Rechte}) = le droit
\item \cle{der Freiwilligendienst} (\emph{die Freiwilligendienste}) = le service volontaire ; \cle{freiwillig} = volontaire, bénévole
\item \all{das Jugendparlament} (\emph{die Jugendparlamente}) = le parlement des jeunes ; \all{die Europawahl} (\emph{die Europawahlen}) = l'élection européenne
\item \all{das Flüchtlingsheim} (\emph{die Flüchtlingsheime}) = le foyer de réfugiés ; \all{das Tierheim} (\emph{die Tierheime}) = le refuge pour animaux
\item \all{die Zivilcourage} (sans pluriel) = le courage civique ; \all{die Saubermachaktion} (\emph{die Saubermachaktionen}) = l'opération de nettoyage
\item \all{das Gesetz} (\emph{die Gesetze}) = la loi ; \all{die Steuer} (\emph{die Steuern}) = l'impôt ; \all{die Gemeinschaft} (\emph{die Gemeinschaften}) = la communauté
\item \all{die Gesellschaft} (\emph{die Gesellschaften}) = la société ; \all{verändern} = changer ; \all{handeln} = agir ; \all{etwas bewegen} = faire bouger les choses
\end{itemize}
\end{definition}

\courssub{Compréhension globale du texte}

Avant d'entrer dans les détails, répondez aux questions essentielles : de quoi parle le texte ? Qui ? Où ?

\begin{exercice}[Compréhension globale]
Répondez en français, puis essayez en allemand avec une phrase simple.
\begin{enumerate}
\item De quel thème parle ce texte ?
\item Qui sont les personnes concernées ?
\item Dans quels pays se situe l'action ?
\item Quel est le message principal de la dernière phrase ?
\end{enumerate}
\end{exercice}

\begin{corrige}[Exercice : compréhension globale]
\begin{enumerate}
\item Le texte parle de l'engagement citoyen des jeunes (\all{Der Text handelt vom Engagement der Jugendlichen.}).
\item Les personnes concernées sont les jeunes, les adolescents (\all{die Jugendlichen}), comme Lisa et Paul.
\item L'action se situe en Allemagne, en Autriche (Vienne) et au Cameroun (Yaoundé et Douala) : \all{in Deutschland, in Österreich und in Kamerun}.
\item Message principal : les jeunes veulent agir dès aujourd'hui, pas seulement attendre l'avenir (\all{Wir wollen auch heute schon etwas bewegen.}).
\end{enumerate}
\end{corrige}

\courssub{Compréhension fine : droits, devoirs et exemples d'engagement}

Cherchons maintenant dans le texte les informations précises. Un bon lecteur sait relever (\all{herausfinden}) les éléments concrets.

\begin{exoresolu}[Relever les droits, les devoirs et les exemples d'engagement]
Énoncé : Relevez dans le texte (a) un droit cité, (b) trois devoirs cités, (c) quatre exemples d'engagement des jeunes.
\tcblower
Solution détaillée :
\begin{enumerate}
\item \textbf{Un droit :} \all{Jugendliche können ab 16 Jahren bei Europawahlen wählen.} — le droit de vote dès 16 ans aux élections européennes.
\item \textbf{Trois devoirs :} \all{Er respektiert die Gesetze, zahlt Steuern und hilft seiner Gemeinschaft.} — respecter les lois, payer les impôts, aider sa communauté.
\item \textbf{Quatre exemples d'engagement :} aider les personnes âgées (\all{alten Menschen helfen}), protéger l'environnement (\all{die Umwelt schützen}), travailler dans un foyer de réfugiés (\all{in Flüchtlingsheimen arbeiten}), organiser des opérations de nettoyage (\all{Saubermachaktionen organisieren}).
\end{enumerate}
Méthode : pour relever, cherchez d'abord les verbes d'action (\all{helfen, schützen, arbeiten, organisieren}) et les noms abstraits (\all{Rechte, Pflichten}).
\end{exoresolu}

\courssub{Observation grammaticale : les structures utiles du texte}

Relisez ces deux phrases tirées du texte :

\all{Viele Jugendliche engagieren sich, weil sie die Gesellschaft verändern wollen.} « Beaucoup de jeunes s'engagent parce qu'ils veulent changer la société. »

\all{Ein guter Bürger denkt an die anderen.} « Un bon citoyen pense aux autres. »

\begin{propriete}[La subordonnée avec \all{weil}]
\all{weil} (parce que) introduit une \cle{subordonnée de cause}. Règle d'or : dans la subordonnée, le \textbf{verbe conjugué va à la fin} de la phrase.
\begin{itemize}
\item \all{Ich engagiere mich, weil ich die Gesellschaft verändern will.} — le verbe conjugué \all{will} est rejeté à la fin, après l'infinitif \all{verändern}.
\item \all{Er wählt, weil er seine Pflicht kennt.} « Il vote parce qu'il connaît son devoir. »
\end{itemize}
Attention : si la subordonnée commence la phrase, le verbe de la principale vient juste après la virgule : \all{Weil er seine Pflicht kennt, wählt er.} (le verbe \all{wählt} reste en deuxième position de la phrase complète). Même règle avec \all{dass} (que), \all{wenn} (si/quand), \all{obwohl} (bien que) : \all{Ich finde es gut, dass junge Leute Zivilcourage zeigen.}
\end{propriete}

\begin{propriete}[Les verbes à préposition]
Certains verbes allemands se construisent avec une préposition fixe, comme en français « penser à », « s'intéresser à ». Il faut les apprendre avec leur préposition et leur cas.
\begin{center}
\begin{tabularx}{\linewidth}{|X|X|X|}
\hline
\rowcolor{popblueL} Verbe + préposition & Français & Exemple \\
\hline
\all{sich engagieren für + Akk.} & s'engager pour & \all{Sie engagiert sich für die Umwelt.} \\
\hline
\all{denken an + Akk.} & penser à & \all{Er denkt an die anderen.} \\
\hline
\all{helfen + Dat.} & aider (quelqu'un) & \all{Wir helfen alten Menschen.} \\
\hline
\all{sich informieren über + Akk.} & s'informer sur & \all{Er informiert sich über die Politik.} \\
\hline
\all{sprechen über + Akk.} & parler de & \all{Wir sprechen über die Wahl.} \\
\hline
\all{sich interessieren für + Akk.} & s'intéresser à & \all{Ich interessiere mich für Politik.} \\
\hline
\end{tabularx}
\end{center}
\end{propriete}

\begin{attention}[Pièges fréquents des francophones]
\begin{itemize}
\item Ne dites jamais \all{weil sie wollen die Gesellschaft verändern} : avec \all{weil}, le verbe conjugué va à la fin. En français, l'ordre ne change pas ; en allemand, si.
\item \all{helfen} se construit avec le \textbf{datif} : \all{ich helfe dem alten Mann} (et non \all{den alten Mann}). Comparez : « j'aide \emph{à} l'homme ».
\item \all{die Wahl} signifie « l'élection / le choix », pas « la valeur ». Faux ami ! « La valeur » se dit \all{der Wert}.
\item N'oubliez pas le pronom réfléchi : \all{ich engagiere \textbf{mich}}, pas \all{ich engagiere} seul.
\item \all{die Politik} est singulier en allemand : \all{Er informiert sich über die Politik} (et non \all{die Politiken}).
\end{itemize}
\end{attention}

\courssub{Carte mentale : le lexique de la citoyenneté}

\begin{popfigure}
\begin{tikzpicture}[mindmap, grow cyclic, every node/.style={concept, text=white, align=center}, concept color=popblue, root concept/.append style={font=\large\bfseries}, level 1/.append style={level distance=3.2cm, sibling angle=90}, level 2/.append style={level distance=2.6cm, sibling angle=45, concept color=popblue!60}]
\node[root concept]{Bürgerschaft}
  child[concept color=popgreen] { node {Rechte} child { node {wählen} } child { node {sich informieren} } }
  child[concept color=poporange] { node {Pflichten} child { node {Gesetze respektieren} } child { node {Steuern zahlen} } }
  child[concept color=poppurple] { node {Engagement} child { node {Freiwilligendienst} } child { node {helfen} } }
  child[concept color=poppink] { node {Wahlen} child { node {die Wahl} } child { node {der Bürger} } };
\end{tikzpicture}
\legende{Carte mentale du vocabulaire de la citoyenneté : quatre familles de mots autour de \all{Bürgerschaft}.}
\end{popfigure}

\coursec{Produire : interview, lettre, petit discours}

La fiche de progression demande de produire trois types de textes sur la citoyenneté : l'interview, la lettre et le petit discours. Chacun a sa structure et ses formules propres.

\courssub{L'interview : das Interview}

Un interview est un dialogue entre un journaliste (\all{der Journalist}) et une personne interviewée. Il alterne questions (\all{die Frage}) et réponses (\all{die Antwort}).

\begin{methode}[Rédiger un interview en cinq étapes]
\begin{enumerate}
\item \textbf{Préparez les questions} avec les mots interrogatifs : \all{Was} (que), \all{Wer} (qui), \all{Wann} (quand), \all{Wo} (où), \all{Warum} (pourquoi), \all{Wie} (comment), \all{Wie lange} (depuis combien de temps). Dans la question, le verbe conjugué est en deuxième position : \all{Warum engagieren Sie sich?}
\item \textbf{Ouvrez l'interview} par une formule de présentation : \all{Guten Tag, Frau Ngono. Danke für das Interview.} « Bonjour Madame Ngono. Merci pour cet interview. »
\item \textbf{Alternez} question et réponse, en notant le nom du locuteur suivi de deux-points : \all{Journalist: ... / Frau Ngono: ...}
\item \textbf{Utilisez le vouvoiement} \all{Sie} (avec majuscule) pour un adulte ou une personnalité ; \all{du} seulement entre jeunes.
\item \textbf{Clôturez} poliment : \all{Vielen Dank für das Gespräch!} « Merci beaucoup pour cet entretien ! »
\end{enumerate}
\end{methode}

\begin{textebox}[Modèle : Interview mit einer Kandidatin]
\textbf{Journalist:} Guten Tag, Amina. Du bist Kandidatin für die Wahl der Klassensprecher. Warum willst du Klassensprecherin werden?\\
\textbf{Amina:} Ich will meiner Klasse helfen. Ich finde, jeder Schüler hat Rechte, aber auch Pflichten.\\
\textbf{Journalist:} Was willst du verändern?\\
\textbf{Amina:} Ich will eine saubere Schule und mehr Sport. Ich engagiere mich auch für die Umwelt.\\
\textbf{Journalist:} Wie willst du das machen?\\
\textbf{Amina:} Ich organisiere Saubermachaktionen, und ich spreche mit den Lehrern. Man muss nicht nur reden, man muss auch handeln!\\
\textbf{Journalist:} Vielen Dank für das Gespräch, Amina. Viel Glück bei der Wahl!\\
\textbf{Amina:} Danke!
\end{textebox}

\begin{definition}[Vocabulaire de l'interview et des élections scolaires]
\begin{itemize}
\item \all{das Interview} (\emph{die Interviews}) = l'interview ; \all{interviewen} = interviewer
\item \all{der Klassensprecher} (\emph{die Klassensprecher}) = le délégué de classe ; \all{die Klassensprecherin} = la déléguée
\item \all{der Kandidat} (\emph{die Kandidaten}) = le candidat ; \all{die Kandidatin} = la candidate
\item \all{die Frage} (\emph{die Fragen}) = la question ; \all{die Antwort} (\emph{die Antworten}) = la réponse
\item \all{das Gespräch} (\emph{die Gespräche}) = l'entretien, la conversation
\item \all{Viel Glück!} = bonne chance !
\end{itemize}
\end{definition}

\courssub{La lettre : der Brief}

La lettre (amicale ou officielle simple) suit un schéma fixe qu'il faut respecter à l'examen.

\begin{propriete}[Structure de la lettre allemande]
\begin{enumerate}
\item \textbf{Lieu et date} en haut à droite : \all{Yaoundé, den 15. März 2025} (attention : \all{den} + date ordinale avec point).
\item \textbf{Formule d'appel} : \all{Lieber Paul,} (à un garçon), \all{Liebe Lisa,} (à une fille), \all{Sehr geehrter Herr Direktor,} (formel, « Monsieur le Directeur »). Après la virgule, la phrase suivante commence par une \textbf{minuscule}.
\item \textbf{Corps de la lettre} : introduction, développement, conclusion.
\item \textbf{Formule de politesse} : \all{Viele Grüße} (amicale) ou \all{Mit freundlichen Grüßen} (formelle), sans virgule après.
\item \textbf{Signature} : prénom ou nom.
\end{enumerate}
\end{propriete}

\begin{textebox}[Modèle : Brief an einen Freund]
Yaoundé, den 15. März 2025

Lieber Paul,\\
danke für deinen Brief. Du fragst, was ein guter Bürger ist. Meiner Meinung nach kennt ein guter Bürger seine Rechte und seine Pflichten. Er respektiert die Gesetze und hilft seiner Gemeinschaft. In meiner Schule engagiere ich mich für die Umwelt: Wir putzen jeden Samstag unser Viertel. Das macht mir Spaß, weil ich etwas bewegen will.\\
Und du? Engagierst du dich auch in deiner Schule? Schreib mir bald!\\
Viele Grüße\\
Dein Etienne
\end{textebox}

\courssub{Le petit discours : die kleine Rede}

Le discours (par exemple un discours de candidat) suit le plan classique : accroche, arguments, appel à l'action.

\begin{exemplebox}[Formules pour le petit discours]
\begin{itemize}
\item \textbf{Ouverture :} \all{Liebe Mitschülerinnen und Mitschüler!} « Chers camarades ! » / \all{Sehr geehrte Damen und Herren!} « Mesdames et Messieurs ! »
\item \textbf{Présentation :} \all{Ich heiße Etienne und ich bin Kandidat für die Wahl.} « Je m'appelle Etienne et je suis candidat à l'élection. »
\item \textbf{Arguments :} \all{Erstens will ich..., zweitens will ich...} « Premièrement je veux..., deuxièmement je veux... »
\item \textbf{Appel à l'action :} \all{Wählt mich!} « Votez pour moi ! » / \all{Helft mit!} « Participez ! »
\item \textbf{Clôture :} \all{Danke für eure Aufmerksamkeit!} « Merci de votre attention ! »
\end{itemize}
\end{exemplebox}

\begin{attention}[Impératif et vouvoiement dans le discours]
\begin{itemize}
\item Impératif à \all{du} : \all{Wähl\textbf{e} mich!} (verbe + \all{-e}, souvent omise : \all{Wähl mich!}).
\item Impératif à \all{ihr} : \all{Wähl\textbf{t} mich!} ; à \all{Sie} : \all{Wähl\textbf{en Sie} mich!}
\item Dans une lettre ou un discours, les pronoms personnels qui désignent le destinataire prennent la majuscule : \all{Du}, \all{Dein}, \all{Ihr}, \all{Euer} (ex. \all{Danke für Deinen Brief}).
\end{itemize}
\end{attention}

\coursec{Révision grammaticale pour la traduction}

Pour traduire correctement, il faut maîtriser quatre points du programme de Première : la place du verbe, les cas, les temps et les verbes modaux.

\courssub{La place du verbe : la règle des trois positions}

\begin{propriete}[Les trois positions du verbe conjugué]
\begin{enumerate}
\item \textbf{Position 2} dans la phrase déclarative : \all{Der Bürger respektiert die Gesetze.} Si un autre élément ouvre la phrase, le sujet passe après le verbe (inversion) : \all{Heute respektiert der Bürger die Gesetze.}
\item \textbf{Position 1} dans la question sans mot interrogatif et à l'impératif : \all{Respektiert der Bürger die Gesetze?} / \all{Respektiere die Gesetze!}
\item \textbf{Position finale} dans la subordonnée (\all{weil, dass, wenn, obwohl}) : \all{..., weil der Bürger die Gesetze respektiert.}
\end{enumerate}
\end{propriete}

\begin{popfigure}
\begin{tikzpicture}[node distance=0.4cm, every node/.style={align=center}]
\node[draw=popblue, fill=popblueL, rounded corners, minimum width=3.4cm, minimum height=0.9cm, align=center] (a){Position 2\\ \all{Er wählt heute.}};
\node[draw=popgreen, fill=popgreenL, rounded corners, minimum width=3.4cm, minimum height=0.9cm, right=1.1cm of a, align=center] (b){Position 1\\ \all{Wählt er heute?}};
\node[draw=poporange, fill=poporangeL, rounded corners, minimum width=3.4cm, minimum height=0.9cm, right=1.1cm of b, align=center] (c){Position finale\\ \all{..., weil er heute wählt.}};
\end{tikzpicture}
\legende{Les trois positions du verbe conjugué en allemand : la clé de toute traduction réussie.}
\end{popfigure}

\courssub{Les cas et les verbes modaux}

\begin{propriete}[Rappel : les quatre cas]
\begin{center}
\begin{tabular}{|l|l|l|l|}
\hline
\rowcolor{popblueL} Cas & der (m.) & die (f.) & das (n.) \\
\hline
Nominatif (sujet) & der Bürger & die Bürgerin & das Recht \\
\hline
Akkusativ (COD) & den Bürger & die Bürgerin & das Recht \\
\hline
Datif (COI) & dem Bürger & der Bürgerin & dem Recht \\
\hline
Génitif (possession) & des Bürgers & der Bürgerin & des Rechts \\
\hline
\end{tabular}
\end{center}
Seul le masculin change à l'accusatif (\all{der} devient \all{den}) : c'est le cas à surveiller en priorité dans vos traductions.
\end{propriete}

\begin{propriete}[Les verbes modaux au présent]
Le modal est conjugué en position 2 ; l'infinitif va à la fin de la phrase : \all{Man muss handeln.} « Il faut agir. »
\begin{center}
\begin{tabular}{|l|l|l|l|}
\hline
\rowcolor{popblueL} Modal & ich/er & Français & Exemple \\
\hline
\all{müssen} & muss & devoir (obligation) & \all{Jeder muss Steuern zahlen.} \\
\hline
\all{können} & kann & pouvoir & \all{Jugendliche können wählen.} \\
\hline
\all{wollen} & will & vouloir & \all{Sie will die Gesellschaft verändern.} \\
\hline
\all{dürfen} & darf & avoir le droit de & \all{Man darf die Gesetze nicht brechen.} \\
\hline
\all{sollen} & soll & devoir (recommandation) & \all{Du sollst den anderen helfen.} \\
\hline
\end{tabular}
\end{center}
\end{propriete}

\courssub{Les temps : Präsens, Perfekt, Futur}

\begin{propriete}[Choisir le bon temps pour traduire]
\begin{itemize}
\item \textbf{Präsens} = présent et futur proche : \all{Ich engagiere mich.} « Je m'engage. »
\item \textbf{Perfekt} = passé composé (langue parlée et écrite courante) : auxiliaire \all{haben} ou \all{sein} + participe II à la fin : \all{Ich habe mich engagiert.} « Je me suis engagé. » ; \all{Wir sind nach Yaoundé gefahren.} « Nous sommes allés à Yaoundé. » (verbes de déplacement : auxiliaire \all{sein}).
\item \textbf{Futur} = \all{werden} + infinitif à la fin : \all{Wir werden die Gesellschaft verändern.} « Nous changerons la société. »
\end{itemize}
\end{propriete}

\begin{attention}[Verbes forts fréquents : participe II]
Apprenez les participes des verbes du texte : \all{helfen — half — geholfen} ; \all{sprechen — sprach — gesprochen} ; \all{verstehen — verstand — verstanden} ; \all{wählen — wählte — gewählt} (verbe faible) ; \all{kennen — kannte — gekannt} (verbe mixte). Les verbes en \all{-ieren} (\all{engagieren, organisieren, informieren}) forment le participe \textbf{sans} \all{ge-} : \all{engagiert, organisiert, informiert}.
\end{attention}

\coursec{Entraînement à la traduction : thème et version}

\courssub{Méthode de traduction}

\begin{methode}[Traduire en cinq étapes]
\begin{enumerate}
\item \textbf{Lisez} la phrase entière et identifiez le sujet, le verbe, les compléments.
\item \textbf{Choisissez} le temps et le mode (présent, Perfekt, futur).
\item \textbf{Placez} le verbe conjugué : position 2 (principale) ou finale (subordonnée).
\item \textbf{Vérifiez} les cas : accusatif après \all{für, durch, ohne}, datif après \all{mit, zu, bei, helfen}.
\item \textbf{Relisez} : majuscules des noms, Umlaute, orthographe, accord du verbe avec le sujet.
\end{enumerate}
\end{methode}

\courssub{Thème : du français vers l'allemand}

\begin{exoresolu}[Thème guidé : phrases sur la citoyenneté]
Énoncé : Traduisez en allemand. (a) Un bon citoyen respecte les lois. (b) Nous aidons les personnes âgées. (c) Les jeunes s'engagent parce qu'ils veulent changer la société. (d) Demain, j'interviewerai le délégué de classe.
\tcblower
Solution détaillée :
\begin{enumerate}
\item \all{Ein guter Bürger respektiert die Gesetze.} — sujet au nominatif, \all{die Gesetze} à l'accusatif pluriel (identique au nominatif).
\item \all{Wir helfen alten Menschen.} — \all{helfen} + datif pluriel : \all{alten Menschen} (datif pluriel en \all{-n} sur l'adjectif ; \all{Menschen} porte déjà le \all{-n} du datif pluriel).
\item \all{Die Jugendlichen engagieren sich, weil sie die Gesellschaft verändern wollen.} — subordonnée avec \all{weil} : verbe conjugué \all{wollen} à la fin, après l'infinitif.
\item \all{Morgen interviewe ich den Klassensprecher.} — inversion après \all{Morgen} ; \all{der Klassensprecher} devient \all{den} à l'accusatif masculin. Futur possible aussi : \all{Morgen werde ich den Klassensprecher interviewen.}
\end{enumerate}
\end{exoresolu}

\begin{exercice}[Thème : à vous]
Traduisez en allemand.
\begin{enumerate}
\item Chaque citoyen doit payer des impôts.
\item Elle s'intéresse à la politique.
\item Nous avons organisé une opération de nettoyage à Douala.
\item Je trouve bien que les jeunes aident leur communauté.
\end{enumerate}
\end{exercice}

\begin{corrige}[Exercice : thème]
\begin{enumerate}
\item \all{Jeder Bürger muss Steuern zahlen.} (modal \all{muss} en position 2, infinitif \all{zahlen} à la fin.)
\item \all{Sie interessiert sich für die Politik.} (\all{sich interessieren für + Akk.})
\item \all{Wir haben eine Saubermachaktion in Douala organisiert.} (Perfekt : auxiliaire \all{haben}, participe \all{organisiert} sans \all{ge-}.)
\item \all{Ich finde es gut, dass die Jugendlichen ihrer Gemeinschaft helfen.} (\all{dass} + verbe à la fin ; \all{helfen} + datif : \all{ihrer Gemeinschaft}.)
\end{enumerate}
\end{corrige}

\courssub{Version : de l'allemand vers le français}

\begin{exoresolu}[Version guidée : un court texte]
Énoncé : Traduisez en français le texte suivant.\\
\all{In Douala engagieren sich viele Schüler in ihrem Viertel. Sie putzen die Straßen und helfen alten Menschen. „Wir wollen unsere Stadt verändern“, sagt Marie, eine Schülerin. „Ein guter Bürger muss handeln, nicht nur reden.“ Die Lehrer finden das Engagement der Jugendlichen sehr wichtig, weil es die Gemeinschaft stark macht.}
\tcblower
Solution détaillée :\\
« À Douala, beaucoup d'élèves s'engagent dans leur quartier. Ils nettoient les rues et aident les personnes âgées. "Nous voulons changer notre ville", dit Marie, une élève. "Un bon citoyen doit agir, pas seulement parler." Les professeurs trouvent l'engagement des jeunes très important, parce qu'il rend la communauté forte. »\\
Points de vigilance : \all{das Viertel} = le quartier ; \all{stark machen} = rendre fort, renforcer ; dans la subordonnée \all{weil es die Gemeinschaft stark macht}, le verbe \all{macht} à la fin se traduit normalement en français.
\end{exoresolu}

\begin{exercice}[Version : à vous]
Traduisez en français.\\
\all{Paul aus Wien sammelt jeden Samstag Müll im Park. Er sagt: „Ich engagiere mich für die Umwelt, weil ich die Natur liebe.“ Seine Freunde helfen ihm oft. Zusammen wollen sie ihre Stadt sauber machen.}
\end{exercice}

\begin{corrige}[Exercice : version]
« Paul, de Vienne, ramasse des déchets dans le parc chaque samedi. Il dit : "Je m'engage pour l'environnement parce que j'aime la nature." Ses amis l'aident souvent. Ensemble, ils veulent rendre leur ville propre. »\\
Rappels : \all{der Müll} = les déchets ; \all{die Umwelt} = l'environnement ; \all{sauber machen} = nettoyer, rendre propre ; \all{ihm} = lui (datif de \all{er} après \all{helfen}).
\end{corrige}

\coursec{Commentaire modèle et production guidée}

\courssub{Méthode : le commentaire de texte en trois parties}

Au baccalauréat et dans les évaluations de Première, on vous demandera souvent de commenter un texte court. Voici la démarche en trois étapes, valable en allemand comme en français.

\begin{methode}[Kommentar : commenter un texte en trois parties]
\begin{enumerate}
\item \textbf{Einleitung (introduction) :} présentez le texte — le thème, le type de texte (article, discours, interview), le contexte. Exemple : \all{Der Text handelt vom Engagement der Jugendlichen in Deutschland und Kamerun.} « Le texte traite de l'engagement des jeunes en Allemagne et au Cameroun. »
\item \textbf{Hauptteil (développement) :} dégagez deux ou trois idées principales et illustrez-les par des citations courtes traduites. Exemple : \all{Der Autor zeigt zuerst die Rechte, dann die Pflichten eines Bürgers. Zum Beispiel: „Er respektiert die Gesetze.“} — citez toujours entre guillemets allemands „…“ et traduisez.
\item \textbf{Schluss (conclusion) :} donnez votre avis personnel avec une formule d'opinion : \all{Meiner Meinung nach ist Engagement sehr wichtig, weil es die Demokratie stark macht.} « À mon avis, l'engagement est très important parce qu'il rend la démocratie forte. » Autres formules utiles : \all{Ich finde, dass...} / \all{Ich bin der Meinung, dass...}
\end{enumerate}
\end{methode}

\begin{exemplebox}[Phrases modèles pour le commentaire]
\begin{itemize}
\item \all{Der Text handelt von den Rechten und Pflichten eines guten Bürgers.} « Le texte traite des droits et des devoirs d'un bon citoyen. »
\item \all{Viele Jugendliche engagieren sich, weil sie die Gesellschaft verändern wollen.} « Beaucoup de jeunes s'engagent parce qu'ils veulent changer la société. »
\item \all{Meiner Meinung nach muss jeder Bürger seine Pflichten kennen.} « À mon avis, chaque citoyen doit connaître ses devoirs. »
\item \all{Ich finde es gut, dass junge Leute in Kamerun Zivilcourage zeigen.} « Je trouve bien que les jeunes au Cameroun fassent preuve de courage civique. »
\end{itemize}
\end{exemplebox}

\courssub{Production guidée : votre discours de candidat}

\begin{exercice}[Production écrite : petit discours de candidat]
Vous êtes candidat(e) à l'élection des délégués de votre lycée à Yaoundé. Rédigez un discours de 8 à 10 phrases en allemand. Respectez le plan : (1) salutation et présentation, (2) deux ou trois projets avec \all{ich will} ou \all{ich werde}, (3) une justification avec \all{weil}, (4) un appel à voter, (5) une formule de clôture. Utilisez au moins trois mots du glossaire et un verbe modal.
\end{exercice}

\begin{corrige}[Exercice : production écrite (corrigé possible)]
\all{Liebe Mitschülerinnen und Mitschüler! Ich heiße Sandrine und ich bin Kandidatin für die Wahl der Klassensprecher. Ich will unserer Klasse helfen, weil jeder Schüler Rechte und Pflichten hat. Erstens will ich eine saubere Schule: Ich organisiere Saubermachaktionen jeden Samstag. Zweitens will ich mehr Sport und mehr Bücher für unsere Klasse. Ich werde mit den Lehrern sprechen, und ich werde euch immer zuhören. Man muss nicht nur reden, man muss auch handeln! Wählt mich, und zusammen können wir unsere Schule verändern. Danke für eure Aufmerksamkeit!}\\
« Chers camarades ! Je m'appelle Sandrine et je suis candidate à l'élection des délégués. Je veux aider notre classe parce que chaque élève a des droits et des devoirs. Premièrement, je veux une école propre : j'organise des opérations de nettoyage chaque samedi. Deuxièmement, je veux plus de sport et plus de livres pour notre classe. Je parlerai avec les professeurs et je vous écouterai toujours. Il ne faut pas seulement parler, il faut aussi agir ! Votez pour moi, et ensemble nous pourrons changer notre école. Merci de votre attention ! »\\
Critères d'évaluation : plan respecté, verbe en position 2, verbe final après \all{weil}, majuscules des noms, au moins un modal (\all{können}) correctement placé.
\end{corrige}

\begin{aretenir}[L'essentiel de la leçon]
\begin{itemize}
\item Le vocabulaire clé : \all{der Bürger / die Bürgerin}, \all{das Recht}, \all{die Pflicht}, \all{die Wahl / wählen}, \all{sich engagieren}, \all{der Freiwilligendienst}.
\item Avec \all{weil}, \all{dass}, \all{wenn}, le verbe conjugué va à la fin de la subordonnée.
\item Les verbes à préposition s'apprennent avec leur cas : \all{helfen + Dat.}, \all{denken an + Akk.}, \all{sich engagieren für + Akk.}
\item Trois textes à savoir produire : l'interview (questions-réponses, vouvoiement), la lettre (date, appel, formule finale), le discours (salutation, arguments, appel à l'action).
\item Pour traduire : identifier sujet et verbe, choisir le temps, placer le verbe, vérifier les cas, relire les majuscules.
\item Un commentaire de texte suit toujours le plan : Einleitung — Hauptteil — Schluss, avec citations traduites et avis personnel (\all{Meiner Meinung nach...}).
\end{itemize}
\end{aretenir}

\begin{savaistu}[Le « Freiwilliges Soziales Jahr » (FSJ)]
En Allemagne, le \all{Freiwilliges Soziales Jahr} (FSJ), « année sociale volontaire », permet aux jeunes de 16 à 27 ans de travailler pendant un an dans le domaine social, de la santé, de la culture ou de l'environnement — par exemple dans un hôpital, une école ou un foyer pour personnes âgées. Les volontaires reçoivent un petit argent de poche et des formations. Ce service est très apprécié : il aide les jeunes à choisir leur métier et montre que la citoyenneté se vit aussi par l'action concrète. Au Cameroun, de nombreux jeunes s'engagent de manière semblable dans des associations de quartier, des clubs de nettoyage ou des actions caritatives pendant les vacances.
\end{savaistu}

\begin{experience}[À vous de parler : mini-débat citoyen]
Par groupes de quatre, répondez à la question : \all{Was ist ein guter Bürger?} « Qu'est-ce qu'un bon citoyen ? » Chaque élève donne une idée avec la structure \all{Ein guter Bürger...} + verbe (ex. \all{Ein guter Bürger respektiert die Gesetze.}), puis une justification avec \all{weil}. Un rapporteur note les idées et les présente à la classe : \all{Wir finden, dass ein guter Bürger...} Prolongement : deux groupes jouent ensuite un interview devant la classe, l'un dans le rôle du journaliste, l'autre dans le rôle d'un jeune engagé de Douala.
\end{experience}

\coursec{Lexique thématique de la citoyenneté}

Dans la section précédente, nous avons lu et commenté le texte \all{Jugendliche engagieren sich}. Nous y avons rencontré des mots comme \all{die Verantwortung}, \all{die Gemeinschaft} ou \all{sich engagieren}. Ces mots appartiennent à un même univers : celui de la \cle{citoyenneté}. Pour pouvoir parler, écrire et surtout \emph{traduire} des textes sur ce thème — à l'examen comme au Baccalauréat —, il faut d'abord maîtriser son vocabulaire de manière précise : chaque nom avec son \cle{article} et son \cle{pluriel}, chaque verbe avec sa \cle{construction}. C'est l'objectif de cette section.

\begin{methode}[Comment apprendre un mot allemand ?]
\begin{enumerate}
\item Apprends chaque \textbf{nom} avec son article (\all{der, die, das}) et son pluriel : \all{die Wahl, -en}.
\item Apprends chaque \textbf{verbe} avec sa construction : \all{sich engagieren für + Akkusativ}.
\item Note toujours un \textbf{exemple complet} traduit en français.
\item Classe les mots par \textbf{champ lexical} (familles de sens) : c'est plus facile à mémoriser.
\end{enumerate}
\end{methode}

\courssub{Champ lexical 1 : la citoyenneté et l'État}

Observons d'abord cette phrase extraite de notre texte support : \all{Junge Bürger lernen, was der Staat und die Demokratie bedeuten.} « Les jeunes citoyens apprennent ce que l'État et la démocratie signifient. » On y trouve trois noms fondamentaux : \all{der Bürger} (le citoyen), \all{der Staat} (l'État), \all{die Demokratie} (la démocratie).

\begin{definition}[Les mots de l'État]
\begin{itemize}
\item \all{die Staatsbürgerschaft, -en} : la nationalité, la citoyenneté (statut juridique)
\item \all{der Bürger, - / die Bürgerin, -nen} : le citoyen / la citoyenne
\item \all{der Staat, -en} : l'État
\item \all{die Verfassung, -en} : la constitution
\item \all{das Gesetz, -e} : la loi
\item \all{die Demokratie, -n} : la démocratie
\item \all{die Regierung, -en} : le gouvernement
\end{itemize}
\end{definition}

\begin{center}
\begin{tabularx}{\linewidth}{|X|X|X|}\hline
\rowcolor{popblueL} Deutsch & Français & Remarque \\ \hline
die Staatsbürgerschaft & la nationalité & mot composé : \all{Staat} + \all{Bürgerschaft} \\ \hline
der Staat (-en) & l'État & en français aussi, majuscule quand on parle de l'institution \\ \hline
die Verfassung (-en) & la constitution & \all{die Verfassung schützt die Bürger} \\ \hline
das Gesetz (-e) & la loi & pluriel en \all{-e} \\ \hline
die Demokratie (-n) & la démocratie & accent sur la dernière syllabe : « dé-mo-kra-ti » \\ \hline
\end{tabularx}
\end{center}

\begin{exemplebox}[Exemples traduits]
\all{Die Verfassung garantiert die Rechte der Bürger.} « La constitution garantit les droits des citoyens. »

\all{Kamerun ist ein demokratischer Staat.} « Le Cameroun est un État démocratique. »

\all{Jeder muss die Gesetze respektieren.} « Chacun doit respecter les lois. »
\end{exemplebox}

\courssub{Champ lexical 2 : droits et devoirs}

Le cœur de la citoyenneté, c'est l'équilibre entre ce que le citoyen \emph{peut} faire (ses droits) et ce qu'il \emph{doit} faire (ses devoirs).

\begin{definition}[Droits et devoirs]
\begin{itemize}
\item \all{das Recht, -e} : le droit ; \all{die Pflicht, -en} : le devoir
\item \all{die Freiheit, -en} : la liberté
\item \all{die Meinungsfreiheit} : la liberté d'expression (littéralement « la liberté d'opinion »)
\item \all{die Verantwortung} : la responsabilité
\item \all{die Gleichheit} : l'égalité
\end{itemize}
\end{definition}

\begin{exemplebox}[Exemples traduits]
\all{Jeder Bürger hat das Recht zu wählen.} « Chaque citoyen a le droit de voter. »

\all{Die Meinungsfreiheit ist ein wichtiges Recht.} « La liberté d'expression est un droit important. »

\all{Mit den Rechten kommen auch die Pflichten.} « Avec les droits viennent aussi les devoirs. »

\all{Eltern tragen die Verantwortung für ihre Kinder.} « Les parents portent la responsabilité de leurs enfants. »
\end{exemplebox}

\begin{popfigure}
\begin{center}
\begin{tikzpicture}
\node[draw=popblue, fill=popblueL, rounded corners, very thick, minimum width=3cm, minimum height=1.2cm, align=center, font=\large\bfseries] (b) at (0,0) {der Bürger\\ \small le citoyen};
\node[draw=popgreen, fill=popgreenL, rounded corners, very thick, minimum width=3cm, minimum height=1cm, align=center] (r) at (-4.5,0) {die Rechte\\ \small les droits};
\node[draw=poporange, fill=poporangeL, rounded corners, very thick, minimum width=3cm, minimum height=1cm, align=center] (p) at (4.5,0) {die Pflichten\\ \small les devoirs};
\draw[-{Stealth}, very thick, popgreen] (r.north east) to[bend left=15] node[above, align=center, font=\small]{wählen, sich äußern} (b.north west);
\draw[-{Stealth}, very thick, popgreen] (b.south west) to[bend left=15] (r.south east);
\draw[-{Stealth}, very thick, poporange] (b.north east) to[bend left=15] node[above, align=center, font=\small]{Gesetze respektieren} (p.north west);
\draw[-{Stealth}, very thick, poporange] (p.south west) to[bend left=15] (b.south east);
\end{tikzpicture}
\end{center}
\legende{Le citoyen entre droits et devoirs : un équilibre fondamental de la démocratie.}
\end{popfigure}

\begin{aretenir}[La formule clé]
\all{das Recht haben, etwas zu tun} = avoir le droit de faire quelque chose.

\all{die Pflicht haben, etwas zu tun} = avoir le devoir de faire quelque chose.

\all{Jeder Bürger hat das Recht zu wählen, aber auch die Pflicht, die Gesetze zu respektieren.} « Chaque citoyen a le droit de voter, mais aussi le devoir de respecter les lois. »
\end{aretenir}

\courssub{Champ lexical 3 : les élections}

\begin{definition}[Le vocabulaire du vote]
\begin{itemize}
\item \all{die Wahl, -en} : l'élection ; le choix
\item \all{wählen (wählte, hat gewählt)} : élire, voter, choisir
\item \all{der Wähler, - / die Wählerin, -nen} : l'électeur / l'électrice
\item \all{der Kandidat, -en / die Kandidatin, -nen} : le candidat / la candidate
\item \all{das Wahllokal, -e} : le bureau de vote
\item \all{die Stimme, -n} : la voix ; le vote (suffrage)
\item \all{abstimmen (über + Akk.)} : voter (lors d'un scrutin), se prononcer sur
\end{itemize}
\end{definition}

\begin{center}
\begin{tabularx}{\linewidth}{|X|X|X|}\hline
\rowcolor{popblueL} Deutsch & Français & Exemple \\ \hline
die Wahl & l'élection / le choix & \all{Die Wahl ist am Sonntag.} \\ \hline
wählen & élire, voter & \all{Die Bürger wählen den Präsidenten.} \\ \hline
der Wähler & l'électeur & \all{Viele Wähler gehen ins Wahllokal.} \\ \hline
die Stimme & la voix, le vote & \all{Jede Stimme zählt.} \\ \hline
abstimmen & voter (scrutin) & \all{Wir stimmen über das Gesetz ab.} \\ \hline
\end{tabularx}
\end{center}

\begin{exemplebox}[Exemples traduits]
\all{Er nimmt an der Wahl teil.} « Il participe à l'élection. »

\all{Die Wähler geben ihre Stimme im Wahllokal ab.} « Les électeurs déposent leur vote au bureau de vote. »

\all{Der Kandidat spricht mit den Bürgern von Douala.} « Le candidat parle avec les habitants de Douala. »
\end{exemplebox}

\begin{attention}[Faux amis : \all{wählen} ne veut pas dire « téléphoner » !]
\begin{itemize}
\item \all{wählen} = élire, voter, choisir : \all{Ich wähle diesen Kandidaten.} « Je vote pour ce candidat. »
\item \all{anrufen} = téléphoner : \all{Ich rufe meine Mutter an.} « Je téléphone à ma mère. »
\item \all{die Wahl} a deux sens : « l'élection » \emph{et} « le choix » : \all{Du hast die Wahl!} « Tu as le choix ! »
\end{itemize}
\end{attention}

\courssub{Champ lexical 4 : l'engagement}

\begin{definition}[S'engager pour les autres]
\begin{itemize}
\item \all{sich engagieren (für + Akk.)} : s'engager, s'investir (pour)
\item \all{der Freiwillige, -n / die Freiwillige, -n} : le/la bénévole, volontaire
\item \all{die Hilfsorganisation, -en} : l'organisation humanitaire (d'aide)
\item \all{die Zivilcourage} : le courage civique
\item \all{die Gemeinschaft, -en} : la communauté
\item \all{helfen (hilft, half, hat geholfen) + Dat.} : aider
\end{itemize}
\end{definition}

\begin{definition}[Die Zivilcourage]
\all{Die Zivilcourage} désigne le \cle{courage civique} : oser agir ou intervenir pour défendre les autres, par exemple quand quelqu'un est attaqué ou traité injustement, au lieu de détourner le regard.

\all{Er zeigt Zivilcourage und hilft dem alten Mann.} « Il fait preuve de courage civique et aide le vieil homme. »
\end{definition}

\begin{attention}[\all{sich engagieren} : n'oublie pas le pronom réfléchi !]
\all{sich engagieren} est un verbe \cle{réfléchi}. En français on dit « je m'engage », en allemand le pronom est obligatoire lui aussi :

\all{Ich engagiere mich für die Umwelt.} « Je m'engage pour l'environnement. »

Faux : \emph{Ich engagiere für die Umwelt.} — il manque \all{mich} !

Conjugaison au présent : \all{ich engagiere mich, du engagierst dich, er/sie/es engagiert sich, wir engagieren uns, ihr engagiert euch, sie/Sie engagieren sich.}
\end{attention}

\begin{exemplebox}[Exemples traduits]
\all{Viele Freiwillige arbeiten für eine Hilfsorganisation in Yaoundé.} « Beaucoup de bénévoles travaillent pour une organisation humanitaire à Yaoundé. »

\all{Die Gemeinschaft hilft den armen Familien.} « La communauté aide les familles pauvres. »

\all{Sie engagiert sich für die Rechte der Kinder.} « Elle s'engage pour les droits des enfants. »
\end{exemplebox}

\courssub{Champ lexical 5 : l'interview et les médias}

\begin{definition}[Le vocabulaire de l'interview]
\begin{itemize}
\item \all{das Interview, -s} : l'interview, l'entretien
\item \all{der Journalist, -en / die Journalistin, -nen} : le/la journaliste
\item \all{die Frage, -n} : la question ; \all{eine Frage stellen} : poser une question
\item \all{die Antwort, -en} : la réponse ; \all{antworten (+ Dat.)} : répondre
\item \all{berichten (über + Akk.)} : rapporter, faire un reportage (sur)
\item \all{die Presse} : la presse ; \all{die Zeitung, -en} : le journal ; \all{der Artikel, -} : l'article
\end{itemize}
\end{definition}

\begin{exemplebox}[Exemples traduits]
\all{Der Journalist stellt dem Bürgermeister viele Fragen.} « Le journaliste pose beaucoup de questions au maire. »

\all{Die Zeitung berichtet über die Wahl in Kamerun.} « Le journal fait un reportage sur l'élection au Cameroun. »

\all{Im Interview antwortet der Kandidat auf alle Fragen.} « Dans l'interview, le candidat répond à toutes les questions. »
\end{exemplebox}

\begin{attention}[\all{antworten} : répondre à quelqu'un, avec le datif]
\all{antworten} se construit avec le \cle{datif} de la personne : \all{Ich antworte dem Journalisten.} « Je réponds au journaliste. »

Pour « répondre à une question », on utilise \all{antworten auf + Akkusativ} : \all{Er antwortet auf die Frage.} « Il répond à la question. »

Faux : \emph{Ich antworte den Journalisten.} — \all{antworten} ne prend jamais l'accusatif de la personne.
\end{attention}

\courssub{Verbes à préposition : les constructions à connaître par cœur}

En allemand, beaucoup de verbes se construisent avec une \cle{préposition fixe} qui ne correspond pas toujours au français. Ces constructions sont indispensables pour la traduction (thème et version).

\begin{center}
\begin{tabularx}{\linewidth}{|X|X|X|}\hline
\rowcolor{popblueL} Verbe + préposition & Français & Exemple traduit \\ \hline
sich engagieren für + Akk. & s'engager pour & \all{Er engagiert sich für die Schule.} « Il s'engage pour l'école. » \\ \hline
kämpfen für + Akk. & lutter pour & \all{Wir kämpfen für die Freiheit.} « Nous luttons pour la liberté. » \\ \hline
kämpfen gegen + Akk. & lutter contre & \all{Sie kämpft gegen die Armut.} « Elle lutte contre la pauvreté. » \\ \hline
sich interessieren für + Akk. & s'intéresser à & \all{Ich interessiere mich für Politik.} « Je m'intéresse à la politique. » \\ \hline
teilnehmen an + Dat. & participer à & \all{Er nimmt an der Wahl teil.} « Il participe à l'élection. » \\ \hline
warten auf + Akk. & attendre & \all{Die Wähler warten auf das Ergebnis.} « Les électeurs attendent le résultat. » \\ \hline
\end{tabularx}
\end{center}

\begin{attention}[\all{teilnehmen} : an + datif, et verbe à particule séparable !]
\all{teilnehmen} se construit avec \all{an + datif} : \all{an der Diskussion teilnehmen} = participer à la discussion.

De plus, c'est un verbe à \cle{particule séparable} : \all{Ich nehme an der Diskussion teil.} « Je participe à la discussion. » La particule \all{teil} va à la fin de la phrase. Au Perfekt : \all{Ich habe an der Diskussion teilgenommen.}
\end{attention}

\begin{savaistu}[La démocratie directe en Suisse]
En Suisse, les citoyens votent plusieurs fois par an par référendum : on appelle cela \all{die Volksabstimmung} (littéralement « le vote du peuple »). Ils se prononcent directement sur des lois, des impôts ou des projets. C'est ce qu'on appelle la \cle{démocratie directe} (\all{die direkte Demokratie}) : le peuple ne choisit pas seulement ses représentants, il décide lui-même. En Allemagne et en Autriche, en revanche, la démocratie est surtout \emph{représentative} : les citoyens élisent des députés qui votent les lois à leur place.
\end{savaistu}

\courssub{Texte de réinvestissement et exercice résolu}

\begin{textebox}[Ein Freiwilliger aus Yaoundé]
\all{Paul ist 19 Jahre alt und wohnt in Yaoundé. Er ist Bürger Kameruns und interessiert sich sehr für die Politik seines Landes. „Jeder Bürger hat Rechte, aber auch Pflichten", sagt er. „Ich habe das Recht zu wählen, und ich nehme an jeder Wahl teil."

Aber Paul engagiert sich auch in seiner Gemeinschaft: Er arbeitet als Freiwilliger für eine Hilfsorganisation. Jeden Samstag hilft er Kindern beim Lesen und Schreiben. „Viele junge Leute warten nur auf die Hilfe des Staates", erklärt er. „Aber wir müssen auch selbst etwas tun. Zivilcourage bedeutet: nicht wegsehen, sondern handeln!"

Letzte Woche hat ein Journalist einer Zeitung ein Interview mit Paul gemacht. Er hat viele Fragen gestellt: Warum engagierst du dich? Was bedeutet Demokratie für dich? Paul hat geantwortet: „Demokratie bedeutet Freiheit und Verantwortung. Wir kämpfen für eine bessere Zukunft, und wir kämpfen gegen die Armut. Jede Stimme zählt, und jeder Mensch kann helfen."}
\end{textebox}

\textbf{Glossaire :} \all{der Bürger, -} (le citoyen) ; \all{die Pflicht, -en} (le devoir) ; \all{die Hilfsorganisation, -en} (l'organisation humanitaire) ; \all{wegsehen} (détourner le regard) ; \all{handeln} (agir) ; \all{die Zukunft} (l'avenir) ; \all{die Armut} (la pauvreté) ; \all{zählen} (compter).

\textbf{Questions de compréhension :}
\begin{enumerate}
\item \all{Warum nimmt Paul an jeder Wahl teil?}
\item \all{Was macht Paul jeden Samstag?}
\item \all{Was bedeutet Zivilcourage für Paul?}
\item \all{Wie definiert Paul die Demokratie?}
\end{enumerate}

\begin{corrige}[Questions de compréhension]
1. \all{Er nimmt an jeder Wahl teil, weil er das Recht zu wählen hat.} « Il participe à chaque élection parce qu'il a le droit de voter. »

2. \all{Jeden Samstag hilft er Kindern beim Lesen und Schreiben.} « Chaque samedi, il aide des enfants à lire et à écrire. »

3. \all{Für Paul bedeutet Zivilcourage: nicht wegsehen, sondern handeln.} « Pour Paul, le courage civique signifie : ne pas détourner le regard, mais agir. »

4. \all{Demokratie bedeutet für ihn Freiheit und Verantwortung.} « Pour lui, la démocratie signifie la liberté et la responsabilité. »
\end{corrige}

\begin{exoresolu}[Traduction (thème) : phrases sur la citoyenneté]
Énoncé : traduis en allemand les phrases suivantes.

1. Chaque citoyen a le droit de voter. 2. Je m'engage pour l'environnement. 3. Elle participe à la discussion. 4. Nous luttons contre la pauvreté. 5. Les électeurs attendent le résultat.
\tcblower
Solution détaillée :

1. \all{Jeder Bürger hat das Recht zu wählen.} — Structure : \all{das Recht haben + zu + infinitif} ; \all{wählen} se place à la fin.

2. \all{Ich engagiere mich für die Umwelt.} — Verbe réfléchi : ne pas oublier \all{mich} ; préposition \all{für + Akkusativ}.

3. \all{Sie nimmt an der Diskussion teil.} — \all{teilnehmen an + Dativ} : \all{an der Diskussion} ; particule \all{teil} à la fin (verbe séparable).

4. \all{Wir kämpfen gegen die Armut.} — « Lutter contre » = \all{kämpfen gegen + Akkusativ}.

5. \all{Die Wähler warten auf das Ergebnis.} — « Attendre » = \all{warten auf + Akkusativ}, jamais \all{warten} seul.
\end{exoresolu}

\begin{exercice}[À toi de jouer !]
1. Traduis en français : \all{Die Presse berichtet über die Wahl, und der Journalist stellt dem Kandidaten viele Fragen.}

2. Traduis en allemand : « La constitution garantit la liberté d'expression. » ; « Beaucoup de jeunes s'intéressent à la politique. »

3. Complète avec la bonne préposition : \all{Er kämpft ... die Freiheit. / Wir nehmen ... der Wahl teil. / Sie wartet ... das Ergebnis.}
\end{exercice}

\begin{corrige}[Exercice « À toi de jouer ! »]
1. « La presse fait un reportage sur l'élection, et le journaliste pose beaucoup de questions au candidat. »

2. \all{Die Verfassung garantiert die Meinungsfreiheit.} ; \all{Viele junge Leute interessieren sich für die Politik.} (Attention au réfléchi \all{sich} et à \all{für + Akkusativ}.)

3. \all{Er kämpft für die Freiheit. / Wir nehmen an der Wahl teil. / Sie wartet auf das Ergebnis.}
\end{corrige}

\begin{aretenir}[L'essentiel du lexique]
\begin{itemize}
\item Chaque nom s'apprend avec article et pluriel : \all{die Wahl, -en ; das Recht, -e ; die Pflicht, -en}.
\item Les verbes à préposition fixe : \all{sich engagieren für, kämpfen für/gegen, sich interessieren für, teilnehmen an (+ Dat.), warten auf (+ Akk.)}.
\item \all{sich engagieren} : verbe réfléchi, le pronom est obligatoire.
\item \all{antworten} : datif de la personne ; \all{antworten auf + Akk.} pour la chose.
\item \all{wählen} = voter/élire ; « téléphoner » = \all{anrufen}.
\item Le citoyen (\all{der Bürger}) vit dans l'équilibre entre \all{Rechte} (droits) et \all{Pflichten} (devoirs).
\end{itemize}
\end{aretenir}

Ce lexique est maintenant ta boîte à outils : dans la section suivante, tu l'utiliseras pour \emph{produire} tes propres textes — une interview, une lettre et un petit discours de citoyen engagé.

\coursec{Produire : interview, lettre, petit discours}

Dans la section précédente, nous avons construit le lexique thématique de la citoyenneté : \all{die Wahl}, \all{der Bürger}, \all{die Rechte und Pflichten}, \all{sich engagieren}... Ces mots ne servent pleinement que lorsqu'on les met en situation. Un citoyen, c'est quelqu'un qui \emph{parle} : il pose des questions à un élu (interview), il écrit à la mairie (lettre formelle), il défend ses idées devant ses camarades (discours). C'est exactement ce que l'examen attend de vous : produire et traduire ces trois types de textes. Nous allons les étudier un par un, avec leur structure, leurs formules fixes et leurs pièges.

\courssub{L'interview : poser des questions comme un journaliste}

\textbf{Observation.} Lisez ce court extrait d'interview réalisée par un élève de Première pour le journal du lycée :

\begin{textebox}[Interview mit dem Bürgermeister — Auszug]
\textbf{Journalist :} Guten Tag, Herr Bürgermeister. Danke, dass Sie Zeit für uns haben. Darf ich Ihnen einige Fragen stellen ?

\textbf{Bürgermeister :} Guten Tag. Ja, natürlich, gern.

\textbf{Journalist :} Wie finden Sie die Idee eines Jugendzentrums in unserer Stadt ?

\textbf{Bürgermeister :} Das ist eine sehr gute Idee. Die Jugendlichen brauchen einen Ort, wo sie sich treffen können.

\textbf{Journalist :} Können Sie uns erklären, warum das Projekt so lange dauert ?

\textbf{Bürgermeister :} Wir haben leider wenig Geld. Aber wir arbeiten daran.

\textbf{Journalist :} Und was denken Sie über die Mitarbeit der Schüler ?

\textbf{Bürgermeister :} Die Schüler können uns sehr helfen, zum Beispiel bei der Planung.

\textbf{Journalist :} Vielen Dank für das Gespräch, Herr Bürgermeister. Auf Wiedersehen !
\end{textebox}

\begin{definition}[Glossaire]
der Bürgermeister, - (le maire) ; das Jugendzentrum, -zentren (le foyer des jeunes) ; sich treffen (se retrouver) ; dauern (durer) ; die Mitarbeit (la collaboration) ; die Planung (la planification) ; das Gespräch, -e (l'entretien, la conversation).
\end{definition}

\textbf{Ce que nous observons :} l'interview suit toujours le même chemin : on salue, on pose des questions ouvertes (qui commencent par un mot interrogatif en \all{W-}), on relance, on remercie, on prend congé. Et le journaliste vouvoie systématiquement son invité : \all{Sie}, \all{Ihnen}, \all{Ihre}.

\begin{propriete}[La structure de l'interview en cinq étapes]
\begin{enumerate}
\item \textbf{Begrüßung} (salutation) : \all{Guten Tag, Frau Ministerin !} / \all{Hallo, Paul !}
\item \textbf{Einleitung} (présentation) : \all{Ich bin Reporter für unsere Schülerzeitung. Darf ich Ihnen einige Fragen stellen ?} « Je suis reporter pour notre journal lycéen. Puis-je vous poser quelques questions ? »
\item \textbf{Fragen und Nachfragen} (questions et relances) : questions ouvertes en \all{W-}, puis relances : \all{Und wie genau...?}, \all{Können Sie ein Beispiel geben ?} « Pouvez-vous donner un exemple ? »
\item \textbf{Dank} (remerciement) : \all{Vielen Dank für das Gespräch / für Ihre Antworten.}
\item \textbf{Abschied} (congé) : \all{Auf Wiedersehen !} / \all{Ich wünsche Ihnen alles Gute !} « Je vous souhaite tout le meilleur ! »
\end{enumerate}
\end{propriete}

\begin{propriete}[Les formules de questions de l'interview]
Trois formules à connaître par cœur :
\begin{itemize}
\item \all{Wie finden Sie...?} (Que pensez-vous de... ? — litt. « comment trouvez-vous ») : \all{Wie finden Sie das neue Sportplatzprojekt ?} « Que pensez-vous du nouveau projet de terrain de sport ? »
\item \all{Was denken Sie über...?} (+ accusatif) : \all{Was denken Sie über das Engagement der Jugendlichen ?} « Que pensez-vous de l'engagement des jeunes ? »
\item \all{Können Sie uns erklären, warum...?} (phrase subordonnée, verbe à la fin) : \all{Können Sie uns erklären, warum die Wahlen so wichtig sind ?} « Pouvez-vous nous expliquer pourquoi les élections sont si importantes ? »
\end{itemize}
Attention à la subordonnée introduite par \all{warum} : le verbe conjugué va \emph{à la fin} : \all{..., warum das Projekt so lange \textbf{dauert}.}
\end{propriete}

\begin{propriete}[Le vouvoiement : Sie, Ihnen, Ihr — toujours avec majuscule]
Dans un interview formel (avec un adulte, un élu, un inconnu), on utilise obligatoirement le \all{Sie} de politesse. Ses formes à tous les cas :
\begin{center}
\begin{tabular}{|l|l|l|}
\hline
\rowcolor{popblueL} Fonction & Forme & Exemple \\
\hline
sujet (nominatif) & Sie & \all{Sie sind Bürgermeister.} « Vous êtes maire. » \\
\hline
COD (accusatif) & Sie & \all{Ich sehe Sie oft im Fernsehen.} « Je vous vois souvent à la télévision. » \\
\hline
COI (datif) & Ihnen & \all{Ich danke Ihnen.} « Je vous remercie. » \\
\hline
possessif & Ihr / Ihre / Ihr & \all{Ihre Antwort, Ihr Projekt, Ihre Ideen} \\
\hline
\end{tabular}
\end{center}
Piège : \all{danken} se construit avec le \textbf{datif} : on dit \all{Ich danke \textbf{Ihnen}} (je vous remercie), jamais \all{ich danke Sie}. Même chose avec \all{antworten} : \all{Ich antworte Ihnen.} « Je vous réponds. »
\end{propriete}

\begin{attention}[La majuscule du Sie de politesse]
En allemand, \all{Sie}, \all{Ihnen}, \all{Ihr(e)} de politesse prennent \textbf{toujours la majuscule}, même au milieu de la phrase : \all{Ich danke \textbf{Ihnen} für \textbf{Ihre} Antwort.} « Je vous remercie de votre réponse. » Ne confondez pas avec \all{sie} (minuscule) qui signifie « elle » ou « ils/elles » : \all{sie kommt} = « elle vient », mais \all{Sie kommen} = « vous venez ». La majuscule change le sens !
\end{attention}

\begin{popfigure}
\begin{center}
\begin{tikzpicture}[node distance=0.45cm,
etape/.style={rectangle, rounded corners=3pt, draw=popblue, fill=popblueL, align=center, minimum width=5.2cm, minimum height=0.85cm, font=\small},
fleche/.style={-{Stealth[length=2.5mm]}, thick, popdark}]
\node[etape, fill=poporangeL, draw=poporange] (beg) {\textbf{1. Begrüßung} — salutation};
\node[etape, below=of beg] (fra) {\textbf{2. Fragen} — questions en W-};
\node[etape, below=of fra] (nach) {\textbf{3. Nachfragen} — relances};
\node[etape, below=of nach] (dank) {\textbf{4. Dank} — remerciement};
\node[etape, below=of dank] (abs) {\textbf{5. Abschied} — prise de congé};
\draw[fleche] (beg) -- (fra);
\draw[fleche] (fra) -- (nach);
\draw[fleche] (nach) -- (dank);
\draw[fleche] (dank) -- (abs);
\end{tikzpicture}
\end{center}
\legende{Organigramme de l'interview : Begrüßung → Fragen → Nachfragen → Dank → Abschied}
\end{popfigure}

\courssub{La lettre formelle : écrire au maire, au ministre, à un journal}

\textbf{Observation.} Voici le début d'une lettre écrite par des élèves à leur maire :

\begin{exemplebox}[Exemple traduit]
\all{Sehr geehrter Herr Bürgermeister,}\\
\all{ich schreibe Ihnen, weil wir einen Sportplatz in unserem Viertel brauchen.}\\
« Monsieur le Maire, je vous écris parce que nous avons besoin d'un terrain de sport dans notre quartier. »
\end{exemplebox}

Remarquez deux choses : l'appel \all{Sehr geehrter Herr Bürgermeister} (litt. « très honoré Monsieur le Maire ») et la première ligne du corps qui commence par une \emph{minuscule}.

\begin{propriete}[La structure de la lettre formelle allemande]
\begin{enumerate}
\item \textbf{Absender} (expéditeur) : nom et adresse de l'auteur, en haut à gauche.
\item \textbf{Empfänger} (destinataire) : nom et adresse du destinataire, en dessous. Sur l'enveloppe, le titre prend l'accusatif : \all{Herrn Bürgermeister Thomas Nkoulou}.
\item \textbf{Ort und Datum} (lieu et date) : à droite : \all{Yaoundé, den 15. März 2025} (attention : \all{den} + date ordinale avec un point : \all{15.} = quinzième).
\item \textbf{Betreff} (objet) : en gras, sans article : \all{\textbf{Betreff: Bau eines Sportplatzes}}.
\item \textbf{Anrede} (appel) : \all{Sehr geehrter Herr Bürgermeister,} / \all{Sehr geehrte Frau Ministerin,} — suivi d'une \textbf{virgule}.
\item \textbf{Text} (corps) : paragraphes courts, connecteurs logiques.
\item \textbf{Gruß} (formule de politesse) : \all{Mit freundlichen Grüßen} (sans virgule après !), puis la signature.
\end{enumerate}
\end{propriete}

\begin{attention}[La minuscule après l'appel]
En français, on écrit « Monsieur le Maire, \textbf{J}e vous écris... » avec une majuscule. En allemand, c'est l'inverse : la première ligne après l'appel commence par une \textbf{minuscule} (sauf si c'est un nom) : \all{Sehr geehrte Frau Ministerin,} / \all{\textbf{i}ch möchte Ihnen schreiben, weil...} C'est une erreur classique de francophone à l'examen !
\end{attention}

\begin{popfigure}
\begin{center}
\begin{tikzpicture}[zone/.style={rectangle, rounded corners=2pt, draw=popdark, fill=popcream, align=left, font=\small, inner sep=4pt},
etiq/.style={font=\footnotesize\itshape, text=popmuted, align=left}]
\node[zone, anchor=north west, minimum width=4.6cm, align=center] (abs) at (0,0){Amina Mbarga\\ Quartier Nkolbisson\\ Yaoundé};
\node[etiq, right=0.15cm of abs.east] {Absender};
\node[zone, anchor=north west, minimum width=4.6cm, below=0.35cm of abs, align=center] (emp){Herrn Bürgermeister\\ Rathaus der Stadt\\ Yaoundé};
\node[etiq, right=0.15cm of emp.east] {Empfänger};
\node[zone, anchor=north west, minimum width=5.2cm, below=0.35cm of emp] (dat) {Yaoundé, den 15. März 2025};
\node[etiq, right=0.15cm of dat.east] {Ort und Datum};
\node[zone, anchor=north west, minimum width=8cm, below=0.35cm of dat] (bet) {\textbf{Betreff: Bau eines Sportplatzes}};
\node[etiq, right=0.15cm of bet.east] {Betreff};
\node[zone, anchor=north west, minimum width=8cm, below=0.3cm of bet] (anr) {Sehr geehrter Herr Bürgermeister,};
\node[etiq, right=0.15cm of anr.east] {Anrede};
\node[zone, anchor=north west, minimum width=11.3cm, minimum height=1.6cm, below=0.3cm of anr, align=center] (txt){ich schreibe Ihnen, weil wir einen Sportplatz\\ brauchen. Erstens ... Außerdem ... Deshalb ...};
\node[etiq, right=0.15cm of txt.east] {Text};
\node[zone, anchor=north west, minimum width=5.5cm, below=0.3cm of txt, align=center] (gru){Mit freundlichen Grüßen\\ Amina Mbarga};
\node[etiq, right=0.15cm of gru.east] {Gruß};
\end{tikzpicture}
\end{center}
\legende{Schéma de la lettre formelle : les sept zones obligatoires}
\end{popfigure}

\courssub{Le petit discours : convaincre son public}

\textbf{Observation.} Écoutez ce que dit une candidate aux élections des délégués de classe :

\begin{exemplebox}[Exemple traduit]
\all{Wählt mich, denn ich werde für eure Rechte kämpfen !}\\
« Élisez-moi, car je me battrai pour vos droits ! »
\end{exemplebox}

On remarque : l'impératif \all{Wählt mich !} (élisez-moi — ici à la 2\textsuperscript{e} personne du pluriel \all{ihr}, car elle s'adresse à ses camarades), la proposition avec \all{denn} (car — verbe en position 2, car \all{denn} est une \emph{coordination}, pas une subordination !), et le futur \all{ich werde kämpfen} pour promettre.

\begin{propriete}[La structure du petit discours]
\begin{enumerate}
\item \textbf{Accroche} : \all{Liebe Mitschülerinnen und Mitschüler !} « Chers camarades ! » (formule inclusive obligatoire en allemand moderne).
\item \textbf{Présentation du problème} : \all{In unserer Schule gibt es ein Problem: ...} « Dans notre école, il y a un problème : ... »
\item \textbf{Arguments} avec connecteurs : \all{Erstens ... Zweitens ... Außerdem ...}
\item \textbf{Appel à l'action} : \all{Wählt mich !} / \all{Kommt alle zur Wahl !} « Venez tous voter ! »
\item \textbf{Remerciement} : \all{Danke für eure Aufmerksamkeit !} « Merci de votre attention ! »
\end{enumerate}
\end{propriete}

\begin{methode}[Plan du discours électoral en 5 étapes]
\begin{enumerate}
\item \textbf{Saluer} le public : \all{Liebe Mitschülerinnen und Mitschüler !}
\item \textbf{Se présenter} : \all{Ich heiße Amina und ich möchte eure Klassensprecherin werden.} « Je m'appelle Amina et je voudrais devenir votre déléguée de classe. »
\item \textbf{Exposer 2 à 3 idées} avec des connecteurs : \all{Erstens brauchen wir mehr Bücher. Zweitens will ich einen Sporttag organisieren. Außerdem höre ich euch immer zu.} « D'abord, nous avons besoin de plus de livres. Ensuite, je veux organiser une journée du sport. De plus, je vous écoute toujours. »
\item \textbf{Appeler à voter} : \all{Wählt mich, denn ich werde für eure Rechte kämpfen !}
\item \textbf{Remercier} : \all{Vielen Dank für eure Aufmerksamkeit !}
\end{enumerate}
\end{methode}

\begin{propriete}[Les connecteurs logiques du discours et de la lettre]
\begin{center}
\begin{tabularx}{\linewidth}{|l|X|X|}
\hline
\rowcolor{popblueL} Deutsch & Français & Position du verbe \\
\hline
\all{zuerst} & d'abord & verbe en 2\textsuperscript{e} position : \all{Zuerst \textbf{danke} ich...} \\
\hline
\all{dann} & ensuite & idem : \all{Dann \textbf{kommt} Paul.} \\
\hline
\all{außerdem} & de plus & idem : \all{Außerdem \textbf{brauchen} wir Geld.} \\
\hline
\all{jedoch} & cependant & idem : \all{Wir haben jedoch wenig Geld.} \\
\hline
\all{deshalb} & c'est pourquoi & verbe en 2\textsuperscript{e} position : \all{Deshalb \textbf{schreibe} ich Ihnen.} \\
\hline
\all{schließlich} & finalement & idem : \all{Schließlich \textbf{bitte} ich euch.} \\
\hline
\all{meiner Meinung nach} & à mon avis & \all{Meiner Meinung nach \textbf{ist} das wichtig.} \\
\hline
\end{tabularx}
\end{center}
Règle d'or : en allemand, le verbe conjugué est \textbf{toujours en deuxième position}. Si le connecteur ouvre la phrase, le sujet passe \emph{après} le verbe : \all{Deshalb schreibe \textbf{ich}} (et non \all{deshalb ich schreibe}).
\end{propriete}

\begin{propriete}[Registre de langue : formel ou courant ?]
\begin{center}
\begin{tabularx}{\linewidth}{|X|X|X|}
\hline
\rowcolor{popblueL} Situation & Registre & Marques \\
\hline
Lettre au maire, discours officiel & formel & \all{Sie/Ihnen/Ihr}, \all{Sehr geehrte...}, \all{Mit freundlichen Grüßen} \\
\hline
Interview avec un élu & formel & \all{Sie}, \all{Guten Tag}, \all{Vielen Dank für das Gespräch} \\
\hline
Interview entre jeunes, discours aux camarades & courant & \all{du/ihr}, \all{Hallo !}, \all{Liebe Mitschüler !}, \all{Tschüss !} \\
\hline
\end{tabularx}
\end{center}
Le discours électoral devant la classe est un cas mixte : accroche courante (\all{Liebe Mitschüler !}, \all{eure}, \all{Wählt mich !}) mais arguments structurés comme un texte formel.
\end{propriete}

\begin{textebox}[Eine Rede für die Schülerwahl — Modelltext]
Liebe Mitschülerinnen und Mitschüler !

Ich heiße Etienne und ich möchte euer Klassensprecher werden. Warum ? Ich erkläre es euch.

In unserer Schule gibt es Probleme. Erstens haben wir zu wenig Bücher in der Bibliothek. Viele Schüler können nicht richtig lernen. Zweitens gibt es keinen Sportplatz, und Sport ist wichtig für die Gesundheit. Außerdem hören die Erwachsenen uns oft nicht zu. Meiner Meinung nach haben auch Jugendliche Rechte — und Ideen !

Was will ich machen ? Zuerst will ich mit dem Direktor über die Bibliothek sprechen. Dann organisiere ich eine Unterschriftenliste für den Sportplatz. Schließlich will ich jeden Monat ein Treffen mit euch machen : Ihr sagt mir eure Probleme, und ich trage sie weiter.

Ich weiß : Das ist nicht einfach. Aber ich bin motiviert, und ich habe Zeit für euch. Deshalb bitte ich euch : Wählt mich, denn ich werde für eure Rechte kämpfen ! Zusammen sind wir stark.

Vielen Dank für eure Aufmerksamkeit !
\end{textebox}

\begin{definition}[Glossaire du discours]
der Klassensprecher, - (le délégué de classe) ; die Bibliothek, -en (la bibliothèque) ; die Gesundheit (la santé) ; die Erwachsenen (les adultes) ; die Unterschriftenliste, -n (la pétition) ; das Treffen, - (la réunion) ; etwas weitertragen (transmettre quelque chose) ; motiviert (motivé) ; zusammen (ensemble).
\end{definition}

\textbf{Questions de compréhension :} 1) Quels sont les trois problèmes cités par Etienne ? 2) Quelles sont ses trois actions promises ? 3) Relevez tous les connecteurs du texte. 4) Pourquoi utilise-t-il \all{euch}, \all{ihr}, \all{eure} et non \all{Ihnen}, \all{Sie} ?

\begin{exoresolu}[Rédiger l'appel et la politesse d'une lettre au maire]
Énoncé : Vous écrivez à la mairesse de Douala pour demander une bibliothèque dans votre quartier. Rédigez en allemand : l'objet (Betreff), l'appel, la première phrase du corps (avec \all{weil}) et la formule de politesse.
\tcblower
Solution détaillée :
\begin{itemize}
\item Betreff : \all{\textbf{Betreff: Bau einer Bibliothek im Quartier Bonabéri}} (pas d'article après « Betreff: »).
\item Anrede : \all{Sehr geehrte Frau Bürgermeisterin,} (féminin : \all{geehrte}, \all{Frau Bürgermeisterin}), suivie d'une virgule.
\item Première phrase, minuscule obligatoire, subordonnée \all{weil} avec verbe à la fin : \all{ich schreibe Ihnen, weil wir eine Bibliothek in unserem Quartier brauchen.}
\item Formule finale : \all{Mit freundlichen Grüßen} (sans virgule), puis la signature : \all{Chantal Ewusi}.
\end{itemize}
Points de vigilance : majuscule à \all{Ihnen}, minuscule en début de corps, verbe \all{brauchen} à la fin de la subordonnée.
\end{exoresolu}

\begin{attention}[Les trois pièges du francophone]
\begin{enumerate}
\item \textbf{L'ordre des mots après un connecteur} : ne dites jamais \all{Deshalb ich schreibe Ihnen} mais \all{Deshalb \textbf{schreibe ich} Ihnen} (verbe en 2\textsuperscript{e} position, sujet après).
\item \textbf{La virgule et la majuscule françaises} : après \all{Sehr geehrter Herr Bürgermeister,} on met une virgule et on continue en \emph{minuscule} ; après \all{Mit freundlichen Grüßen} on ne met \emph{pas} de virgule.
\item \textbf{Le tutoiement intempestif} : dans une lettre ou un interview formel, un seul \all{du} peut coûter des points. Vérifiez chaque pronom : \all{Sie}, \all{Ihnen}, \all{Ihr} avec majuscule.
\end{enumerate}
\end{attention}

\begin{savaistu}[Voter dès 16 ans en Autriche]
En Autriche, le droit de vote est acquis dès \textbf{16 ans} pour toutes les élections — une première en Europe. L'idée : celui qui s'intéresse tôt à la politique reste un citoyen actif toute sa vie. En Allemagne, on vote à 18 ans pour les élections nationales, mais dès 16 ans dans certains Länder pour les élections locales et européennes. Et vous, à quel âge devrait-on voter, selon vous ? \all{Meiner Meinung nach...}
\end{savaistu}

\begin{exercice}[À vous de produire]
1) Rédigez trois questions d'interview (avec \all{Wie finden Sie...?}, \all{Was denken Sie über...?}, \all{Können Sie uns erklären, warum...?}) pour interviewer le maire de Yaoundé sur la propreté de la ville.\\
2) Écrivez le Betreff, l'Anrede et la première phrase d'une lettre au ministre de la Jeunesse pour demander un centre culturel.\\
3) Rédigez l'accroche et l'appel à l'action d'un discours pour l'élection du délégué de votre classe, avec au moins deux connecteurs.
\end{exercice}

\begin{corrige}[Exercice « À vous de produire »]
1) Exemples : \all{Wie finden Sie die Situation in unserer Stadt ?} ; \all{Was denken Sie über die Arbeit der Jugendlichen ?} ; \all{Können Sie uns erklären, warum die Straßen oft schmutzig sind ?}\\
2) \all{\textbf{Betreff: Bau eines Kulturzentrums}} / \all{Sehr geehrter Herr Minister,} / \all{ich schreibe Ihnen, weil die Jugendlichen in unserer Stadt ein Kulturzentrum brauchen.}\\
3) \all{Liebe Mitschülerinnen und Mitschüler ! Ich heiße ... und ich möchte euer Klassensprecher werden. Erstens will ich euch immer zuhören, außerdem werde ich eure Ideen weitertragen. Deshalb : Wählt mich ! Danke für eure Aufmerksamkeit !}
\end{corrige}

\begin{aretenir}[L'essentiel de la section]
Interview en 5 étapes : \all{Begrüßung → Fragen → Nachfragen → Dank → Abschied}. Lettre formelle en 7 zones : \all{Absender, Empfänger, Datum, Betreff, Anrede, Text, Gruß}. Discours en 5 temps : saluer, se présenter, argumenter avec connecteurs, appeler à l'action, remercier. Trois règles d'or : \all{Sie/Ihnen/Ihr} avec majuscule ; minuscule après l'appel ; verbe conjugué toujours en deuxième position après un connecteur.
\end{aretenir}

\coursec{Révision grammaticale pour la traduction}

Dans la section précédente, vous avez appris à produire trois types de textes : l'interview, la lettre et le petit discours. Pour réussir ces productions — et surtout pour traduire correctement dans les deux sens —, il faut maintenant réviser les grandes structures grammaticales du programme de Première. Cette section est votre boîte à outils : chaque règle est d'abord observée sur des exemples, puis énoncée rigoureusement, comparée au français, illustrée et entraînée. Gardez-la sous les yeux chaque fois que vous traduisez ou rédigez.

\courssub{L'ordre des mots dans la phrase principale : le verbe en position 2}

\begin{exemplebox}[Observer avant de comprendre]
Comparez ces deux phrases allemandes et leur traduction :

\all{Ich gehe heute in die Schule.} « Je vais aujourd'hui à l'école. »

\all{Heute gehe ich in die Schule.} « Aujourd'hui, je vais à l'école. »

Que remarquez-vous ? Dans les deux cas, le verbe conjugué \all{gehe} occupe exactement la \cle{deuxième place} de la phrase. Quand un complément (\all{heute}) passe en tête, le sujet et le verbe \textbf{s'inversent} : \all{heute gehe ich}, et non \all{*heute ich gehe}.
\end{exemplebox}

\begin{propriete}[La règle d'or : V2]
Dans toute phrase principale affirmative allemande, le \cle{verbe conjugué} occupe la \textbf{deuxième position}. La première position peut être occupée par le sujet, un complément de temps, de lieu ou tout autre élément — mais jamais par deux éléments à la fois. Si la première place n'est pas le sujet, le sujet vient \textbf{immédiatement après} le verbe (inversion sujet-verbe).

\all{Morgen fährt mein Vater nach Douala.} « Demain, mon père part pour Douala. »

\all{Nach der Schule spielen wir Fußball.} « Après l'école, nous jouons au football. »
\end{propriete}

\begin{attention}[Erreur typique du francophone]
En français on dit « Hier, je suis allé au marché ». Ne traduisez jamais \all{*Gestern ich bin zum Markt gegangen}. Dès qu'un élément autre que le sujet ouvre la phrase, le verbe reste en position 2 et le sujet suit : \all{Gestern bin ich zum Markt gegangen.} Retenez : \textbf{un seul élément avant le verbe}.
\end{attention}

\courssub{La phrase à cadre (Satzklammer)}

\begin{exemplebox}[Observer]
\all{Ich habe gestern einen Brief geschrieben.} « J'ai écrit une lettre hier. »

\all{Wir müssen die Gesetze respektieren.} « Nous devons respecter les lois. »

Dans les deux phrases, le verbe se dédouble : une partie conjuguée en position 2 (\all{habe}, \all{müssen}) et une partie non conjuguée — participe II (\all{geschrieben}) ou infinitif (\all{respektieren}) — rejetée \textbf{à la toute fin}. Les deux parties forment un « cadre » qui entoure le reste de la phrase.
\end{exemplebox}

\begin{propriete}[Le cadre verbal]
Au Perfekt, au Futur, au passif et avec les verbes modaux, la phrase principale allemande forme un \cle{cadre} : verbe conjugué en \textbf{position 2} + participe II ou infinitif en \textbf{dernière position}. Tous les compléments se placent à l'intérieur du cadre.

\all{Sie hat ihre Pflichten als Bürgerin erfüllt.} « Elle a accompli ses devoirs de citoyenne. »

\all{Der Bürgermeister wird morgen eine Rede halten.} « Le maire fera un discours demain. »
\end{propriete}

\begin{popfigure}
\begin{center}
\begin{tikzpicture}[font=\small]
\node[draw=popblue, fill=popblueL, rounded corners, align=center] (v2) at (0,0) {\textbf{V2}\\ \all{habe}};
\node[draw=popmuted, fill=popcream, rounded corners, align=center] (mitte) at (3.2,0) {compléments\\ \all{gestern einen Brief}};
\node[draw=poporange, fill=poporangeL, rounded corners, align=center] (ende) at (7,0) {\textbf{fin de phrase}\\ \all{geschrieben}};
\draw[-{Stealth}, very thick, poppurple, bend left=35] (v2.north) to node[above, align=center]{le cadre se referme} (ende.north);
\node[align=center, popdark] at (3.5,-1.3) {\all{Ich \textbf{habe} gestern einen Brief \textbf{geschrieben}.}};
\end{tikzpicture}
\end{center}
\legende{La phrase à cadre : le verbe conjugué ouvre, le participe ou l'infinitif ferme.}
\end{popfigure}

\begin{attention}[Ne rien placer après le cadre]
En français, « hier » peut terminer la phrase. En allemand, \textbf{rien} ne peut se placer après le participe ou l'infinitif final (sauf une subordonnée). Faux : \all{*Ich habe geschrieben gestern einen Brief.} Juste : \all{Ich habe gestern einen Brief geschrieben.}
\end{attention}

\courssub{Les subordonnées : le verbe conjugué à la fin}

\begin{exemplebox}[Observer]
\all{Ich lerne Deutsch, weil ich in Deutschland studieren möchte.} « J'apprends l'allemand parce que je veux étudier en Allemagne. »

\all{Der Journalist sagt, dass die Jugendlichen sehr engagiert sind.} « Le journaliste dit que les jeunes sont très engagés. »

Dans chaque subordonnée, le verbe conjugué (\all{möchte}, \all{sind}) est rejeté \textbf{à la dernière place}.
\end{exemplebox}

\begin{propriete}[Verbe final dans les subordonnées]
Dans une subordonnée introduite par \all{weil} (parce que), \all{dass} (que), \all{wenn} (quand, si), \all{obwohl} (bien que), \all{damit} (afin que) ou \all{als} (quand, au passé), le verbe conjugué va à la \cle{dernière place} de la subordonnée. La subordonnée est toujours séparée de la principale par une \textbf{virgule}.

\all{..., weil ich meine Pflichten kenne.} « ..., parce que je connais mes devoirs. »

\all{Obwohl er jung ist, engagiert er sich politisch.} « Bien qu'il soit jeune, il s'engage politiquement. »

\all{Wir wählen, damit unsere Stimme zählt.} « Nous votons afin que notre voix compte. »
\end{propriete}

\begin{center}
\begin{tabularx}{\linewidth}{|X|X|X|}
\hline
\rowcolor{popblueL} Conjonction & Sens & Exemple \\
\hline
\all{weil} & parce que & \all{..., weil ich wählen will.} \\
\hline
\all{dass} & que & \all{Ich denke, dass Wahlen wichtig sind.} \\
\hline
\all{wenn} & quand, si (présent, futur, répétition) & \all{Wenn ich 18 bin, wähle ich.} \\
\hline
\all{obwohl} & bien que & \all{..., obwohl es regnet.} \\
\hline
\all{damit} & afin que & \all{..., damit alle verstehen.} \\
\hline
\all{als} & quand (passé ponctuel) & \all{Als ich Kind war, ...} \\
\hline
\end{tabularx}
\end{center}

\begin{attention}[\all{als} ou \all{wenn} ?]
Ne confondez jamais \all{als} et \all{wenn} pour traduire « quand ». \all{als} s'emploie pour un événement \textbf{ponctuel au passé} : \all{Als ich Kind war, wohnten wir in Yaoundé.} « Quand j'étais enfant, nous habitions à Yaoundé. » \all{wenn} s'emploie au présent, au futur et pour la \textbf{répétition} (chaque fois que) : \all{Wenn ich Zeit habe, helfe ich meinen Eltern.} « Quand j'ai le temps, j'aide mes parents. »
\end{attention}

\begin{popfigure}
\begin{center}
\begin{tabular}{|l|l|}
\hline
\rowcolor{popgreenL} \textbf{Principale (V2)} & \textbf{Subordonnée (verbe final)} \\
\hline
\all{\textcolor{popblue}{Ich kenne} meine Pflichten.} & \all{..., weil ich meine Pflichten \textcolor{poporange}{kenne}.} \\
\hline
\all{\textcolor{popblue}{Er hat} gestern gewählt.} & \all{..., dass er gestern \textcolor{poporange}{gewählt hat}.} \\
\hline
\all{\textcolor{popblue}{Wir müssen} die Gesetze achten.} & \all{..., weil wir die Gesetze \textcolor{poporange}{achten müssen}.} \\
\hline
\end{tabular}
\end{center}
\legende{Principale : verbe conjugué en position 2. Subordonnée : verbe conjugué à la fin.}
\end{popfigure}

Remarquez la deuxième ligne : au Perfekt dans une subordonnée, c'est l'auxiliaire conjugué (\all{hat}) qui va en dernier, \textbf{après} le participe : \all{..., dass er gestern gewählt hat.} De même, avec un modal (troisième ligne), l'ordre final est : infinitif + modal conjugué : \all{..., achten müssen}.

\courssub{Les pronoms relatifs : der, die, das, den, dem}

\begin{exemplebox}[Observer]
\all{Der Kandidat, den wir gewählt haben, ist jung.} « Le candidat que nous avons élu est jeune. »

\all{Die Frau, die neben mir wohnt, ist Lehrerin.} « La femme qui habite à côté de chez moi est enseignante. »

\all{Das Buch, das ich lese, ist interessant.} « Le livre que je lis est intéressant. »
\end{exemplebox}

\begin{propriete}[Accord du pronom relatif]
Le pronom relatif allemand s'accorde en \cle{genre} et en \cle{nombre} avec son antécédent (le nom auquel il se rapporte), et prend le \cle{cas} de sa fonction dans la subordonnée relative. La relative est une subordonnée : le verbe conjugué va à la fin, et la relative est encadrée de virgules.
\end{propriete}

\begin{center}
\begin{tabular}{|l|l|l|l|}
\hline
\rowcolor{popblueL} Cas & Masculin & Féminin & Neutre \\
\hline
Nominatif (sujet) & \all{der} & \all{die} & \all{das} \\
\hline
Accusatif (COD) & \all{den} & \all{die} & \all{das} \\
\hline
Datif & \all{dem} & \all{der} & \all{dem} \\
\hline
Génitif & \all{dessen} & \all{deren} & \all{dessen} \\
\hline
\end{tabular}
\end{center}

\begin{exemplebox}[Exemples traduits]
\all{Der Mann, der dort spricht, ist unser Bürgermeister.} « L'homme qui parle là-bas est notre maire. » (nominatif masculin : \all{der})

\all{Der Kandidat, den wir gewählt haben, ist jung.} « Le candidat que nous avons élu est jeune. » (accusatif masculin : \all{den})

\all{Die Bürgerin, der ich geholfen habe, dankt mir.} « La citoyenne que j'ai aidée me remercie. » (datif féminin : \all{der}, car \all{helfen} + datif)
\end{exemplebox}

\begin{attention}[Pas de « que » universel en allemand]
Le français utilise « que » partout ; l'allemand exige le bon genre et le bon cas. Avant d'écrire le relatif, posez-vous deux questions : quel est le genre de l'antécédent ? quelle est la fonction du relatif dans sa propre phrase (sujet, COD, complément d'un verbe à datif) ?
\end{attention}

\courssub{Les cas après les prépositions}

\begin{propriete}[Prépositions à cas fixe]
\begin{itemize}
\item \textbf{Accusatif} : \all{für} (pour), \all{gegen} (contre), \all{ohne} (sans), \all{durch} (à travers), \all{um} (autour de). Ex. : \all{Wir kämpfen für unsere Rechte.} « Nous luttons pour nos droits. »
\item \textbf{Datif} : \all{mit} (avec), \all{zu} (à, vers), \all{bei} (chez), \all{von} (de), \all{nach} (après, vers), \all{aus} (de), \all{seit} (depuis). Ex. : \all{Ich spreche mit dem Journalisten.} « Je parle avec le journaliste. »
\end{itemize}
\end{propriete}

\begin{propriete}[Prépositions à deux cas : in, an, auf]
Les prépositions \all{in}, \all{an}, \all{auf}, \all{über}, \all{unter}, \all{vor}, \all{hinter}, \all{neben}, \all{zwischen} se construisent avec le \textbf{datif} si l'on répond à la question \all{wo ?} (où ? — lieu, position), et avec l'\textbf{accusatif} si l'on répond à \all{wohin ?} (vers où ? — mouvement).

\all{Wir sind in der Schule.} « Nous sommes à l'école. » (où ? → datif)

\all{Wir gehen in die Schule.} « Nous allons à l'école. » (vers où ? → accusatif)

\all{Das Plakat hängt an der Wand.} « L'affiche est accrochée au mur. » / \all{Ich hänge das Plakat an die Wand.} « J'accroche l'affiche au mur. »
\end{propriete}

\begin{center}
\begin{tabularx}{\linewidth}{|X|X|X|}
\hline
\rowcolor{popblueL} Français & Deutsch & Remarque \\
\hline
à l'école (on y est) & \all{in der Schule} & datif : lieu \\
\hline
à l'école (on y va) & \all{in die Schule} & accusatif : mouvement \\
\hline
avec le maire & \all{mit dem Bürgermeister} & datif après \all{mit} \\
\hline
pour la patrie & \all{für das Vaterland} & accusatif après \all{für} \\
\hline
chez le médecin & \all{beim Arzt} & \all{bei dem} → \all{beim} \\
\hline
\end{tabularx}
\end{center}

\courssub{Le choix des temps en traduction}

\begin{propriete}[Perfekt, Präteritum, Futur I]
\begin{itemize}
\item \textbf{Perfekt} (\all{haben/sein} + participe II) : temps du passé à l'\textbf{oral} et dans les lettres. \all{Ich habe einen Brief geschrieben.} « J'ai écrit une lettre. » Auxiliaire \all{sein} pour les verbes de mouvement et de changement d'état : \all{Wir sind nach Douala gefahren.} « Nous sommes allés à Douala. »
\item \textbf{Präteritum} : temps du passé à l'\textbf{écrit} (récits, articles), et \textbf{toujours} pour \all{sein} (\all{war}), \all{haben} (\all{hatte}) et les modaux (\all{musste}, \all{wollte}, \all{konnte}), même à l'oral. \all{Als ich Kind war, musste ich früh aufstehen.} « Quand j'étais enfant, je devais me lever tôt. »
\item \textbf{Futur I} : \all{werden} + infinitif en fin de phrase. \all{Ich werde nächstes Jahr wählen.} « Je voterai l'année prochaine. »
\end{itemize}
\end{propriete}

\begin{attention}[Ne dites pas \all{*ich bin gewesen} pour « j'étais »]
Pour traduire « j'étais », « j'avais », « je devais » au passé, utilisez le Präteritum : \all{ich war}, \all{ich hatte}, \all{ich musste}. Le Perfekt de ces verbes (\all{bin gewesen}, \all{habe gehabt}) est lourd et marque un autre sens.
\end{attention}

\courssub{L'infinitif avec \all{zu} et le passif}

\begin{propriete}[L'infinitif avec \all{zu}]
Après certaines tournures (\all{es ist wichtig}, \all{ich hoffe}, \all{versuchen}, \all{vergessen}...), l'infinitif est précédé de \all{zu} et rejeté en fin de proposition, après une virgule. Avec un verbe à particule séparable, \all{zu} s'insère \textbf{entre la particule et le verbe} : \all{aufzustehen}, \all{teilzunehmen}.

\all{Es ist wichtig, die Gesetze zu respektieren.} « Il est important de respecter les lois. »

\all{Ich hoffe, an der Wahl teilzunehmen.} « J'espère participer à l'élection. »
\end{propriete}

\begin{propriete}[Le passif : \all{werden} + participe II]
Le passif se forme avec l'auxiliaire \all{werden} (conjugué) + participe II (en fin de phrase). L'agent est introduit par \all{von} + datif.

\all{Die Gesetze werden respektiert.} « Les lois sont respectées. »

\all{Der Bürgermeister wird von den Bürgern gewählt.} « Le maire est élu par les citoyens. »

Au Perfekt passif : \all{Die Gesetze sind respektiert worden.} « Les lois ont été respectées. » (attention : \all{worden}, pas \all{geworden}).
\end{propriete}

\begin{attention}[Faux amis et tournures à ne pas traduire mot à mot]
Ne traduisez jamais « Je suis d'accord avec toi » mot à mot : dites \all{Ich bin mit dir einverstanden} ou \all{Ich stimme dir zu.} (attention, \all{zustimmen} + datif). De même, « avoir raison » = \all{Recht haben} (\all{Du hast Recht.}), « avoir 18 ans » = \all{18 Jahre alt sein} (\all{Ich bin 18 Jahre alt.}), et non \all{*Ich habe 18 Jahre}.
\end{attention}

\begin{methode}[Traduire une phrase en cinq étapes]
\begin{enumerate}
\item \textbf{Souligner} le verbe français et identifier son temps et sa voix (active ou passive).
\item \textbf{Choisir} le temps allemand : Perfekt (oral, lettre), Präteritum (récit écrit, \all{sein/haben}/modaux), Futur I (\all{werden} + infinitif).
\item \textbf{Placer} le verbe conjugué en position 2 dans la principale ; vérifier l'inversion si un complément ouvre la phrase.
\item \textbf{Construire le cadre} : participe II ou infinitif en dernière position ; verbe final dans les subordonnées ; virgule avant toute subordonnée.
\item \textbf{Relire} : majuscules aux noms, cas après prépositions, genre des relatifs, orthographe (ß, Umlaute).
\end{enumerate}
\end{methode}

\begin{exoresolu}[Thème : traduire en allemand]
Énoncé : Traduisez — « Hier, les jeunes ont discuté avec le maire, parce qu'ils veulent participer à la vie de la cité. »

\tcblower

Solution détaillée :

Étape 1 : deux verbes — « ont discuté » (passé composé, contexte oral/lettre → Perfekt) et « veulent » (présent). « Parce que » introduit une subordonnée → \all{weil} + verbe final, précédée d'une virgule.

Étape 2 : « hier » en tête de phrase → inversion sujet-verbe obligatoire : \all{Gestern haben die Jugendlichen ...}

Étape 3 : « avec le maire » → \all{mit} + datif : \all{mit dem Bürgermeister}. « Participer à » → \all{teilnehmen an} + datif, verbe à particule séparable ; après \all{wollen}, infinitif en fin de subordonnée : \all{teilnehmen wollen}.

Traduction : \all{Gestern haben die Jugendlichen mit dem Bürgermeister diskutiert, weil sie am Leben der Stadt teilnehmen wollen.}

Vérification : verbe \all{haben} en position 2, participe \all{diskutiert} en fin de principale ; dans la subordonnée, l'ordre final est infinitif + modal conjugué : \all{teilnehmen wollen} ; \all{am} = \all{an dem} (datif).
\end{exoresolu}

\begin{exoresolu}[Version : traduire en français]
Énoncé : Traduisez — \all{Obwohl viele Bürger jung sind, werden sie von den Politikern oft nicht gehört.}

\tcblower

Solution : \all{obwohl} = « bien que » (subordonnée, verbe \all{sind} en fin, confirmé). La principale est au passif : \all{werden ... gehört} = « sont entendus » ; \all{von den Politikern} = « par les hommes politiques ». Traduction : « Bien que de nombreux citoyens soient jeunes, ils ne sont souvent pas écoutés par les hommes politiques. » Piège évité : ne pas traduire \all{werden} par « deviendront » ici — c'est l'auxiliaire du passif.
\end{exoresolu}

\begin{exercice}[À vous de jouer]
Traduisez en allemand : 1. « Demain, j'écrirai une lettre au président de la République. » 2. « Le candidat que nous avons élu parle avec les jeunes. » 3. « Quand j'étais enfant, je devais aider mes parents. » 4. « Il est important de respecter les droits des autres. »

Traduisez en français : 5. \all{Die Schüler, die an dem Projekt teilnehmen, werden von der Schule belohnt.}
\end{exercice}

\begin{corrige}[Exercice « À vous de jouer »]
1. \all{Morgen werde ich einen Brief an den Staatspräsidenten schreiben.} (Futur I, cadre \all{werde ... schreiben}, \all{an} + accusatif : destinataire d'un envoi.)

2. \all{Der Kandidat, den wir gewählt haben, spricht mit den Jugendlichen.} (relatif accusatif masculin \all{den} ; \all{mit} + datif pluriel \all{den Jugendlichen}.)

3. \all{Als ich Kind war, musste ich meinen Eltern helfen.} (passé ponctuel → \all{als} ; Präteritum du modal \all{musste} ; \all{helfen} + datif.)

4. \all{Es ist wichtig, die Rechte der anderen zu respektieren.} (infinitif avec \all{zu} en fin, virgule avant l'infinitive.)

5. « Les élèves qui participent au projet sont récompensés par l'école. » (relatif nominatif pluriel \all{die} ; passif \all{werden belohnt}.)
\end{corrige}

\begin{aretenir}[Les cinq réflexes du traducteur]
\begin{itemize}
\item Principale : verbe conjugué en \textbf{position 2}, toujours.
\item Perfekt, Futur, passif, modaux : \textbf{cadre} (participe ou infinitif en fin).
\item Subordonnée (\all{weil}, \all{dass}, \all{wenn}, \all{obwohl}, \all{damit}, \all{als}) : verbe conjugué \textbf{en dernier}, et virgule avant la subordonnée.
\item Relatif : genre de l'antécédent, cas de la fonction.
\item Prépositions : accusatif (\all{für}, \all{gegen}, \all{ohne}), datif (\all{mit}, \all{zu}, \all{bei}, \all{von}, \all{nach}) ; \all{in/an/auf} : datif = lieu, accusatif = mouvement.
\end{itemize}
\end{aretenir}

\coursec{Entraînement à la traduction : thème et version}

Dans la section précédente, nous avons révisé les grandes structures grammaticales du programme : la place du verbe conjugué en deuxième position, le cadre verbal, la subordination, les cas et les temps. Ces connaissances ne sont pas une fin en soi : elles trouvent leur application la plus exigeante dans la \cle{traduction}. À l'examen de fin d'année comme plus tard au baccalauréat, deux exercices vous attendent : le \cle{thème} (traduire du français vers l'allemand) et la \cle{version} (traduire de l'allemand vers le français). Cette section vous donne une méthode rigoureuse, pas à pas, puis vous entraîne sur des phrases et un texte liés à la citoyenneté.

\courssub{Le thème : traduire du français vers l'allemand}

\textbf{Observer avant de traduire}\par

Lisez d'abord ces deux phrases et comparez-les :

\all{Die Jugendlichen müssen am Leben ihres Viertels teilnehmen.} « Les jeunes doivent participer à la vie de leur quartier. »

\all{Wir haben den Bürgermeister interviewt.} « Nous avons interviewé le maire. »

Que remarquez-vous ? Dans la première phrase, le verbe conjugué \all{müssen} est en deuxième position et l'infinitif \all{teilnehmen} est renvoyé \textbf{à la fin}. Dans la seconde, l'auxiliaire \all{haben} est en deuxième position et le participe \all{interviewt} ferme le cadre. Le français ne fonctionne pas ainsi : c'est toute la difficulté du thème. On ne traduit jamais mot à mot ; on \textbf{reconstruit} la phrase selon les lois de l'allemand.

\begin{methode}[Les 5 étapes du thème]
\begin{enumerate}
\item \textbf{Lire toute la phrase} française, jusqu'au point final, sans écrire un seul mot d'allemand. Identifier le sens global.
\item \textbf{Repérer le verbe et son temps} : présent, passé composé, futur ? Cela décide du temps allemand (Präsens, Perfekt, Futur I).
\item \textbf{Choisir la construction} : phrase principale (verbe en position 2, cadre verbal) ou subordonnée (verbe conjugué à la fin) ? Y a-t-il une relative, une infinitive, un modal ?
\item \textbf{Traduire} en plaçant d'abord le sujet et le verbe conjugué, puis en fermant le cadre avec l'infinitif, le participe ou la particule séparable.
\item \textbf{Relire} en vérifiant trois points : la place du verbe, les cas après les prépositions et les verbes, les majuscules à tous les noms.
\end{enumerate}
\end{methode}

\begin{popfigure}
\begin{tikzpicture}[
  node distance=0.55cm and 0.9cm,
  box/.style={draw=popblue, fill=popblueL, rounded corners, thick, align=center, text width=3.1cm, minimum height=1.1cm, font=\small},
  fleche/.style={-{Stealth[length=3mm]}, thick, poporange}
]
\node[box, fill=poporangeL, draw=poporange, align=center] (depart){Phrase française\\ (sens global)};
\node[box, below=of depart, align=center] (analyse){Analyse\\ temps, cas,\\ type de phrase};
\node[box, below=of analyse, align=center] (construction){Construction\\ V2 + cadre verbal\\ ou subordonnée};
\node[box, below=of construction, align=center] (verif){Vérification\\ verbe, cas,\\ majuscules};
\node[box, fill=popgreenL, draw=popgreen, below=of verif, align=center] (fin){Phrase allemande\\ correcte};
\draw[fleche] (depart) -- (analyse);
\draw[fleche] (analyse) -- (construction);
\draw[fleche] (construction) -- (verif);
\draw[fleche] (verif) -- (fin);
\draw[fleche, dashed, popmuted] (verif.west) to[out=180,in=180] node[left, font=\footnotesize, align=center]{erreur ?\\ on recommence} (analyse.west);
\end{tikzpicture}
\legende{Organigramme de la méthode du thème : de la phrase française à la phrase allemande.}
\end{popfigure}

\textbf{Thème guidé : cinq phrases progressives}\par

Appliquons la méthode à cinq phrases de difficulté croissante, toutes sur le thème de la citoyenneté.

\begin{exemplebox}[Phrase 1 — présent simple]
\textbf{Français :} « Les jeunes doivent participer à la vie de leur quartier. »

\textbf{Analyse :} verbe « devoir » + infinitif « participer » → verbe modal \all{müssen} + infinitif \all{teilnehmen} ; présent ; phrase principale ; \all{an etwas teilnehmen} exige le datif.

\textbf{Allemand :} \all{Die Jugendlichen müssen am Leben ihres Viertels teilnehmen.}

Notez le cadre : \all{müssen} en position 2, \all{teilnehmen} à la toute fin, et \all{am Leben} (= \all{an dem Leben}) au datif.
\end{exemplebox}

\begin{exemplebox}[Phrase 2 — subordonnée]
\textbf{Français :} « Je pense que le bénévolat est important pour la société. »

\textbf{Analyse :} « que » introduit une subordonnée → \all{dass} + verbe conjugué \textbf{à la fin}.

\textbf{Allemand :} \all{Ich denke, dass das Ehrenamt wichtig für die Gesellschaft ist.}

Piège évité : le verbe \all{ist} ne reste pas après le sujet comme en français ; il part en fin de subordonnée.
\end{exemplebox}

\begin{exemplebox}[Phrase 3 — passé composé]
\textbf{Français :} « Nous avons interviewé le maire de notre ville. »

\textbf{Analyse :} passé composé → Perfekt ; verbe \all{interviewen} (faible, participe \all{interviewt}, sans \all{ge-} car le verbe se termine en \all{-ieren}) ; « le maire » est complément direct → accusatif \all{den Bürgermeister}.

\textbf{Allemand :} \all{Wir haben den Bürgermeister unserer Stadt interviewt.}
\end{exemplebox}

\begin{exemplebox}[Phrase 4 — proposition relative]
\textbf{Français :} « Le jeune homme qui travaille comme volontaire habite à Bonn. »

\textbf{Analyse :} « qui » reprend « le jeune homme », sujet du verbe de la relative → \all{der} (nominatif masculin) ; verbe de la relative \textbf{à la fin} ; la relative est encadrée par des virgules.

\textbf{Allemand :} \all{Der junge Mann, der als Freiwilliger arbeitet, wohnt in Bonn.}
\end{exemplebox}

\begin{exemplebox}[Phrase 5 — futur]
\textbf{Français :} « Demain, les élèves voteront pour leurs délégués. »

\textbf{Analyse :} futur → Futur I : \all{werden} + infinitif ; « demain » en tête → inversion sujet/verbe ; « élire, voter pour quelqu'un » = \all{wählen} + accusatif.

\textbf{Allemand :} \all{Morgen werden die Schüler ihre Klassensprecher wählen.}

Notez l'inversion : \all{Morgen werden die Schüler ...} et non \all{Morgen die Schüler werden ...}.
\end{exemplebox}

\begin{attention}[Au Perfekt, le participe passe en fin de phrase]
L'erreur n° 1 des francophones est de coller le participe à l'auxiliaire, comme en français (« nous avons interviewé »). En allemand, le participe ferme le cadre, quelle que soit la longueur de la phrase :

\all{Wir haben den Bürgermeister interviewt.} « Nous avons interviewé le maire. »

\all{Wir haben gestern in der Schule den Bürgermeister unserer Stadt interviewt.} « Nous avons interviewé hier à l'école le maire de notre ville. »

Jamais \all{Wir haben interviewt den Bürgermeister}. Même règle pour les verbes à particule : \all{Er nimmt am Projekt teil.} « Il participe au projet. »
\end{attention}

\courssub{La version : traduire de l'allemand vers le français}

\textbf{Observer avant de traduire}\par

Lisez cette phrase :

\all{Wer die Demokratie stärken will, muss sich engagieren.}

Un mot à mot donnerait : « Qui la démocratie renforcer veut, doit s'engager » — incompréhensible. La bonne traduction est : « Qui veut renforcer la démocratie doit s'engager. » En version, on \textbf{démonte} la machine allemande, puis on \textbf{reconstruit} un français naturel.

\begin{methode}[Les 4 étapes de la version]
\begin{enumerate}
\item \textbf{Repérer le verbe conjugué} de chaque proposition : il est en position 2 dans la principale, à la fin dans les subordonnées.
\item \textbf{Découper la phrase} en propositions : principale, relatives (\all{der, die, das, wer ...}), subordonnées (\all{dass, weil, wenn, obwohl ...}).
\item \textbf{Traduire chaque bloc} en respectant les cas : le nominatif donne le sujet, l'accusatif le complément direct, le datif le complément indirect.
\item \textbf{Rédiger en français naturel} : ordre sujet–verbe–complément, participes collés à l'auxiliaire, tournures idiomatiques. Relire à voix haute : si ce n'est pas du français correct, recommencer.
\end{enumerate}
\end{methode}

\begin{attention}[« man » se traduit par « on »]
Le pronom indéfini \all{man} (toujours à la 3\textsuperscript{e} personne du singulier) correspond au français « on » :

\all{Man muss die Gesetze respektieren.} « On doit respecter les lois. »

\all{In einer Demokratie darf man seine Meinung sagen.} « Dans une démocratie, on peut dire son opinion. »

Ne traduisez jamais \all{man} par « l'homme » (c'est \all{der Mann}) et ne le confondez pas avec \all{Mann}.
\end{attention}

\textbf{Version guidée : un texte sur un jeune volontaire à Bonn}\par

\begin{textebox}[Ein Freiwilliger in Bonn]
Lukas ist neunzehn Jahre alt und wohnt in Bonn. Nach dem Abitur hat er ein freiwilliges soziales Jahr gemacht. Jeden Morgen ist er um acht Uhr aufgestanden und hat in einem Jugendzentrum gearbeitet. Dort hat er Kindern bei den Hausaufgaben geholfen und Sport mit ihnen gemacht. „Man lernt hier, Verantwortung zu tragen“, sagt Lukas. „Ich habe auch gelernt, dass Demokratie mit kleinen Taten anfängt.“ Am Ende des Jahres hat der Bürgermeister von Bonn alle Freiwilligen zu einem Fest eingeladen. Lukas sagt, dass er dieses Jahr nie vergessen wird. Jetzt studiert er Politik, und später will er sich für seine Stadt engagieren.
\end{textebox}

\textbf{Glossaire :} das Abitur (baccalauréat) ; das freiwillige soziale Jahr (année de service volontaire) ; aufstehen (se lever, verbe à particule séparable) ; das Jugendzentrum, die Jugendzentren (le centre de jeunes) ; helfen + datif (aider) ; die Verantwortung tragen (assumer la responsabilité) ; die Tat, die Taten (l'acte, l'action) ; einladen (inviter) ; vergessen (oublier) ; sich engagieren (s'engager).

\textbf{Traduction de référence :} « Lukas a dix-neuf ans et habite à Bonn. Après le bac, il a fait une année de service volontaire. Chaque matin, il se levait à huit heures et travaillait dans un centre de jeunes. Là, il aidait les enfants à faire leurs devoirs et faisait du sport avec eux. „Ici, on apprend à assumer des responsabilités“, dit Lukas. „J'ai aussi appris que la démocratie commence par de petites actions.“ À la fin de l'année, le maire de Bonn a invité tous les volontaires à une fête. Lukas dit qu'il n'oubliera jamais cette année. Maintenant, il étudie les sciences politiques et, plus tard, il veut s'engager pour sa ville. »

Remarquez le choix des temps français : le Präteritum allemand (\all{ist ... aufgestanden}, \all{hat ... gearbeitet} au Perfekt dans le texte) se rend ici par l'imparfait pour les habitudes (« il se levait », « il aidait ») et par le passé composé pour les actions ponctuelles (« il a fait », « a invité »). La version exige donc aussi une bonne maîtrise des temps \textbf{français}.

\textbf{Questions de traduction ciblées :}
\begin{enumerate}
\item Traduisez : \all{Nach dem Abitur hat er ein freiwilliges soziales Jahr gemacht.} — Attention à \all{nach} + datif et au participe final.
\item Traduisez : \all{Dort hat er Kindern bei den Hausaufgaben geholfen.} — Pourquoi \all{Kindern} est-il au datif ?
\item Traduisez : \all{Lukas sagt, dass er dieses Jahr nie vergessen wird.} — Où se place le verbe dans la subordonnée ?
\end{enumerate}

\begin{corrige}[Questions de traduction]
\begin{enumerate}
\item « Après le bac, il a fait une année de service volontaire. » \all{Nach} exige le datif (\all{nach dem Abitur}) ; le participe \all{gemacht} ferme le cadre.
\item « Là, il a aidé les enfants à faire leurs devoirs. » Le verbe \all{helfen} exige toujours le \textbf{datif} : \all{er hilft den Kindern}, contrairement au français « aider quelqu'un ».
\item « Lukas dit qu'il n'oubliera jamais cette année. » Dans la subordonnée introduite par \all{dass}, l'auxiliaire \all{wird} passe à la fin ; en français, on reconstruit : « qu'il n'oubliera jamais ».
\end{enumerate}
\end{corrige}

\courssub{Les pièges fréquents et la grille d'auto-correction}

\begin{center}
\begin{tabularx}{\linewidth}{|X|X|X|}
\hline
\rowcolor{popblueL} \textbf{Piège} & \textbf{Faux} & \textbf{Correct} \\
\hline
Majuscule oubliée & \all{die jugendlichen wählen} & \all{Die Jugendlichen wählen.} \\
\hline
Verbe mal placé & \all{Morgen die Schüler wählen} & \all{Morgen wählen die Schüler.} \\
\hline
Participe collé & \all{Wir haben interviewt den Bürgermeister} & \all{Wir haben den Bürgermeister interviewt.} \\
\hline
Cas faux après préposition & \all{mit den Bus} & \all{mit dem Bus} (datif) \\
\hline
Calque du français & \all{Ich habe recht} → « j'ai droit » & \all{Ich habe recht.} = « j'ai raison » \\
\hline
\all{man} mal traduit & « l'homme doit ... » & « on doit ... » \\
\hline
\end{tabularx}
\end{center}

\begin{propriete}[Grille d'auto-correction en cinq points]
Après chaque traduction, relisez en cochant mentalement :
\begin{itemize}
\item \textbf{Majuscules} : chaque nom allemand commence-t-il par une majuscule ?
\item \textbf{Place du verbe} : verbe conjugué en position 2 (principale) ou en fin (subordonnée) ?
\item \textbf{Cas} : nominatif pour le sujet, accusatif après \all{durch, für, ohne}, datif après \all{mit, nach, zu, bei} et après \all{helfen, danken} ?
\item \textbf{Temps} : le temps allemand correspond-il au temps français (présent, Perfekt, Futur) ?
\item \textbf{Sens global} : la phrase traduite dit-elle exactement la même chose, sans ajout ni oubli ?
\end{itemize}
\end{propriete}

\begin{exoresolu}[Thème : phrase complète à traduire]
\textbf{Énoncé :} Traduisez en allemand : « Hier, les citoyens ont discuté avec le maire parce qu'ils veulent améliorer leur quartier. »

\tcblower
\textbf{Solution détaillée :}
\begin{enumerate}
\item \textbf{Lecture} : deux propositions, une principale et une causale avec « parce que ».
\item \textbf{Verbes et temps} : « ont discuté » → Perfekt de \all{diskutieren} ; « veulent » → présent de \all{wollen} + infinitif \all{verbessern}.
\item \textbf{Construction} : « hier » en tête → inversion ; « parce que » → \all{weil} + verbe à la fin ; « avec » → \all{mit} + datif : \all{mit dem Bürgermeister}.
\item \textbf{Traduction} : \all{Gestern haben die Bürger mit dem Bürgermeister diskutiert, weil sie ihr Viertel verbessern wollen.}
\item \textbf{Vérification} : majuscules (\all{Bürger, Bürgermeister, Viertel}) ; \all{haben} en position 2, \all{diskutiert} en fin ; datif \all{dem Bürgermeister} ; \all{wollen} en fin de subordonnée ; sens fidèle.
\end{enumerate}
\end{exoresolu}

\begin{exercice}[À vous de traduire]
\textbf{Thème :} 1. « On doit respecter les lois de son pays. » 2. « La semaine dernière, nous avons aidé une association de notre ville. » 3. « Je crois que chaque citoyen doit voter. »

\textbf{Version :} 4. \all{Wer sich nicht engagiert, darf sich nicht beschweren.} 5. \all{Nach dem Fest haben die Freiwilligen dem Bürgermeister gedankt.}
\end{exercice}

\begin{corrige}[Exercice « À vous de traduire »]
1. \all{Man muss die Gesetze seines Landes respektieren.} (cadre \all{muss ... respektieren} ; génitif \all{seines Landes}).

2. \all{Letzte Woche haben wir einem Verein unserer Stadt geholfen.} (inversion après \all{Letzte Woche} ; \all{helfen} + datif : \all{einem Verein} ; participe \all{geholfen} en fin).

3. \all{Ich glaube, dass jeder Bürger wählen muss.} (verbe \all{muss} en fin de subordonnée).

4. « Qui ne s'engage pas n'a pas le droit de se plaindre. » (\all{sich beschweren} = se plaindre ; ne pas calquer « qui ne s'engage pas, ne doit pas se plaindre » sans reformuler).

5. « Après la fête, les volontaires ont remercié le maire. » (\all{danken} + datif : \all{dem Bürgermeister} ; participe \all{gedankt} en fin).
\end{corrige}

\begin{savaistu}[Comment la traduction est-elle notée au Bac ?]
Au baccalauréat, la traduction est notée sur \textbf{deux critères combinés} : la \textbf{fidélité du sens} (avez-vous tout dit, sans déformation ?) et la \textbf{correction grammaticale} (place du verbe, cas, temps, orthographe). Une phrase au sens juste mais mal construite — par exemple un participe collé à l'auxiliaire ou un datif oublié — perd des points, et une phrase grammaticalement parfaite mais infidèle au sens aussi. D'où l'importance de la grille d'auto-correction : deux minutes de relecture peuvent sauver plusieurs points.
\end{savaistu}

\begin{aretenir}[L'essentiel de la section]
\begin{itemize}
\item \textbf{Thème} : lire, analyser (temps, cas, construction), construire le cadre verbal, traduire, relire. Jamais de mot à mot.
\item \textbf{Version} : repérer les verbes conjugués, découper les propositions, traduire chaque bloc, rédiger un français naturel.
\item Le verbe conjugué est en position 2 dans la principale, à la fin dans la subordonnée ; le participe et l'infinitif ferment toujours le cadre.
\item \all{man} = « on » ; \all{helfen} et \all{danken} exigent le datif ; majuscule à tous les noms.
\item La relecture systématique (majuscules, verbe, cas, temps, sens) est la clé de la réussite.
\end{itemize}
\end{aretenir}

\coursec{Commentaire modèle et production guidée}

Dans la section précédente, nous avons traduit des phrases et de courts textes sur la citoyenneté, du français vers l'allemand (thème) et de l'allemand vers le français (version). Traduire, c'est déjà produire : on choisit des mots, on construit des phrases, on respecte l'ordre des mots allemand. Nous allons maintenant franchir une étape supplémentaire : rédiger nos propres textes. Dans cette dernière section, tu vas d'abord observer un \cle{commentaire modèle} du texte support \all{Jugendliche engagieren sich}, relever les phrases témoins réutilisables, puis produire toi-même un mini-discours et des questions d'interview, avant de corriger la production d'un camarade à l'aide d'une grille de relecture.

\courssub{Le commentaire modèle : observer avant d'imiter}

Un commentaire de texte en allemand, au niveau Première, suit toujours la même architecture en trois temps : une \cle{introduction} (présenter le texte), un \cle{développement} en deux paragraphes (analyser le contenu avec des citations traduites), une \cle{conclusion} (donner son avis personnel justifié). Lis d'abord le début du commentaire modèle ci-dessous : l'introduction et les deux paragraphes du développement.

\begin{textebox}[Commentaire modèle : introduction et développement]
Der Text „Jugendliche engagieren sich“ handelt vom Engagement junger Menschen in Deutschland und in Kamerun. Der Autor zeigt, dass Jugendliche nicht nur Rechte, sondern auch Pflichten haben.

Zuerst beschreibt der Text die Rechte der jungen Bürger: Sie dürfen ihre Meinung sagen, sie können wählen gehen und sie haben das Recht auf Bildung. Der Autor schreibt: „Junge Menschen sind die Zukunft unseres Landes.“ Das bedeutet: Die Jugendlichen von heute sind die Bürger von morgen.

Dann erklärt der Text die Pflichten. Ein guter Bürger respektiert die Gesetze, hilft den anderen und engagiert sich in der Schule oder im Quartier. In Kamerun zum Beispiel putzen viele Jugendliche am Wochenende ihr Viertel. In Deutschland arbeiten viele freiwillig im Verein oder bei der Feuerwehr.
\end{textebox}

\textbf{Glossaire :} \all{das Engagement, -s} (l'engagement), \all{das Recht, -e} (le droit), \all{die Pflicht, -en} (le devoir), \all{der Bürger, -} (le citoyen), \all{wählen gehen} (aller voter), \all{die Bildung} (l'éducation, la formation), \all{das Gesetz, -e} (la loi), \all{sich engagieren} (s'engager), \all{das Quartier, -e} (le quartier), \all{putzen} (nettoyer), \all{das Viertel, -} (le quartier), \all{freiwillig} (bénévolement), \all{der Verein, -e} (le club, l'association), \all{die Feuerwehr, -en} (les pompiers).

\begin{definition}[Le commentaire de texte]
Un \cle{commentaire de texte} est une production écrite qui présente un texte, en analyse les idées principales à l'aide de citations, et se termine par une opinion personnelle argumentée. En allemand, il se construit avec des \cle{phrases témoins} fixes que l'on réutilise d'un texte à l'autre : seule la matière change, le squelette reste identique.
\end{definition}

\courssub{Les phrases témoins à réutiliser}

Observons maintenant les phrases-clés du modèle. Elles sont \emph{transposables} : tu peux les réemployer dans n'importe quel commentaire, en changeant seulement le contenu.

\begin{center}
\begin{tabularx}{\linewidth}{|X|X|X|}
\hline
\rowcolor{popblueL} Fonction & Phrase témoin (Deutsch) & Traduction \\
\hline
Présenter le sujet & \all{Der Text handelt von...} & Le texte traite de... \\
\hline
Présenter la thèse & \all{Der Autor zeigt, dass...} & L'auteur montre que... \\
\hline
Annoncer le plan & \all{Zuerst..., dann..., schließlich...} & D'abord..., ensuite..., enfin... \\
\hline
Citer le texte & \all{Der Autor schreibt: „...“} & L'auteur écrit : « ... » \\
\hline
Expliquer une citation & \all{Das bedeutet: ...} & Cela signifie : ... \\
\hline
Donner un exemple & \all{zum Beispiel} & par exemple \\
\hline
Donner son avis & \all{Meiner Meinung nach ist Engagement wichtig, weil...} & À mon avis, l'engagement est important, parce que... \\
\hline
Conclure & \all{Zusammenfassend kann man sagen, dass...} & En résumé, on peut dire que... \\
\hline
\end{tabularx}
\end{center}

\begin{propriete}[Règle de construction des phrases témoins]
\begin{itemize}
\item Après \all{handeln von} + \textbf{datif} : \all{Der Text handelt vom Engagement.} (von + dem = vom).
\item Après \all{dass} et \all{weil}, le verbe conjugué va \textbf{à la fin} de la subordonnée : \all{Der Autor zeigt, dass Jugendliche Pflichten \textbf{haben}.} / \all{...weil die Stimme ein Recht \textbf{ist}.}
\item \all{Meiner Meinung nach} est suivi du verbe en \textbf{deuxième position} : \all{Meiner Meinung nach \textbf{ist} Engagement wichtig.}
\item Avec un modal (\all{können}, \all{müssen}, \all{sollen}, \all{dürfen}) dans une subordonnée, l'ordre à la fin est : infinitif puis modal conjugué : \all{..., dass jeder Bürger wählen gehen \textbf{kann}.}
\end{itemize}
\end{propriete}

\begin{exemplebox}[Exemples traduits]
\all{Meiner Meinung nach sollte jeder Bürger wählen gehen, denn die Stimme ist ein Recht und eine Pflicht.} « À mon avis, chaque citoyen devrait aller voter, car le vote est un droit et un devoir. »

\all{Der Text handelt von Jugendlichen, die sich in der Schule engagieren.} « Le texte traite de jeunes qui s'engagent à l'école. »

\all{Das bedeutet: Jeder kann etwas für die Gemeinschaft tun.} « Cela signifie : chacun peut faire quelque chose pour la communauté. »
\end{exemplebox}

\begin{methode}[Rédiger la conclusion du commentaire]
Pour la conclusion, donne toujours ton avis avec une justification, en trois étapes :
\begin{enumerate}
\item Annonce ton opinion : \all{Ich finde Engagement wichtig, weil...} ou \all{Meiner Meinung nach...}
\item Justifie avec \all{weil} (verbe à la fin) ou \all{denn} (verbe en deuxième position) : \all{...weil junge Menschen die Zukunft sind.} / \all{...denn junge Menschen sind die Zukunft.}
\item Termine par une phrase de synthèse : \all{Zusammenfassend kann man sagen, dass Rechte und Pflichten zusammengehören.}
\end{enumerate}
\end{methode}

\begin{attention}[Varier les débuts de phrases]
Dans une production, les francophones commencent presque toutes leurs phrases par \all{Ich}. C'est lourd et répétitif. Varie les débuts : \all{Außerdem finde ich...}, \all{Deshalb engagiere ich mich...}, \all{Zuerst helfe ich...}, \all{In meinem Quartier gibt es...}. Rappel : après un élément placé en tête, le verbe reste en \textbf{deuxième position} et le sujet suit le verbe (inversion). Ne dis jamais \all{Deshalb ich engagiere mich...} : c'est un calque du français « c'est pourquoi je... ».
\end{attention}

\courssub{Production guidée n°1 : le mini-discours du candidat délégué}

Tu es candidat(e) au poste de délégué(e) de classe. Rédige un mini-discours de 8 à 10 lignes en respectant le plan imposé ci-dessous.

\begin{methode}[Plan imposé du discours]
\begin{enumerate}
\item Salutation et présentation : \all{Liebe Mitschüler, ich heiße... und ich möchte Klassensprecher werden.}
\item Deux arguments (avec \all{zuerst} / \all{außerdem}) : ce que tu veux améliorer dans la classe.
\item Une promesse concrète : \all{Ich verspreche, dass ich...} (attention : verbe à la fin de la subordonnée !)
\item Appel au vote et remerciement : \all{Wählt mich! Vielen Dank für eure Aufmerksamkeit.}
\end{enumerate}
\end{methode}

\begin{exoresolu}[Discours modèle corrigé]
Énoncé : rédige le discours complet en suivant le plan imposé.
\tcblower
Solution détaillée :

\all{Liebe Mitschüler, liebe Mitschülerinnen, ich heiße Amina und ich möchte Klassensprecherin werden. Zuerst finde ich, dass unsere Klasse sauberer sein muss. Deshalb schlage ich vor, dass wir jeden Freitag zusammen putzen. Außerdem möchte ich einen Bücherschrank für die Klasse organisieren, denn Lesen ist wichtig für unsere Bildung. Ich verspreche, dass ich immer eure Meinung höre und mit den Lehrern spreche. Meiner Meinung nach ist ein Klassensprecher ein Bürger in der Schule: Er hat Rechte, aber auch Pflichten. Wählt mich! Vielen Dank für eure Aufmerksamkeit.}

Points de grammaire vérifiés : verbe en fin de subordonnée après \all{dass} (\all{sein muss}, \all{putzen}, \all{höre}, \all{spreche}), inversion après \all{deshalb} et \all{außerdem} (\all{schlage ich}, \all{möchte ich}), majuscules aux noms (\all{Bücherschrank}, \all{Bildung}, \all{Bürger}, \all{Meinung}), impératif pluriel \all{Wählt mich!}, datif pluriel \all{mit den Lehrern}.
\end{exoresolu}

\courssub{Production guidée n°2 : cinq questions d'interview}

Ton école reçoit un correspondant allemand de Bonn. Prépare \textbf{cinq questions} sur l'engagement des jeunes. Rappel de la règle : dans une question fermée (réponse oui/non), le verbe conjugué est en \textbf{première position} ; dans une question ouverte, le mot interrogatif (\all{Was}, \all{Warum}, \all{Wie}, \all{Wann}, \all{Wo}) est en première position et le verbe conjugué en deuxième. Avec un adulte ou un invité, on vouvoie : \all{Sie} (toujours avec majuscule).

\begin{exemplebox}[Modèles de questions]
\all{Engagieren sich die Jugendlichen in Deutschland?} « Les jeunes s'engagent-ils en Allemagne ? » (question fermée, verbe en tête)

\all{Warum engagieren Sie sich?} « Pourquoi vous engagez-vous ? » (question ouverte)

\all{Was kann ein junger Mensch für die Umwelt tun?} « Que peut faire un jeune pour l'environnement ? »

\all{Wie wichtig ist das Wählen für junge Bürger?} « Quelle importance a le vote pour les jeunes citoyens ? »

\all{Gibt es in Ihrer Stadt viele Vereine für Jugendliche?} « Y a-t-il dans votre ville beaucoup d'associations pour les jeunes ? »
\end{exemplebox}

\begin{exercice}[À toi de jouer]
Rédige tes cinq questions : deux questions fermées (verbe en tête) et trois questions ouvertes avec \all{Was}, \all{Warum}, \all{Wie}. Utilise le vocabulaire du chapitre : \all{das Engagement}, \all{der Verein}, \all{wählen}, \all{die Pflicht}, \all{freiwillig helfen}.
\end{exercice}

\begin{corrige}[Exercice : cinq questions d'interview]
Proposition de corrigé (d'autres formulations sont possibles) :

1. \all{Engagieren sich viele Jugendliche in Ihrer Stadt?} 2. \all{Gehen die jungen Deutschen wählen?} 3. \all{Was machen die Jugendlichen in den Vereinen?} 4. \all{Warum ist freiwillige Arbeit wichtig für Sie?} 5. \all{Wie kann die Schule das Engagement fördern?}

Vérification : questions 1 et 2 = verbe en première position ; questions 3, 4, 5 = mot interrogatif + verbe en deuxième position ; majuscules aux noms (\all{Jugendliche}, \all{Deutschen}, \all{Vereine}, \all{Arbeit}, \all{Schule}, \all{Engagement}) ; vouvoiement \all{Ihrer}, \all{Sie} avec majuscule.
\end{corrige}

\courssub{Relecture croisée : corriger avec la grille}

Échange ta production (discours ou questions) avec un camarade. Corrige son texte avec la grille suivante : coche chaque critère et note une erreur trouvée pour chaque ligne.

\begin{center}
\begin{tabularx}{\linewidth}{|X|X|X|}
\hline
\rowcolor{popgreenL} Critère & Question de contrôle & Exemple d'erreur typique \\
\hline
Majuscules & Tous les noms ont-ils une majuscule ? & \all{der bürger} $\rightarrow$ \all{der Bürger} \\
\hline
Place du verbe & Verbe en 2e position ? À la fin après \all{dass}/\all{weil} ? & \all{weil ich habe Rechte} $\rightarrow$ \all{weil ich Rechte habe} \\
\hline
Cas & Datif après \all{mit}, \all{von} ; accusatif après \all{für} ? & \all{mit die Klasse} $\rightarrow$ \all{mit der Klasse} \\
\hline
Connecteurs & Y a-t-il \all{zuerst}, \all{außerdem}, \all{deshalb} ? & texte sans connecteurs = monotone \\
\hline
Orthographe & Umlaute (\all{ä}, \all{ö}, \all{ü}) et \all{ß} corrects ? & \all{mochte} $\rightarrow$ \all{möchte} \\
\hline
\end{tabularx}
\end{center}

\begin{attention}[Pièges de la relecture]
Ne corrige pas seulement l'orthographe ! Les trois erreurs les plus graves des francophones sont : le verbe mal placé dans la subordonnée, l'oubli de la majuscule aux noms, et le mauvais cas après une préposition. Relis ton texte \textbf{trois fois}, une fois par critère : d'abord les verbes, ensuite les majuscules, enfin les cas. Cette méthode de relecture en trois passages est aussi valable pour l'épreuve de thème au Baccalauréat.
\end{attention}

\courssub{Bilan visuel et auto-évaluation}

La carte mentale ci-dessous résume tout le chapitre : la citoyenneté mobilise quatre compétences que tu as travaillées.

\begin{popfigure}
\begin{tikzpicture}[mindmap, grow cyclic, every node/.style={concept, align=center}, root concept/.append style={concept color=popblue, text=white, font=\bfseries}, level 1/.append style={level distance=3.2cm, sibling angle=90}]
\node[root concept]{Citoyenneté}
  child[concept color=poporange, text=white] { node[align=center]{Lexique\\ \all{der Bürger}\\ \all{das Recht}\\ \all{die Pflicht}} }
  child[concept color=popgreen, text=white] { node[align=center]{Grammaire\\ verbe en 2e\\ position\\ \all{dass}/\all{weil} + fin} }
  child[concept color=poppurple, text=white] { node[align=center]{Production\\ interview\\ lettre\\ discours} }
  child[concept color=popgold, text=white] { node[align=center]{Traduction\\ thème\\ version\\ relecture} };
\end{tikzpicture}
\legende{Carte mentale récapitulative du chapitre « Citoyenneté »}
\end{popfigure}

\begin{aretenir}[Ce que je sais faire maintenant]
Auto-évaluation : coche mentalement chaque compétence acquise.
\begin{itemize}
\item Je sais présenter un texte : \all{Der Text handelt von...}
\item Je sais analyser avec une citation : \all{Der Autor zeigt, dass...}
\item Je sais donner mon avis justifié : \all{Meiner Meinung nach..., weil...}
\item Je sais rédiger un discours avec un plan et des connecteurs.
\item Je sais poser des questions fermées (verbe en tête) et ouvertes (mot interrogatif + verbe).
\item Je sais relire et corriger : majuscules, place du verbe, cas, connecteurs.
\end{itemize}
\end{aretenir}

\begin{savaistu}[D'où vient le mot \all{Bürger} ?]
Le mot \all{der Bürger} (le citoyen) vient de \all{die Burg} (le château fort). Au Moyen Âge, le \all{Bürger} était celui qui habitait dans la ville protégée par le château et qui participait à sa défense. Le citoyen avait donc déjà des droits (la protection de la ville) et des devoirs (la défendre) : l'idée de citoyenneté est née avec ce lien entre droits et devoirs. La famille de mots est encore visible aujourd'hui : \all{die Burg}, \all{der Bürger}, \all{das Bürgeramt} (le service administratif municipal), \all{bürgerlich} (civil, citoyen).
\end{savaistu}

\newpage

\coursec{Activité d'intégration}

\courssub{Objectifs de l'activité}

À la fin de cette activité d'intégration, l'élève sera capable de :

\begin{itemize}
\item rédiger et prononcer un \cle{mini-discours électoral} de 8 à 10 lignes en allemand, avec connecteurs logiques et appel à l'action ;
\item préparer et mener un \cle{interview} de cinq questions (vouvoiement, \all{W-Fragen}, relances) ;
\item \cle{traduire} de courtes réponses de l'allemand vers le français pour un journal bilingue ;
\item réinvestir le \cle{lexique de la citoyenneté} (vote, droits, devoirs, engagement) et la grammaire révisée (place du verbe, cas, temps, verbes à particule) ;
\item évaluer la production des camarades à l'aide d'une grille simple et argumenter son choix.
\end{itemize}

\courssub{Situation}

Le lycée de Yaoundé organise l'élection des délégués de classe. Un correspondant d'un journal scolaire allemand, \all{Lukas Berger}, est en visite d'échange au Cameroun. Il veut écrire un article sur la vie démocratique dans les lycées camerounais. Votre classe se transforme en salle de presse : chaque groupe de trois élèves joue le rôle d'un candidat, d'un journaliste et d'un traducteur. Voici le message que Lukas a affiché au tableau :

\begin{textebox}[Message de Lukas, correspondant allemand]
Liebe Schülerinnen und Schüler,

ich heiße Lukas und ich komme aus München. Ich schreibe für unsere Schülerzeitung „Die Brücke“. Nächste Woche wählt ihr eure Klassensprecher. Das ist spannend! Ich möchte gern wissen: Was versprechen die Kandidaten? Was denken die Wähler? Wer hat die besten Ideen für die Schule?

Bitte helft mir: Ein Kandidat hält eine kleine Rede, ein Journalist stellt fünf Fragen, und ein Übersetzer übersetzt die besten Antworten ins Französische. Am Ende wählen wir das beste Team. Viel Glück!

Euer Lukas
\end{textebox}

\textbf{Glossaire :} \all{die Schülerzeitung, -en} (le journal scolaire) ; \all{der Klassensprecher, -} (le délégué de classe) ; \all{wählen} (élire, voter) ; \all{versprechen} (promettre) ; \all{der Wähler, -} (l'électeur) ; \all{die Rede, -n} (le discours) ; \all{die Antwort, -en} (la réponse) ; \all{übersetzen} (traduire).

\courssub{Consignes}

Travaillez par groupes de trois : un \cle{candidat}, un \cle{journaliste}, un \cle{traducteur}. Suivez les étapes dans l'ordre.

\begin{enumerate}
\item \textbf{Préparation du discours (candidat, 15 min).} Rédigez un mini-discours électoral de 8 à 10 lignes. Structure obligatoire : salutation, présentation, deux promesses concrètes pour la classe, appel à l'action. Utilisez au moins deux connecteurs (\all{zuerst}, \all{außerdem}, \all{deshalb}, \all{schließlich}) et une phrase avec \all{weil} ou \all{wenn}.

\item \textbf{Préparation de l'interview (journaliste, 15 min).} Écrivez cinq questions : au moins trois \all{W-Fragen} (avec \all{was}, \all{warum}, \all{wie}, \all{wann}) et une question fermée avec inversion. Vouvoyez le candidat (\all{Sie}). Prévoyez une relance : \all{Können Sie das genauer erklären?} (Pouvez-vous expliquer cela plus précisément ?).

\item \textbf{Jeu de rôle (10 min par groupe).} Le candidat prononce son discours devant la classe, puis le journaliste mène l'interview. Le traducteur prend des notes en allemand.

\item \textbf{Traduction (traducteur, 10 min).} Choisissez les trois meilleures réponses de l'interview et traduisez-les en français correct pour le journal bilingue du lycée. Attention à l'ordre des mots et aux cas en allemand, et au naturel en français.

\item \textbf{Évaluation croisée.} Pendant chaque prestation, les autres élèves remplissent la grille ci-dessous (de 0 à 2 points par critère).

\item \textbf{Vote final.} La classe vote à main levée pour le meilleur duo interview/discours. Justifiez votre vote en une phrase allemande : \all{Ich wähle Gruppe 2, weil ...}
\end{enumerate}

\begin{center}
\begin{tabularx}{\linewidth}{|X|c|}
\hline
\rowcolor{popblueL} \textbf{Critère} & \textbf{Points (0--2)} \\
\hline
Lexique de la citoyenneté (au moins 5 mots du chapitre) & \\
\hline
Place du verbe (position 2, verbe final après \all{weil}) & \\
\hline
Cas corrects (Akkusativ / Dativ) & \\
\hline
Questions bien formées (inversion, \all{W-Fragen}) & \\
\hline
Fluidité et prononciation & \\
\hline
Qualité de la traduction française & \\
\hline
\end{tabularx}
\end{center}

\courssub{Ressources à mobiliser}

\begin{aretenir}[Boîte à outils de l'activité]
\textbf{Lexique de la citoyenneté :} \all{die Wahl, -en} (l'élection) ; \all{wählen} (élire) ; \all{der Kandidat, -en} (le candidat) ; \all{das Recht, -e} (le droit) ; \all{die Pflicht, -en} (le devoir) ; \all{sich engagieren} (s'engager) ; \all{die Meinung, -en} (l'opinion) ; \all{versprechen} (promettre) ; \all{gemeinsam} (ensemble) ; \all{die Verantwortung} (la responsabilité).

\textbf{Structures du discours :} \all{Liebe Mitschüler!} (Chers camarades !) ; \all{Ich möchte eure Stimme, weil ...} (Je veux votre voix parce que...) ; \all{Wenn ihr mich wählt, ...} (Si vous m'élisez, ...) ; \all{Wählt mich!} (Élisez-moi !).

\textbf{Structures de l'interview :} \all{Warum wollen Sie Klassensprecher werden?} ; \all{Was wollen Sie für die Klasse tun?} ; \all{Wie wollen Sie das machen?} ; \all{Haben Sie Erfahrung?}

\textbf{Grammaire :} verbe en position 2 ; verbe conjugué à la fin après \all{weil}, \all{dass}, \all{wenn} ; verbes à particule séparable (\all{sich engagieren} : \all{Ich engagiere mich für die Klasse}) ; Akkusativ après \all{für}, Dativ après \all{mit}, \all{helfen}.
\end{aretenir}

\begin{attention}[Pièges à éviter]
\begin{itemize}
\item Ne dites pas \all{Ich will gewählt werden weil ich bin gut} : après \all{weil}, le verbe va à la fin : \all{..., weil ich gut bin}.
\item Le verbe \all{helfen} exige le Datif : \all{Ich helfe den Schülern} (et non \all{die Schüler}).
\item \all{wählen} = élire/voter ; \all{treffen} = rencontrer. Ne confondez pas avec le français « choisir une rencontre ».
\item Dans le discours, on peut tutoyer les camarades (\all{ihr}), mais dans l'interview formelle, le journaliste vouvoie (\all{Sie}).
\end{itemize}
\end{attention}

\courssub{Proposition de production et corrigé commenté}

\begin{exoresolu}[Modèle complet : discours, interview et traduction]
Produisez le discours du candidat, les cinq questions du journaliste et la traduction française de trois réponses.

\tcblower

\textbf{1. Mini-discours électoral (candidat)}

\begin{textebox}[Die Rede von Amina]
Liebe Mitschülerinnen und Mitschüler,

ich heiße Amina und ich möchte eure Klassensprecherin werden. Zuerst verspreche ich: Ich höre euch zu und ich bringe eure Probleme zur Direktion. Außerdem organisiere ich jede Woche eine Deutschstunde für alle, die Hilfe brauchen. Wenn ihr mich wählt, arbeite ich jeden Tag für unsere Klasse, denn gemeinsam sind wir stark. Schließlich habe ich schon Erfahrung: Letztes Jahr habe ich den Sporttag organisiert. Wählt mich, denn eure Stimme ist wichtig! Danke schön!
\end{textebox}

\emph{Commentaire :} la salutation \all{Liebe Mitschülerinnen und Mitschüler} ouvre correctement le discours. Les connecteurs \all{zuerst}, \all{außerdem}, \all{schließlich} structurent les promesses. La subordonnée \all{Wenn ihr mich wählt, arbeite ich ...} place bien le verbe conjugué en tête de la principale (inversion). L'appel à l'action final \all{Wählt mich!} utilise l'impératif pluriel. Le Perfekt \all{habe ich ... organisiert} est correct pour raconter une expérience passée.

\textbf{2. Interview (journaliste)}

\begin{textebox}[Das Interview]
\textbf{Lukas:} Guten Tag, Frau Amina. Warum wollen Sie Klassensprecherin werden?

\textbf{Amina:} Ich will das Amt, weil ich meiner Klasse helfen möchte.

\textbf{Lukas:} Was wollen Sie zuerst für die Schüler tun?

\textbf{Amina:} Zuerst möchte ich eine Deutschgruppe organisieren.

\textbf{Lukas:} Wie wollen Sie die Probleme der Klasse lösen?

\textbf{Amina:} Ich spreche mit den Lehrern und der Direktion.

\textbf{Lukas:} Haben Sie schon Erfahrung?

\textbf{Amina:} Ja, letztes Jahr habe ich den Sporttag organisiert.

\textbf{Lukas:} Können Sie das genauer erklären?

\textbf{Amina:} Natürlich! Ich habe die Teams gebildet und das Programm geschrieben.

\textbf{Lukas:} Vielen Dank für das Interview!
\end{textebox}

\emph{Commentaire :} le journaliste vouvoie (\all{Sie}) et alterne \all{W-Fragen} (\all{warum}, \all{was}, \all{wie}) et question fermée avec inversion (\all{Haben Sie ...?}). La relance \all{Können Sie das genauer erklären?} montre l'écoute active. Notez le Datif après \all{helfen} : \all{meiner Klasse helfen}.

\textbf{3. Traduction des trois meilleures réponses (traducteur)}

\begin{center}
\begin{tabularx}{\linewidth}{|X|X|}
\hline
\rowcolor{popblueL} Deutsch & Français \\
\hline
\all{Ich will das Amt, weil ich meiner Klasse helfen möchte.} & Je veux ce poste parce que je voudrais aider ma classe. \\
\hline
\all{Zuerst möchte ich eine Deutschgruppe organisieren.} & D'abord, je voudrais organiser un groupe d'allemand. \\
\hline
\all{Ja, letztes Jahr habe ich den Sporttag organisiert.} & Oui, l'année dernière, j'ai organisé la journée du sport. \\
\hline
\end{tabularx}
\end{center}

\emph{Commentaire de traduction :} \all{das Amt} se rend ici par « ce poste » (et non « l'office ») ; \all{möchte} devient « je voudrais » (conditionnel français pour la politesse) ; \all{letztes Jahr} se place naturellement après « oui » en français. Le Perfekt allemand se traduit par le passé composé.
\end{exoresolu}

\courssub{Bilan de l'activité}

Cette activité a permis de réunir toutes les compétences du chapitre dans une tâche de communication authentique. Le \cle{candidat} a réinvesti la production écrite et orale : structurer un discours, employer les connecteurs, l'impératif et les subordonnées. Le \cle{journaliste} a pratiqué l'interaction orale : former des questions, vouvoyer, relancer. Le \cle{traducteur} a exercé la version allemand--français sur des phrases courtes, en veillant au naturel du français. L'évaluation croisée et le vote final ont développé l'écoute critique et la conscience des critères de qualité : lexique citoyen, place du verbe, cas et fluidité.

\begin{experience}[Prolongement possible]
Pour la séance suivante, le journal bilingue du lycée publie le meilleur duo : le discours et l'interview en allemand en page 1, la traduction française en page 2. Chaque groupe peut aussi enregistrer son interview sur téléphone et réécouter la place du verbe et la prononciation pour s'auto-corriger.
\end{experience}

\newpage

\coursec{Méthodes et exercices résolus}

\courssub{Méthode 1 : Rédiger une interview sur la citoyenneté}

\begin{methode}[Rédiger une interview (das Interview)]
Une interview est un dialogue écrit entre un journaliste (\all{der Journalist, -en}) et une personne interviewée. Pour la réussir :
\begin{enumerate}
\item \cle{Préparer le cadre} : présenter brièvement la personne (nom, âge, activité) en une ou deux phrases d'introduction.
\item \cle{Construire les questions} : utiliser les mots interrogatifs (\all{Wie? Was? Warum? Wann? Wo? Seit wann?}) avec le verbe en deuxième position, ou des questions fermées avec inversion du sujet (\all{Engagieren Sie sich...?}).
\item \cle{Soigner les réponses} : phrases complètes, verbe conjugué en position 2, vocabulaire du chapitre (\all{sich engagieren, helfen, die Freiwilligenarbeit, die Verantwortung}).
\item \cle{Varier les verbes de parole} : \all{antworten, erklären, meinen, finden, berichten} au lieu de répéter \all{sagen}.
\item \cle{Conclure} : une phrase de remerciement (\all{Vielen Dank für das Gespräch!}).
\item \cle{Relire} : majuscules aux noms, place du verbe, cas après les prépositions.
\end{enumerate}
Structure type : \all{Frage: Warum engagieren Sie sich? — Antwort: Ich engagiere mich, weil ich anderen helfen will.} « Question : Pourquoi vous engagez-vous ? — Réponse : Je m'engage parce que je veux aider les autres. »
\end{methode}

\begin{exoresolu}[Rédaction d'une interview]
Énoncé : Vous êtes journaliste au journal du lycée de Yaoundé. Rédigez une interview (6 échanges minimum) avec Amina, 17 ans, élève de Première, qui fait du bénévolat dans un orphelinat. Utilisez au moins trois mots interrogatifs différents et un verbe réfléchi.
\tcblower
Solution détaillée :

\textbf{Étape 1 — Plan.} Introduction (qui ?) → questions sur l'activité, les motivations, l'organisation, les difficultés, les projets → remerciement.

\textbf{Étape 2 — Choix grammaticaux.} Questions ouvertes : mot interrogatif + verbe + sujet. Verbe réfléchi : \all{sich engagieren} (ich engagiere mich, du engagierst dich). Subordonnée avec \all{weil} : verbe conjugué à la fin.

\textbf{Étape 3 — Rédaction.}

\all{Unsere Reporterin: Amina, du bist 17 Jahre alt und gehst in die erste Klasse. Seit wann arbeitest du im Waisenhaus?}

\all{Amina: Seit einem Jahr. Ich helfe dort jeden Samstag.}

\all{Reporterin: Was machst du genau mit den Kindern?}

\all{Amina: Ich spiele mit ihnen, ich helfe ihnen bei den Hausaufgaben, und manchmal koche ich auch.}

\all{Reporterin: Warum engagierst du dich?}

\all{Amina: Ich engagiere mich, weil ich anderen Menschen helfen will. Die Kinder brauchen Liebe und Aufmerksamkeit.}

\all{Reporterin: Ist es schwer, Schule und Freiwilligenarbeit zu verbinden?}

\all{Amina: Ja, manchmal. Aber ich lerne, meine Zeit gut zu organisieren.}

\all{Reporterin: Was sind deine Pläne für die Zukunft?}

\all{Amina: Ich möchte Sozialarbeiterin werden und weiter für Kinder arbeiten.}

\all{Reporterin: Vielen Dank für das Gespräch, Amina!}

\all{Amina: Gern geschehen!}

« Notre journaliste : Amina, tu as 17 ans et tu es en classe de Première. Depuis quand travailles-tu à l'orphelinat ? — Amina : Depuis un an. J'y aide tous les samedis. — Journaliste : Que fais-tu exactement avec les enfants ? — Amina : Je joue avec eux, je les aide à faire leurs devoirs, et parfois je cuisine aussi. — Journaliste : Pourquoi t'engages-tu ? — Amina : Je m'engage parce que je veux aider les autres. Les enfants ont besoin d'amour et d'attention. — Journaliste : Est-il difficile de combiner l'école et le bénévolat ? — Amina : Oui, parfois. Mais j'apprends à bien organiser mon temps. — Journaliste : Quels sont tes projets d'avenir ? — Amina : Je voudrais devenir assistante sociale et continuer à travailler pour les enfants. — Journaliste : Merci beaucoup pour cet entretien, Amina ! — Amina : De rien ! »

\textbf{Étape 4 — Vérification.} Trois interrogatifs différents (\all{Seit wann? Was? Warum?}) ; verbe réfléchi présent (\all{engagierst du dich}) ; subordonnée \all{weil} avec verbe final (\all{..., weil ich anderen Menschen helfen will}) ; tous les noms avec majuscule (\all{Waisenhaus, Kinder, Freiwilligenarbeit, Hausaufgaben}).

\textbf{Conclusion.} L'interview respecte le cadre demandé : 7 échanges, registre courtois, vocabulaire de la citoyenneté.
\end{exoresolu}

\courssub{Méthode 2 : Rédiger une lettre formelle}

\begin{methode}[Rédiger une lettre (der Brief)]
Pour écrire une lettre formelle (à un maire, une association, un journal) :
\begin{enumerate}
\item \cle{Lieu et date} en haut à droite : \all{Yaoundé, den 15. März 2025} (attention : \all{den} + date ordinale avec un point : \all{den 15. März} = « le 15 mars »).
\item \cle{Formule d'appel} : \all{Sehr geehrte Frau Mbarga,} / \all{Sehr geehrter Herr Bürgermeister,} / \all{Sehr geehrte Damen und Herren,} (virgule obligatoire, puis \textbf{minuscule} à la ligne suivante).
\item \cle{Introduction} : dire qui on est et pourquoi on écrit (\all{Ich schreibe Ihnen, weil...}).
\item \cle{Développement} : un paragraphe par idée, avec des connecteurs (\all{zuerst, dann, außerdem, deshalb, schließlich}).
\item \cle{Demande ou proposition} claire : \all{Ich bitte Sie um...} / \all{Wir schlagen vor, ... zu organisieren.}
\item \cle{Formule de politesse} : \all{Mit freundlichen Grüßen} (\textbf{sans virgule} !), puis la signature à la ligne.
\item \cle{Relire} : vouvoiement \all{Sie} et \all{Ihnen/Ihrer} avec majuscule partout, verbe en position 2, verbe final dans les subordonnées.
\end{enumerate}
\end{methode}

\begin{exoresolu}[Rédaction d'une lettre formelle]
Énoncé : Écrivez une lettre (80–100 mots) au maire de votre commune pour proposer une journée de nettoyage du quartier organisée par les élèves de votre lycée. Donnez la date, le programme et demandez du matériel.
\tcblower
Solution détaillée :

\textbf{Étape 1 — Éléments obligatoires.} Date, appel, présentation, proposition, programme, demande, formule finale.

\textbf{Étape 2 — Points de grammaire.} Subordonnées avec \all{weil} (verbe final) ; datif après \all{mit} (\all{mit Ihrer Hilfe}) ; accusatif après \all{um} (\all{um Material}) ; infinitive avec \all{zu} (\all{einen Reinigungstag zu organisieren}).

\textbf{Étape 3 — Rédaction.}

\all{Yaoundé, den 15. März 2025}

\all{Sehr geehrter Herr Bürgermeister,}

\all{ich bin Schülerin der ersten Klasse am Gymnasium Yaoundé. Ich schreibe Ihnen, weil meine Klasse und ich unser Viertel sauberer machen wollen.}

\all{Wir schlagen vor, am Samstag, den 29. März, einen Reinigungstag zu organisieren. Zuerst sammeln wir den Müll in den Straßen, dann putzen wir den Marktplatz, und schließlich pflanzen wir Blumen vor der Schule.}

\all{Wir bitten Sie um Material: Besen, Handschuhe und Müllsäcke. Mit Ihrer Hilfe können wir unser Viertel schöner und gesünder machen.}

\all{Mit freundlichen Grüßen}

\all{Clarisse Ngo Bell}

« Yaoundé, le 15 mars 2025. Monsieur le Maire, je suis élève de Première au lycée de Yaoundé. Je vous écris parce que ma classe et moi voulons rendre notre quartier plus propre. Nous proposons d'organiser une journée de nettoyage le samedi 29 mars. D'abord nous ramasserons les déchets dans les rues, ensuite nous nettoierons la place du marché, et enfin nous planterons des fleurs devant l'école. Nous vous demandons du matériel : balais, gants et sacs poubelles. Avec votre aide, nous pouvons rendre notre quartier plus beau et plus sain. Cordialement, Clarisse Ngo Bell. »

\textbf{Étape 4 — Vérification.} \all{den 15. März} correct ; minuscule après la virgule de l'appel (\all{ich bin...}) ; \all{Herr Bürgermeister} avec majuscules ; subordonnée \all{weil... wollen} avec verbe final ; \all{um} + accusatif (\all{um Material}) ; \all{mit} + datif (\all{mit Ihrer Hilfe}) ; environ 95 mots.

\textbf{Conclusion.} La lettre est complète, polie et grammaticalement correcte : c'est le modèle à imiter.
\end{exoresolu}

\courssub{Méthode 3 : Rédiger un petit discours}

\begin{methode}[Rédiger un petit discours (die Rede)]
Un discours s'adresse directement à un public. Démarche :
\begin{enumerate}
\item \cle{Saluer le public} : \all{Liebe Mitschülerinnen und Mitschüler,} / \all{Sehr geehrte Gäste,}
\item \cle{Annoncer le sujet} : \all{Heute spreche ich über das Thema Zusammenleben.} « Aujourd'hui je parle du thème de la vie en société. »
\item \cle{Développer 2 ou 3 arguments} avec des exemples concrets (école, quartier, pays) et des connecteurs (\all{erstens, zweitens, außerdem}).
\item \cle{Utiliser le « nous »} (\all{wir}) pour créer la solidarité : \all{Wir müssen zusammenarbeiten.} « Nous devons travailler ensemble. »
\item \cle{Lancer un appel à l'action} : \all{Lasst uns gemeinsam handeln!} « Agissons ensemble ! » (impératif à la 1\iere{} personne du pluriel : \all{lasst uns} + infinitif).
\item \cle{Remercier} : \all{Ich danke euch für eure Aufmerksamkeit.} « Je vous remercie de votre attention. »
\end{enumerate}
Réflexes : phrases courtes, verbes modaux (\all{müssen, können, sollen}) pour convaincre, répétitions volontaires pour le rythme.
\end{methode}

\begin{exoresolu}[Rédaction d'un petit discours]
Énoncé : Pour la Journée de la citoyenneté de votre lycée, rédigez un discours de 60–80 mots intitulé \all{„Jugendliche für eine bessere Stadt“}. Saluez le public, donnez deux arguments et terminez par un appel à l'action.
\tcblower
Solution détaillée :

\textbf{Étape 1 — Structure.} Salutation → annonce du sujet → argument 1 (propreté) → argument 2 (solidarité) → appel → remerciement.

\textbf{Étape 2 — Grammaire.} Impératif \all{lasst uns} + infinitif ; modaux \all{können, müssen} avec infinitif en fin de phrase ; possessifs simples (\all{unsere Stadt, unser Zuhause}).

\textbf{Étape 3 — Rédaction.}

\all{Liebe Mitschülerinnen und Mitschüler, liebe Lehrer,}

\all{heute spreche ich über ein wichtiges Thema: Jugendliche für eine bessere Stadt.}

\all{Erstens können wir unsere Stadt sauber halten. Wir dürfen den Müll nicht auf die Straße werfen, sondern wir müssen ihn in die Mülleimer werfen. Zweitens können wir alten Menschen und Kindern helfen. Ein freundliches Wort kostet nichts, aber es bedeutet viel.}

\all{Unsere Stadt ist unser Zuhause. Lasst uns gemeinsam handeln! Lasst uns unsere Stadt jeden Tag ein bisschen besser machen!}

\all{Ich danke euch für eure Aufmerksamkeit.}

« Chers camarades, chers professeurs, aujourd'hui je parle d'un sujet important : les jeunes pour une ville meilleure. Premièrement, nous pouvons garder notre ville propre. Nous ne devons pas jeter les déchets dans la rue, mais dans les poubelles. Deuxièmement, nous pouvons aider les personnes âgées et les enfants. Un mot gentil ne coûte rien, mais il signifie beaucoup. Notre ville est notre maison. Agissons ensemble ! Rendons notre ville chaque jour un peu meilleure ! Je vous remercie de votre attention. »

\textbf{Étape 4 — Vérification.} Environ 80 mots ; deux arguments marqués par \all{erstens / zweitens} ; appel avec \all{Lasst uns...!} ; verbe en position 2 partout ; noms avec majuscule (\all{Stadt, Müll, Mülleimer, Zuhause, Aufmerksamkeit}) ; \all{helfen} + datif (\all{alten Menschen und Kindern}).

\textbf{Conclusion.} Le discours est bref, rythmé et convaincant : exactement ce que l'examinateur attend.
\end{exoresolu}

\courssub{Méthode 4 : Traduire du français vers l'allemand (thème)}

\begin{methode}[Réussir le thème (français → allemand)]
\begin{enumerate}
\item \cle{Lire toute la phrase} et repérer le verbe principal et son temps (présent → \all{Präsens}, passé composé → \all{Perfekt}, futur proche → présent avec complément de temps).
\item \cle{Identifier le sujet} et placer le verbe conjugué en position 2 : \all{Sujet + verbe + ...} ; si la phrase commence par un complément (\all{Gestern...}), le sujet passe après le verbe : \all{Gestern habe ich...}.
\item \cle{Traduire les compléments} en respectant les cas : accusatif pour le COD (\all{Ich sehe den Mann}), datif pour le COI et après \all{helfen, danken} (\all{Ich helfe dem Kind}).
\item \cle{Subordonnées} : conjonction + verbe conjugué à la fin (\all{..., weil ich meiner Mutter helfe.}).
\item \cle{Verbes à particule} : particule détachée en fin de proposition principale (\all{Er steht um 7 Uhr auf.}).
\item \cle{Relire} : majuscules aux noms, genre et cas des articles, Umlaute, \all{ß/ss}.
\end{enumerate}
\end{methode}

\begin{exoresolu}[Thème : phrases sur la citoyenneté]
Énoncé : Traduisez en allemand :
1. Les jeunes doivent respecter les lois de leur pays.
2. Hier, nous avons aidé une vieille dame au marché.
3. Je vote parce que c'est un droit et un devoir.
4. Ma sœur se lève tôt pour travailler avec les bénévoles.
\tcblower
Solution détaillée :

\textbf{Phrase 1.} « doivent » → modal \all{müssen} (verbe conjugué en position 2) + infinitif \all{respektieren} à la fin. « les lois » = COD pluriel → accusatif \all{die Gesetze} (die Gesetz, -e). « de leur pays » → génitif \all{ihres Landes} (das Land). Réponse : \all{Die Jugendlichen müssen die Gesetze ihres Landes respektieren.}

\textbf{Phrase 2.} Passé composé → \all{Perfekt} avec l'auxiliaire \all{haben} (verbe transitif) : \all{wir haben ... geholfen}. Attention : \all{helfen} exige le \textbf{datif} ! « une vieille dame » → \all{einer alten Dame}. « au marché » → \all{auf dem Markt} (datif, lieu). La phrase commence par « hier » → inversion du sujet. Réponse : \all{Gestern haben wir einer alten Dame auf dem Markt geholfen.}

\textbf{Phrase 3.} Subordonnée avec \all{weil} : verbe conjugué \all{ist} à la fin. « un droit et un devoir » → \all{ein Recht und eine Pflicht} (das Recht, -e ; die Pflicht, -en). Réponse : \all{Ich wähle, weil es ein Recht und eine Pflicht ist.}

\textbf{Phrase 4.} « se lever » = \all{aufstehen} : verbe à particule \textbf{non réfléchi} en allemand (ne pas ajouter \all{sich} !). Particule \all{auf} détachée en fin de principale : \all{Meine Schwester steht früh auf}. « pour travailler » → \all{um ... zu} + infinitif ; « avec les bénévoles » → datif pluriel \all{mit den Freiwilligen}. Réponse : \all{Meine Schwester steht früh auf, um mit den Freiwilligen zu arbeiten.}

\textbf{Conclusion.} Chaque phrase illustre un piège classique : modal + infinitif final, datif après \all{helfen}, verbe final après \all{weil}, particule séparable et tournure \all{um ... zu}.
\end{exoresolu}

\courssub{Méthode 5 : Traduire de l'allemand vers le français (version)}

\begin{methode}[Réussir la version (allemand → français)]
\begin{enumerate}
\item \cle{Repérer le verbe conjugué} (position 2) et le sujet : c'est le squelette de la phrase.
\item \cle{Aller chercher l'infinitif ou le participe en fin de phrase} (modaux, Perfekt, Futur) avant de traduire : \all{Er hat den Text gestern geschrieben} = « Il a écrit le texte hier ».
\item \cle{Subordonnées} : identifier la conjonction (\all{weil, dass, wenn, obwohl}) et le verbe final, puis traduire la subordonnée à sa place logique en français.
\item \cle{Attention aux faux amis} : \all{bekommen} = recevoir (pas « devenir »), \all{aktuell} = actuel, \all{der Chef} = le patron, \all{die Gesellschaft} = la société.
\item \cle{Traduire les cas} par les prépositions françaises : datif → « à », génitif → « de ».
\item \cle{Rédiger un français naturel}, pas un calque mot à mot.
\end{enumerate}
\end{methode}

\begin{exoresolu}[Version : court texte sur l'engagement]
Énoncé : Traduisez en français le texte suivant :

\all{Viele junge Leute in Deutschland engagieren sich in Vereinen. Sie helfen zum Beispiel alten Menschen oder sammeln Geld für Kinder. Lukas, 17 Jahre alt, sagt: „Ich habe letztes Jahr gelernt, dass man mit wenig Zeit viel bewegen kann. Jeder kann etwas für die Gesellschaft tun.“}

\tcblower
Solution détaillée :

\textbf{Analyse phrase par phrase.}

Phrase 1 : \all{engagieren sich} = « s'engagent » (verbe réfléchi) ; \all{in Vereinen} = « dans des associations » (datif pluriel après \all{in}, der Verein, -e). Traduction : « Beaucoup de jeunes en Allemagne s'engagent dans des associations. »

Phrase 2 : \all{zum Beispiel} = « par exemple » ; \all{alten Menschen} = datif pluriel après \all{helfen} → « aident des personnes âgées » ; \all{Geld sammeln für} = « collectent de l'argent pour ». Traduction : « Ils aident par exemple des personnes âgées ou collectent de l'argent pour les enfants. »

Phrase 3 : \all{letztes Jahr} = « l'année dernière » ; \all{ich habe gelernt, dass...} = Perfekt + subordonnée \all{dass} avec verbe final (\all{..., dass man ... viel bewegen kann}) ; \all{mit wenig Zeit viel bewegen} = « avec peu de temps, faire bouger beaucoup de choses » ; \all{die Gesellschaft} = « la société » (faux ami : pas « la compagnie »). Traduction : « Lukas, 17 ans, dit : “L'année dernière, j'ai appris qu'avec peu de temps on peut faire bouger beaucoup de choses. Chacun peut faire quelque chose pour la société.” »

\textbf{Conclusion.} La version exige de reconstruire chaque phrase autour du verbe conjugué, puis de rédiger un français fluide sans calquer l'ordre allemand.
\end{exoresolu}

\courssub{Méthode 6 : Relire et corriger sa production}

\begin{methode}[La relecture efficace en 5 passages]
Ne relisez jamais « en général » : faites \cle{cinq passages ciblés}.
\begin{enumerate}
\item \cle{Passage 1 — Le verbe} : est-il en position 2 dans chaque principale ? À la fin dans chaque subordonnée ? Particules séparables détachées ?
\item \cle{Passage 2 — Les majuscules} : tous les noms commencent-ils par une majuscule ? \all{Sie} (vous) avec majuscule ?
\item \cle{Passage 3 — Les cas} : accusatif après \all{für, durch, ohne, um} ; datif après \all{mit, zu, von, bei, nach, aus} ; datif après \all{helfen, danken, gefallen}.
\item \cle{Passage 4 — L'orthographe} : Umlaute (\all{ä, ö, ü}), \all{ß} après voyelle longue (\all{die Straße}), \all{ss} après voyelle courte (\all{der Fluss}), \all{sch/sp/st}.
\item \cle{Passage 5 — Le sens} : ai-je répondu à la consigne ? Nombre de mots respecté ? Faux amis évités ?
\end{enumerate}
Réflexe d'examen : garder 5 minutes pour la relecture, c'est souvent 2 points gagnés.
\end{methode}

\begin{exoresolu}[Correction d'un texte fautif]
Énoncé : Le texte suivant contient 10 erreurs (place du verbe, cas, majuscules, verbe réfléchi). Trouvez-les et corrigez-les.

\all{ich bin schüler in Yaoundé und ich engagiere in einem verein. Letzte Woche ich habe ein alten Mann geholfen. Wir haben geld für die kinder gesammelt, weil sie brauchen hilfe. Mein freund sagt, dass bürger müssen zusammenarbeiten.}
\tcblower
Solution détaillée — les 10 erreurs :

\begin{enumerate}
\item \all{ich} → \all{Ich} : majuscule en début de phrase.
\item \all{schüler} → \all{Schüler} : majuscule au nom.
\item \all{ich engagiere} → \all{ich engagiere mich} : \all{sich engagieren} est un verbe réfléchi.
\item \all{verein} → \all{Verein} : majuscule au nom.
\item \all{Letzte Woche ich habe} → \all{Letzte Woche habe ich} : le verbe conjugué reste en position 2 ; après un complément de temps initial, le sujet passe après le verbe (inversion).
\item \all{ein alten Mann} → \all{einem alten Mann} : \all{helfen} exige le datif (\all{einem alten Mann}).
\item \all{geld} → \all{Geld} : majuscule au nom.
\item \all{kinder} → \all{Kinder} : majuscule au nom.
\item \all{weil sie brauchen hilfe} → \all{weil sie Hilfe brauchen} : verbe conjugué à la fin dans la subordonnée \all{weil} + majuscule au nom \all{Hilfe}.
\item \all{Mein freund sagt, dass bürger müssen zusammenarbeiten} → \all{Mein Freund sagt, dass Bürger zusammenarbeiten müssen} : majuscules (\all{Freund, Bürger}) et, dans la subordonnée \all{dass}, l'infinitif précède le modal à la toute fin (\all{zusammenarbeiten müssen}).
\end{enumerate}

\textbf{Texte corrigé :} \all{Ich bin Schüler in Yaoundé und ich engagiere mich in einem Verein. Letzte Woche habe ich einem alten Mann geholfen. Wir haben Geld für die Kinder gesammelt, weil sie Hilfe brauchen. Mein Freund sagt, dass Bürger zusammenarbeiten müssen.}

« Je suis élève à Yaoundé et je m'engage dans une association. La semaine dernière, j'ai aidé un vieil homme. Nous avons collecté de l'argent pour les enfants parce qu'ils ont besoin d'aide. Mon ami dit que les citoyens doivent travailler ensemble. »

\textbf{Conclusion.} Une relecture méthodique par catégories (verbe, majuscules, cas, orthographe, sens) permet de repérer toutes les fautes, même les plus discrètes.
\end{exoresolu}

\begin{attention}[Les 5 erreurs qui coûtent des points au Bac]
\begin{enumerate}
\item \cle{Oublier la majuscule aux noms} : \all{die schule, das kind} sont fautifs. En allemand, \textbf{tous} les noms prennent une majuscule : \all{die Schule, das Kind, die Freiheit}. C'est la faute la plus sanctionnée.
\item \cle{Mal placer le verbe} : ne dites jamais \all{Gestern ich habe...} mais \all{Gestern habe ich...} (verbe en position 2). Dans une subordonnée (\all{weil, dass, wenn}), le verbe conjugué va \textbf{à la fin} : \all{..., weil ich meiner Mutter helfe.}
\item \cle{Confondre les cas} : \all{helfen, danken, gefallen} exigent le \textbf{datif} : \all{Ich helfe dem Kind} (et non \all{das Kind}). Après \all{mit, zu, von, bei} : datif ; après \all{für, ohne, um, durch} : accusatif.
\item \cle{Calquer le français} : « Je suis d'accord avec toi » se dit \all{Ich bin mit dir einverstanden} ou \all{Ich bin deiner Meinung} ; « avoir 17 ans » = \all{17 Jahre alt sein} (avec \all{sein}, jamais \all{haben}) ; « se lever » = \all{aufstehen}, sans réfléchi.
\item \cle{Tomber dans les faux amis} : \all{bekommen} = recevoir (devenir = \all{werden}) ; \all{aktuell} = actuel ; \all{die Gesellschaft} = la société ; \all{der Chef} = le patron. Vérifiez toujours avant d'écrire.
\end{enumerate}
\end{attention}

\newpage

\coursec{Exercices}

\courssub{Je vérifie mes connaissances}

\begin{exercice}[QCM : grammaire et vocabulaire de la citoyenneté]
Choisissez la bonne réponse parmi les trois propositions.
\begin{enumerate}
    \item Pour exprimer une hypothèse au présent (Konjunktiv II), on utilise souvent la forme du \emph{Präteritum} avec umlaut. Quelle est la forme correcte pour \all{können} (il/elle) ?
    \begin{enumerate}
        \item \all{könnte}
        \item \all{könnt}
        \item \all{konnte}
    \end{enumerate}
    \item Dans une lettre formelle au \all{Bürgermeister}, quelle formule d'appel est correcte ?
    \begin{enumerate}
        \item \all{Lieber Herr Müller,}
        \item \all{Sehr geehrter Herr Bürgermeister,}
        \item \all{Hallo Bürgermeister,}
    \end{enumerate}
    \item Le mot \all{die Staatsbürgerschaft} signifie :
    \begin{enumerate}
        \item la nationalité / la citoyenneté
        \item le parlement
        \item l'élection
    \end{enumerate}
    \item Complétez avec la préposition correcte (et le bon cas) : \all{Wir diskutieren \_\_\_\_\_ die neuen Gesetze.} (Nous discutons \emph{des} nouvelles lois.)
    \begin{enumerate}
        \item \all{über} + Akkusativ
        \item \all{mit} + Dativ
        \item \all{an} + Dativ
    \end{enumerate}
    \item Transformez au discours indirect (Konjunktiv I) : \all{Er sagt: „Ich engagiere mich.“} $\rightarrow$ \all{Er sagt, er \_\_\_\_\_ sich.}
    \begin{enumerate}
        \item \all{engagiert}
        \item \all{engagiere}
        \item \all{engagierte}
    \end{enumerate}
\end{enumerate}
\end{exercice}

\begin{exercice}[Vrai ou Faux ? Justifiez en allemand]
Lisez le court texte suivant et décidez si les affirmations sont \all{richtig} ou \all{falsch}. Justifiez votre réponse par une phrase complète en allemand (citation ou reformulation).

\begin{textebox}[Kurzer Text: Jugendparlament]
In vielen deutschen Städten gibt es Jugendparlamente. Junge Menschen zwischen 14 und 18 Jahren können dort ihre Ideen für die Stadt einbringen. Sie werden nicht vom Bürgermeister ernannt, sondern von anderen Jugendlichen gewählt. Das Parlament trifft sich einmal im Monat im Rathaus.
\end{textebox}

\begin{enumerate}
    \item \all{Das Jugendparlament besteht nur aus Erwachsenen.}
    \item \all{Die Mitglieder werden von den Jugendlichen gewählt.}
    \item \all{Das Parlament trifft sich jede Woche.}
    \item \all{Junge Menschen können ihre Ideen für die Stadt einbringen.}
\end{enumerate}
\end{exercice}

\begin{exercice}[Vocabulaire thématique : Citoyenneté et entretien]
Complétez le tableau en associant le mot allemand (avec article et pluriel) à sa définition française. Un mot allemand est déjà donné en exemple.

\begin{center}
\begin{tabularx}{\linewidth}{|X|X|X|}
\hline
\rowcolor{popblueL} \textbf{Deutsch (Artikel, Plural)} & \textbf{Français} & \textbf{Contexte / Exemple} \\
\hline
\all{die Wahl, die Wahlen} & l'élection & \all{Die Wahl findet am Sonntag statt.} \\
\hline
& le droit de vote & \\
\hline
\all{der Abgeordnete, die Abgeordneten} & & \all{Der Abgeordnete hält eine Rede.} \\
\hline
& l'engagement (bénévole, civique) & \\
\hline
\all{die Versammlung, die Versammlungen} & & \all{Die Versammlung beginnt um 18 Uhr.} \\
\hline
& le bulletin de vote & \\
\hline
\all{der Wahlkreis, die Wahlkreise} & & \\
\hline
\end{tabularx}
\end{center}
\end{exercice}

\begin{exercice}[Conjugaison à compléter : révision des temps et des modes]
Conjuguez les verbes entre parenthèses au temps ou au mode indiqué (Präsens, Perfekt, Präteritum, Futur I, Konjunktiv II). Attention aux verbes à particule séparable et aux verbes forts.

\begin{enumerate}
    \item Wenn ich Zeit \all{haben} (Konjunktiv II), \all{sich engagieren} (Konjunktiv II, ich) in einem Verein.
    \item Der Bürgermeister \all{eröffnen} (Präteritum) gestern das neue Jugendzentrum.
    \item Wir \all{diskutieren} (Perfekt) schon lange über das Wahlalter.
    \item \all{Können} (Konjunktiv II, Sie) Sie mir bitte helfen, den Brief zu formulieren?
    \item Die Schüler \all{müssen} (Präsens) zu der Versammlung \all{gehen} (Infinitiv).
    \item Ich \all{schreiben} (Futur I) morgen den Brief an die Schulleitung.
\end{enumerate}
\end{exercice}

\courssub{Je m'entraîne}

\begin{exercice}[Traduction : phrases courtes (thème et version)]
Traduisez les phrases suivantes. Respectez l'ordre des mots (place du verbe, compléments), les majuscules aux noms et la ponctuation allemande.

\textbf{Thème (Français $\rightarrow$ Allemand) :}
\begin{enumerate}
    \item Les jeunes s'engagent pour l'environnement.
    \item Avez-vous déjà voté à une élection ?
    \item Si j'étais maire, je construirais un centre de jeunesse. (Konjunktiv II)
    \item On dit que le parlement se réunit demain. (Discours indirect / Konjunktiv I)
    \item Je lui ai écrit une lettre formelle hier.
\end{enumerate}

\textbf{Version (Allemand $\rightarrow$ Français) :}
\begin{enumerate}
    \item \all{Das Wahlalter wurde auf 16 Jahre gesenkt.}
    \item \all{Wir müssen uns für unsere Rechte einsetzen.}
    \item \all{Der Reporter fragt den Kandidaten, ob er gewählt werden will.}
    \item \all{Nach der Versammlung haben wir einen Antrag gestellt.}
    \item \all{Würdest du dich als Klassensprecher aufstellen lassen?}
\end{enumerate}
\end{exercice}

\begin{exercice}[Transformation : discours direct $\rightarrow$ discours indirect]
Transformez les questions et réponses de cet entretien (discours direct) en un compte rendu au discours indirect (subordonnées avec \all{dass} ou \all{ob}, Konjunktiv I). Adaptez les pronoms et les possessifs.

\begin{textebox}[Interview: Direkte Rede]
Reporter: „Warum engagieren Sie sich in der Lokalpolitik?“
Kandidat: „Ich möchte die Situation für junge Menschen verbessern.“
Reporter: „Haben Sie konkrete Pläne?“
Kandidat: „Ja, ich will einen Jugendrat gründen.“
Reporter: „Werden Sie bei der nächsten Wahl antreten?“
Kandidat: „Das habe ich noch nicht entschieden.“
\end{textebox}

\textit{Consigne :} rédigez le compte rendu en allemand. Début : \all{Der Reporter fragt den Kandidaten, warum er sich in der Lokalpolitik engagiere. Der Kandidat antwortet, dass er\ldots}
\end{exercice}

\begin{exercice}[Texte à trous : lettre formelle au Bürgermeister]
Complétez la lettre suivante avec les éléments grammaticaux demandés entre parenthèses (cas, prépositions, conjugaison, pronoms relatifs, connecteurs logiques).

\begin{textebox}[Lückentext: Brief an den Bürgermeister]
Sehr geehrter Herr Bürgermeister,

hiermit möchte ich mich \_\_\_\_\_ (Präposition: bei) Ihnen für das Gespräch neulich bedanken. Bei dem Treffen, \_\_\_\_\_ (Relativpronomen, Akkusativ, Neutrum) wir hatten, ging es um einen Schüleraustausch zwischen unserem Gymnasium in Yaoundé und einer Schule in Bonn.

Ich \_\_\_\_\_ (Konjunktiv II: bitten) Sie, die Möglichkeit einer Partnerschaft zu prüfen. Viele Schüler \_\_\_\_\_ (Perfekt: sich interessieren) schon lange für Deutschland. Wenn der Austausch \_\_\_\_\_ (Konjunktiv II: stattfinden), \_\_\_\_\_ (Konjunktiv II: profitieren) beide Seiten davon.

\_\_\_\_\_ (Präposition: auf) eine positive Antwort würde ich mich sehr freuen.

Mit freundlichen Grüßen

Moussa Njoya
\end{textebox}
\end{exercice}

\begin{exercice}[Appariements : situations de communication et formulations]
Associez chaque situation de communication (1--6) à la formulation allemande la plus appropriée (A--F). Écrivez les paires (ex. 1--B).

\begin{center}
\begin{tabularx}{\linewidth}{|X|X|}
\hline
\rowcolor{popblueL} \textbf{Situation} & \textbf{Formulation allemande} \\
\hline
1. Débuter un petit discours officiel devant le conseil municipal. & A. \all{Könnten Sie mir bitte sagen, was Ihre Meinung zum Wahlalter ist?} \\
\hline
2. Poser une question polie lors d'une interview (vouvoiement). & B. \all{Meine sehr geehrten Damen und Herren, liebe Mitschülerinnen und Mitschüler, ...} \\
\hline
3. Exprimer son accord argumenté dans un débat. & C. \all{Ich bin ganz Ihrer Meinung, denn die Jugend ist die Zukunft.} \\
\hline
4. Refuser poliment une invitation à se présenter aux élections. & D. \all{Das ist ein interessanter Punkt, aber ich sehe das anders.} \\
\hline
5. Exprimer un désaccord nuancé. & E. \all{Ich bedanke mich für das Vertrauen, muss aber ablehnen.} \\
\hline
6. Conclure une lettre formelle de demande. & F. \all{Ich freue mich auf Ihre Antwort und verbleibe mit freundlichen Grüßen.} \\
\hline
\end{tabularx}
\end{center}
\end{exercice}

\courssub{Je me prépare au Bac}

\begin{exercice}[Compréhension et analyse linguistique : les Jugendparlamente (type Bac)]
Lisez attentivement le texte suivant (environ 160 mots) et répondez aux questions en allemand (questions 1 à 3) ou en français (questions 4 à 7).

\begin{textebox}[Jugendparlamente: Eine Stimme für die Jugend]
In Deutschland gibt es in vielen Städten und Gemeinden Jugendparlamente, Jugendbeiräte oder Schülervertretungen. Diese Gremien geben jungen Menschen die Möglichkeit, politische Prozesse kennenzulernen und eigene Interessen einzubringen. Die Mitglieder sind meist zwischen 14 und 21 Jahre alt und werden für ein oder zwei Jahre gewählt. Wahlberechtigt sind oft alle Jugendlichen, die im jeweiligen Wahlkreis wohnen, unabhängig von der Staatsbürgerschaft.

Ein Beispiel ist das Jugendparlament in Bonn. Es beschließt einen Haushaltsplan, den die Mitglieder selbst aufstellen, und kann Anträge in den Stadtrat einbringen. Kürzlich hat es sich erfolgreich für einen neuen Skatepark und für kostenloses WLAN in öffentlichen Parks eingesetzt. Es spricht auch Jugendliche an, denen Politik sonst fern erscheint. „Wir werden endlich gehört“, sagt die 17-jährige Sprecherin Lena Meyer. „Politik ist nicht nur etwas für Alte.“

Kritiker bemängeln jedoch, dass die Parlamente oft nur beratende Funktion haben und keine echten Entscheidungsbefugnisse besitzen. Dennoch: Wer früh lernt, Verantwortung zu übernehmen, stärkt die Demokratie von morgen.
\end{textebox}

\textbf{Questions de compréhension (réponses en allemand, phrases complètes) :}
\begin{enumerate}
    \item \all{Wie alt sind die Mitglieder der Jugendparlamente meistens und wie lange werden sie gewählt?}
    \item \all{Welche Erfolge hat das Jugendparlament in Bonn kürzlich erzielt? Nennen Sie zwei Beispiele.}
    \item \all{Was kritisieren die Kritiker an den Jugendparlamenten?}
\end{enumerate}

\textbf{Questions de langue (réponses en français) :}
\begin{enumerate}
    \setcounter{enumi}{3}
    \item Citez dans le texte \textbf{deux verbes conjugués au Perfekt} et donnez leur infinitif ainsi que leur traduction française.
    \item Trouvez dans le texte \textbf{une subordonnée relative au cas accusatif} et \textbf{une subordonnée relative au cas datif}. Recopiez-les, soulignez le pronom relatif et indiquez son antécédent.
    \item Traduisez en français le passage : \all{„Wir werden endlich gehört“, sagt die 17-jährige Sprecherin Lena Meyer. „Politik ist nicht nur etwas für Alte.“}
    \item La dernière phrase du texte est au présent de l'indicatif. Reformulez-la en une phrase hypothétique avec \all{wenn} + \textbf{Konjunktiv II} (modèle : \all{Wenn man früh Verantwortung übernehmen würde, ...}).
\end{enumerate}
\end{exercice}

\begin{exercice}[Traduction thème et version (type Bac)]

\textbf{Partie A : Thème (Français $\rightarrow$ Allemand)}
Traduisez les 5 phrases suivantes en allemand correct (orthographe, cas, ordre des mots, temps). Chaque phrase vaut 2 points.

\begin{enumerate}
    \item La citoyenneté active renforce la démocratie dans notre pays.
    \item Depuis l'année dernière, les jeunes de 16 ans ont le droit de vote aux élections communales.
    \item Le maire a demandé aux élèves s'ils voulaient organiser un projet pour l'environnement.
    \item Si les politiciens écoutaient davantage la jeunesse, beaucoup de problèmes seraient résolus. (Konjunktiv II)
    \item Nous avons décidé d'écrire une lettre au ministre de l'Éducation pour proposer un échange scolaire.
\end{enumerate}

\textbf{Partie B : Version (Allemand $\rightarrow$ Français)}
Traduisez en français le passage suivant (extrait d'un article de presse sur le vote à 16 ans).

\begin{textebox}[Auszug: Wahlalter 16]
Die Debatte über die Senkung des Wahlalters auf 16 Jahre gewinnt an Fahrt. Befürworter argumentieren, dass Jugendliche bereits mit 16 Jahren strafmündig sind und Steuern zahlen, wenn sie arbeiten. Daher sei es nur gerecht, ihnen auch das Wahlrecht zu geben. Studien zeigen, dass die Wahlbeteiligung bei Erstwählern oft höher ist als bei älteren Altersgruppen. Gegner hingegen fürchten, dass 16-Jährige noch zu leicht beeinflussbar seien und die Tragweite politischer Entscheidungen nicht voll erfassen könnten. In mehreren Bundesländern ist die Wahl ab 16 bei Kommunal- und Landtagswahlen bereits Realität. Auf Bundesebene hingegen gilt weiterhin das Mindestalter von 18 Jahren für die Bundestagswahl.
\end{textebox}
\end{exercice}

\begin{exercice}[Production écrite guidée : lettre formelle / petit discours (type Bac)]
Choisissez \textbf{UN} des deux sujets suivants. Respectez scrupuleusement les conventions du genre épistolaire ou du discours, le registre de langue (formel), la structure (introduction, développement, conclusion) et le nombre de lignes indiqué (environ 15--18 lignes manuscrites $\approx$ 120--150 mots). Écrivez \textbf{uniquement en allemand}.

\textbf{Sujet 1 : Lettre formelle (Brief an den Bürgermeister)}

Vous êtes délégué(e) de la classe de Première A du Lycée de Yaoundé. Écrivez une lettre formelle au \all{Bürgermeister} de Bonn (Allemagne) pour solliciter un partenariat d'échange scolaire entre votre lycée et un \all{Gymnasium} de Bonn. Votre lettre doit contenir :
\begin{itemize}
    \item L'en-tête (expéditeur, destinataire, date, lieu) et l'objet (\all{Betreff}).
    \item La présentation de votre établissement et de votre classe.
    \item Les objectifs de l'échange (découverte culturelle, pratique de l'allemand, citoyenneté).
    \item Une proposition concrète (durée, période, nombre d'élèves).
    \item Une formule de politesse de conclusion adaptée.
\end{itemize}

\textbf{Sujet 2 : Petit discours (Kurze Rede)}

Vous êtes candidat(e) au poste de \all{Schulsprecher(in)} (délégué(e) général(e) de l'établissement). Prononcez un discours de 2 minutes devant l'assemblée des élèves pour vous faire élire. Votre discours doit contenir :
\begin{itemize}
    \item Une accroche (salutation formelle et informelle).
    \item Votre motivation personnelle (pourquoi vous ?).
    \item Trois propositions concrètes pour améliorer la vie scolaire (ex. : cantine, clubs, représentation auprès de la direction, événements culturels).
    \item Un appel au vote (\all{Stimmt für mich!}).
    \item Une conclusion engageante.
\end{itemize}
\end{exercice}

\begin{methode}[Réussir la lettre formelle et le petit discours]
\begin{enumerate}
    \item \textbf{Lettre formelle} : en-tête en haut à gauche (expéditeur), destinataire en dessous, lieu et date à droite (\all{Yaoundé, den 12. März 2025}), objet sans article après \all{Betreff:}, appel \all{Sehr geehrter Herr Bürgermeister,} (virgule, puis minuscule), paragraphes courts, conclusion \all{Mit freundlichen Grüßen} (sans virgule) + signature.
    \item \textbf{Discours} : accroche (\all{Meine sehr geehrten Damen und Herren, liebe Mitschülerinnen und Mitschüler,}), annonce du plan, propositions au Futur I ou avec \all{wollen}, phrases courtes et rythmées, appel au vote, remerciement final (\all{Ich danke euch für eure Aufmerksamkeit.}).
    \item \textbf{Langue} : vouvoiement \all{Sie} dans la lettre, \all{ihr/euch} possible dans le discours devant les élèves ; Konjunktiv II pour les demandes polies (\all{Ich würde mich sehr freuen, ...}).
    \item \textbf{Relecture} : verbe en 2\textsuperscript{e} position, infinitif à la fin avec \all{würde/werden}, majuscules aux noms, cas après prépositions (\all{für, mit, bei, von}).
\end{enumerate}
\end{methode}

\begin{exemplebox}[Modèle de début de lettre formelle]
\all{Moussa Njoya, Klassensprecher der Klasse 1\textsuperscript{ère} A, Lycée de Yaoundé, Yaoundé, Kamerun}

\all{An den Bürgermeister der Stadt Bonn, Rathaus, Bonn, Deutschland}

\all{Yaoundé, den 12. März 2025}

\all{Betreff: Partnerschaft für einen Schüleraustausch}

\all{Sehr geehrter Herr Bürgermeister,}

\all{als Klassensprecher der Klasse 1\textsuperscript{ère} A des Lycée de Yaoundé schreibe ich Ihnen, um eine Partnerschaft zwischen unserer Schule und einem Gymnasium in Bonn vorzuschlagen. Unsere Klasse lernt seit drei Jahren Deutsch und interessiert sich sehr für die deutsche Kultur und das demokratische Leben in Ihrem Land.}

\all{Ein Austausch von zwei Wochen im nächsten Schuljahr mit etwa zwanzig Schülern würde beiden Seiten viele Vorteile bringen: Die Schüler könnten ihre Sprachkenntnisse verbessern, andere Kulturen entdecken und als junge Bürger Verantwortung übernehmen.}

\all{Über eine positive Antwort würde ich mich sehr freuen.}

\all{Mit freundlichen Grüßen}

\all{Moussa Njoya}
\end{exemplebox}

\begin{attention}[Erreurs fréquentes des francophones]
\begin{itemize}
    \item \textbf{Place du verbe} : ne dites pas \all{*Ich habe geschrieben einen Brief} mais \all{Ich habe einen Brief geschrieben} (participe à la fin).
    \item \textbf{Discours indirect} : \all{Er sagt, dass er sich engagiere} (Konjunktiv I), pas \all{*dass er sich engagiert} si l'on veut marquer le discours rapporté.
    \item \textbf{Faux amis} : \all{die Wahl} = l'élection, mais \all{die Wahl treffen} = faire un choix ; \all{der Antrag} = la motion/la demande officielle, pas « l'antre » ; \all{aktiv engagiert} se dit \all{engagiert}, « actuellement » se dit \all{aktuell} (et non \all{aktiv}).
    \item \textbf{Majuscules} : tous les noms prennent la majuscule : \all{die Demokratie, das Wahlrecht, der Bürgermeister}.
    \item \textbf{Prépositions figées} : \all{sich engagieren für + Akk.}, \all{sich einsetzen für + Akk.}, \all{sich freuen auf + Akk.} (futur), \all{sich bedanken bei + Dat.}, \all{diskutieren über + Akk.}
    \item \textbf{Lettre} : après \all{Sehr geehrter Herr Bürgermeister,} on met une virgule et on continue en minuscule ; \all{Mit freundlichen Grüßen} ne prend jamais de virgule.
\end{itemize}
\end{attention}

\newpage

\coursec{Corrigés des exercices}

\begin{corrige}[Exercice 1 — QCM : grammaire et vocabulaire de la citoyenneté]
\begin{enumerate}
    \item \textbf{a} — \all{könnte}. Formation du Konjunktiv II : Präteritum \all{konnte} + umlaut sur la voyelle radicale $\rightarrow$ \all{könnte}. \all{könnt} est une forme du Präsens (2\textsuperscript{e} pers. pl.), \all{konnte} est le Präteritum de l'indicatif.
    \item \textbf{b} — \all{Sehr geehrter Herr Bürgermeister,} est la formule d'appel formelle. \all{Lieber Herr Müller,} convient à une connaissance ; \all{Hallo Bürgermeister,} est familier et incorrect dans un courrier officiel.
    \item \textbf{a} — \all{die Staatsbürgerschaft, -en} = la nationalité, la citoyenneté. Le parlement se dit \all{das Parlament} ; l'élection \all{die Wahl}.
    \item \textbf{a} — \all{diskutieren über + Akkusativ} : \all{Wir diskutieren über die neuen Gesetze.} (\all{die Gesetze} : accusatif pluriel, article inchangé.)
    \item \textbf{b} — \all{Er sagt, er engagiere sich.} Konjunktiv I, 3\textsuperscript{e} pers. sg. : radical + \all{-e} $\rightarrow$ \all{er engagiere}. Le pronom réfléchi \all{sich} est conservé.
\end{enumerate}
\textbf{Points de vigilance} : ne pas confondre Konjunktiv I (discours indirect) et Konjunktiv II (hypothèse) ; dans le QCM, toujours vérifier la construction du verbe avant de répondre.
\end{corrige}

\begin{corrige}[Exercice 2 — Vrai ou Faux ?]
\begin{enumerate}
    \item \all{Falsch. Das Jugendparlament besteht aus jungen Menschen zwischen 14 und 18 Jahren, nicht aus Erwachsenen.}
    \item \all{Richtig. Die Mitglieder werden von anderen Jugendlichen gewählt.}
    \item \all{Falsch. Das Parlament trifft sich einmal im Monat im Rathaus, nicht jede Woche.}
    \item \all{Richtig. Junge Menschen können dort ihre Ideen für die Stadt einbringen.}
\end{enumerate}
\textbf{Points de vigilance} : la justification doit être une phrase complète (verbe conjugué en 2\textsuperscript{e} position) ; repérer les pièges de quantité et de fréquence (\all{nur}, \all{einmal im Monat} vs. \all{jede Woche}).
\end{corrige}

\begin{corrige}[Exercice 3 — Vocabulaire thématique]
\begin{center}
\begin{tabularx}{\linewidth}{|X|X|X|}
\hline
\rowcolor{popblueL} \textbf{Deutsch} & \textbf{Français} & \textbf{Exemple} \\
\hline
\all{die Wahl, die Wahlen} & l'élection & \all{Die Wahl findet am Sonntag statt.} \\
\hline
\all{das Wahlrecht, die Wahlrechte} & le droit de vote & \all{Das Wahlrecht bekommt man mit 18 Jahren.} \\
\hline
\all{der Abgeordnete, die Abgeordneten} & le député & \all{Der Abgeordnete hält eine Rede.} \\
\hline
\all{das Engagement, die Engagements} & l'engagement (civique) & \all{Ihr Engagement im Verein ist vorbildlich.} \\
\hline
\all{die Versammlung, die Versammlungen} & la réunion, l'assemblée & \all{Die Versammlung beginnt um 18 Uhr.} \\
\hline
\all{der Stimmzettel, die Stimmzettel} & le bulletin de vote & \all{Der Stimmzettel kommt in die Urne.} \\
\hline
\all{der Wahlkreis, die Wahlkreise} & la circonscription électorale & \all{Unser Wahlkreis hat einen neuen Kandidaten.} \\
\hline
\end{tabularx}
\end{center}
\textbf{Points de vigilance} : toujours apprendre le nom avec son article et son pluriel ; \all{der Abgeordnete} est un adjectif substantivé (déclinaison faible : \all{der Abgeordnete, den Abgeordneten, die Abgeordneten}).
\end{corrige}

\begin{corrige}[Exercice 4 — Conjugaison à compléter]
\begin{enumerate}
    \item \all{Wenn ich Zeit hätte, würde ich mich in einem Verein engagieren.} (Konjunktiv II de \all{haben} : \all{hätte} ; verbe pronominal séparé : \all{würde} + infinitif \all{engagieren} à la fin, \all{mich} après le sujet.)
    \item \all{Der Bürgermeister eröffnete gestern das neue Jugendzentrum.} (verbe faible au Präteritum : radical + \all{-te}.)
    \item \all{Wir haben schon lange über das Wahlalter diskutiert.} (Perfekt : auxiliaire \all{haben} en 2\textsuperscript{e} position, participe \all{diskutiert} à la fin.)
    \item \all{Könnten Sie mir bitte helfen, den Brief zu formulieren?} (Konjunktiv II de \all{können} : \all{könnten} ; question directe : verbe en 1\textsuperscript{re} position.)
    \item \all{Die Schüler müssen zu der Versammlung gehen.} (modal conjugué en 2\textsuperscript{e} position, infinitif \all{gehen} à la fin.)
    \item \all{Ich werde morgen den Brief an die Schulleitung schreiben.} (Futur I : \all{werden} conjugué + infinitif à la fin.)
\end{enumerate}
\textbf{Points de vigilance} : avec \all{würde}, \all{werden} et les modaux, l'infinitif se place toujours en fin de phrase ; \all{hätte} garde l'umlaut.
\end{corrige}

\begin{corrige}[Exercice 5 — Traduction : phrases courtes]
\textbf{Thème (Français $\rightarrow$ Allemand) :}
\begin{enumerate}
    \item \all{Die Jugendlichen engagieren sich für die Umwelt.} (\all{sich engagieren für + Akkusativ} ; pronom réfléchi \all{sich} après le verbe conjugué.)
    \item \all{Haben Sie schon einmal bei einer Wahl gewählt?} (Perfekt en question directe : auxiliaire en 1\textsuperscript{re} position ; \all{bei einer Wahl} = à une élection.)
    \item \all{Wenn ich Bürgermeister wäre, würde ich ein Jugendzentrum bauen.} (Konjunktiv II : \all{wäre} ; après la subordonnée en \all{wenn}, le verbe de la principale vient en 1\textsuperscript{re} position : \all{würde ich}.)
    \item \all{Man sagt, dass das Parlament sich morgen versammle.} (Konjunktiv I : \all{versammle} ; verbe à la fin de la subordonnée en \all{dass}.)
    \item \all{Ich habe ihm gestern einen formellen Brief geschrieben.} (\all{jemandem (Dat.) einen Brief (Akk.) schreiben} : \all{ihm} ; participe \all{geschrieben} à la fin.)
\end{enumerate}
\textbf{Version (Allemand $\rightarrow$ Français) :}
\begin{enumerate}
    \item L'âge du droit de vote a été abaissé à 16 ans. (passif Präteritum : \all{wurde gesenkt}.)
    \item Nous devons nous engager pour / défendre nos droits. (\all{sich einsetzen für} = se battre pour.)
    \item Le journaliste demande au candidat s'il veut être élu. (\all{ob} = si ; \all{gewählt werden} = être élu, infinitif passif.)
    \item Après la réunion, nous avons déposé une motion. (\all{einen Antrag stellen} = déposer une motion / une demande officielle.)
    \item Accepterais-tu de te porter candidat au poste de délégué de classe ? (\all{sich aufstellen lassen} = se porter candidat.)
\end{enumerate}
\textbf{Points de vigilance} : vérifier la place du verbe, les majuscules aux noms, le cas après la préposition ; en version, rendre le passif allemand par « être + participe » ou « on ».
\end{corrige}

\begin{corrige}[Exercice 6 — Transformation : discours indirect]
\all{Der Reporter fragt den Kandidaten, warum er sich in der Lokalpolitik engagiere. Der Kandidat antwortet, dass er die Situation für junge Menschen verbessern möchte. Der Reporter fragt weiter, ob er konkrete Pläne habe. Der Kandidat bejaht und sagt, dass er einen Jugendrat gründen wolle. Auf die Frage, ob er bei der nächsten Wahl antreten werde, erklärt der Kandidat, dass er das noch nicht entschieden habe.}

\textbf{Méthode appliquée} :
\begin{itemize}
    \item Question avec mot interrogatif (\all{warum}) $\rightarrow$ subordonnée avec le même mot : \all{..., warum er sich ... engagiere.}
    \item Question sans mot interrogatif (\all{Haben Sie...?}, \all{Werden Sie...?}) $\rightarrow$ subordonnée avec \all{ob}.
    \item Konjunktiv I, 3\textsuperscript{e} pers. sg. : \all{er engagiere, er habe, er wolle, er werde} ; quand le Konjunktiv I ressemble à l'indicatif, on peut utiliser \all{würde} + infinitif.
    \item \all{Ich möchte} $\rightarrow$ \all{er möchte} (Konjunktiv I identique au Präteritum, forme acceptée) ; \all{ich will} $\rightarrow$ \all{er wolle} ; \all{ich habe entschieden} $\rightarrow$ \all{er habe entschieden} (Konjunktiv I au passé : \all{habe} + participe).
    \item Pronoms adaptés : \all{ich} $\rightarrow$ \all{er} ; verbe conjugué toujours en fin de subordonnée.
\end{itemize}
\end{corrige}

\begin{corrige}[Exercice 7 — Texte à trous : lettre formelle]
\begin{enumerate}
    \item \all{bei} : \all{sich bei jemandem bedanken} (datif) $\rightarrow$ \all{hiermit möchte ich mich bei Ihnen für das Gespräch neulich bedanken.}
    \item \all{das} : pronom relatif accusatif neutre, antécédent \all{das Treffen} $\rightarrow$ \all{Bei dem Treffen, das wir hatten, ...}
    \item \all{bäte} : Konjunktiv II de \all{bitten} (Präteritum \all{bat} + umlaut) $\rightarrow$ \all{Ich bäte Sie, die Möglichkeit einer Partnerschaft zu prüfen.} La forme \all{würde bitten} est acceptée.
    \item \all{haben sich interessiert} : Perfekt de \all{sich interessieren} $\rightarrow$ \all{Viele Schüler haben sich schon lange für Deutschland interessiert.}
    \item \all{stattfände} : Konjunktiv II de \all{stattfinden} (verbe fort séparable : \all{fand statt} $\rightarrow$ \all{fände statt}) $\rightarrow$ \all{Wenn der Austausch stattfände, ...}
    \item \all{würden ... profitieren} : \all{..., würden beide Seiten davon profitieren.} (après la subordonnée, verbe en 1\textsuperscript{re} position de la principale.)
    \item \all{Auf} : \all{sich freuen auf + Akkusativ} (joie à venir) $\rightarrow$ \all{Auf eine positive Antwort würde ich mich sehr freuen.}
\end{enumerate}
\textbf{Points de vigilance} : distinguer \all{sich freuen auf} (avenir) et \all{sich freuen über} (présent/passé) ; les verbes forts au Konjunktiv II prennent l'umlaut quand c'est possible (\all{bat $\rightarrow$ bäte}, \all{fand $\rightarrow$ fände}).
\end{corrige}

\begin{corrige}[Exercice 8 — Appariements]
\begin{itemize}
    \item \textbf{1--B} : \all{Meine sehr geehrten Damen und Herren, liebe Mitschülerinnen und Mitschüler, ...} — salutation solennelle typique de l'ouverture d'un discours officiel.
    \item \textbf{2--A} : \all{Könnten Sie mir bitte sagen, was Ihre Meinung zum Wahlalter ist?} — question polie au Konjunktiv II avec vouvoiement, typique de l'interview.
    \item \textbf{3--C} : \all{Ich bin ganz Ihrer Meinung, denn die Jugend ist die Zukunft.} — accord argumenté avec \all{denn} (car).
    \item \textbf{4--E} : \all{Ich bedanke mich für das Vertrauen, muss aber ablehnen.} — remerciement puis refus poli.
    \item \textbf{5--D} : \all{Das ist ein interessanter Punkt, aber ich sehe das anders.} — concession puis désaccord nuancé.
    \item \textbf{6--F} : \all{Ich freue mich auf Ihre Antwort und verbleibe mit freundlichen Grüßen.} — formule de clôture de lettre formelle.
\end{itemize}
\textbf{Points de vigilance} : repérer les marqueurs de registre (\all{Sie} / \all{Ihrer} = formel) ; \all{verbleibe mit freundlichen Grüßen} est une formule de fin très soutenue.
\end{corrige}

\begin{corrige}[Exercice 9 — Compréhension et analyse linguistique (type Bac)]
\textbf{Compréhension (en allemand) :}
\begin{enumerate}
    \item \all{Die Mitglieder sind meist zwischen 14 und 21 Jahre alt. Sie werden für ein oder zwei Jahre gewählt.}
    \item \all{Das Jugendparlament in Bonn hat sich erfolgreich für einen neuen Skatepark und für kostenloses WLAN in öffentlichen Parks eingesetzt.}
    \item \all{Die Kritiker bemängeln, dass die Parlamente oft nur beratende Funktion haben und keine echten Entscheidungsbefugnisse besitzen.}
\end{enumerate}
\textbf{Questions de langue (en français) :}
\begin{enumerate}
    \setcounter{enumi}{3}
    \item Deux verbes au Perfekt : \all{hat ... eingesetzt} (infinitif \all{sich einsetzen} : s'engager, se battre pour) et \all{hat ... erzielt} (dans la question 2 ; infinitif \all{erzielen} : obtenir, remporter). On accepte aussi \all{hat ... angesprochen} si l'élève reformule la phrase du texte (\all{Es spricht ... an} est au Präsens dans le texte ; ne pas l'accepter sans transformation). Attention : \all{werden gehört} est un présent passif, pas un Perfekt.
    \item Relatif à l'accusatif : \all{einen Haushaltsplan, den die Mitglieder selbst aufstellen} — \all{den} : masculin accusatif, antécédent \all{der Haushaltsplan}. Relatif au datif : \all{Jugendliche, denen Politik sonst fern erscheint} — \all{denen} : pluriel datif, antécédent \all{die Jugendlichen}.
    \item « On nous écoute enfin », dit la porte-parole Lena Meyer, 17 ans. « La politique n'est pas seulement une affaire de vieux. » (\all{wir werden gehört} : passif présent, rendu par « on » ; \all{etwas für Alte} : une chose pour les personnes âgées.)
    \item \all{Wenn man früh lernen würde, Verantwortung zu übernehmen, (dann) würde das die Demokratie von morgen stärken.} Autre possibilité : \all{Wenn man früh Verantwortung übernähme, stärkte das die Demokratie von morgen.} (\all{lernen} est faible : on préfère \all{würde lernen} ; \all{übernehmen} est fort : \all{übernähme} possible.)
\end{enumerate}
\textbf{Barème indicatif (20 pts)} : compréhension 3 $\times$ 2 pts = 6 ; Perfekt 2 pts ; relatifs 2 $\times$ 2 pts = 4 ; traduction 4 pts ; reformulation au Konjunktiv II 4 pts.
\textbf{Points de vigilance} : répondre par des phrases complètes en allemand ; ne pas confondre passif (\all{werden} + participe) et Perfekt (\all{haben/sein} + participe) ; le pronom relatif prend le genre et le nombre de l'antécédent, mais le cas de sa fonction dans la subordonnée.
\end{corrige}

\begin{corrige}[Exercice 10 — Traduction thème et version (type Bac)]
\textbf{Partie A : Thème}
\begin{enumerate}
    \item \all{Die aktive Bürgerschaft stärkt die Demokratie in unserem Land.} (accepté aussi : \all{Das aktive bürgerschaftliche Engagement stärkt die Demokratie in unserem Land.})
    \item \all{Seit dem letzten Jahr haben die Jugendlichen ab 16 Jahren das Wahlrecht bei den Kommunalwahlen.} (\all{seit} + datif : \all{seit dem letzten Jahr} ; \all{bei} + datif.)
    \item \all{Der Bürgermeister hat die Schüler gefragt, ob sie ein Projekt für die Umwelt organisieren wollten.} (\all{jemanden fragen, ob ...} ; discours indirect au passé : Konjunktiv II \all{wollten}.)
    \item \all{Wenn die Politiker die Jugend mehr hören würden, würden viele Probleme gelöst.} (ou \all{..., wären viele Probleme gelöst.} Passif au Konjunktiv II : \all{würden gelöst} ; participe à la fin.)
    \item \all{Wir haben beschlossen, dem Bildungsminister einen Brief zu schreiben, um einen Schüleraustausch vorzuschlagen.} (\all{beschließen, etwas zu tun} ; \all{um ... vorzuschlagen} : infinitif avec \all{zu} inséré dans la particule séparable.)
\end{enumerate}
\textbf{Partie B : Version}

Le débat sur l'abaissement de l'âge du droit de vote à 16 ans prend de l'ampleur. Les partisans font valoir que les jeunes sont déjà pénalement responsables à 16 ans et paient des impôts lorsqu'ils travaillent. Il serait donc juste de leur accorder aussi le droit de vote. Des études montrent que la participation électorale chez les primo-votants est souvent plus élevée que dans les tranches d'âge plus âgées. Les opposants craignent en revanche que les jeunes de 16 ans soient encore trop facilement influençables et ne puissent pas saisir pleinement la portée des décisions politiques. Dans plusieurs Länder, le vote à partir de 16 ans est déjà une réalité aux élections communales et régionales. Au niveau fédéral, en revanche, l'âge minimum de 18 ans reste en vigueur pour l'élection du Bundestag.

\textbf{Points de langue} : \all{an Fahrt gewinnen} = prendre de l'ampleur ; \all{strafmündig} = pénalement responsable ; \all{sei / seien / könnten} : Konjunktiv I du discours indirect, rendu par le conditionnel ou le subjonctif ; \all{die Wahlbeteiligung} = la participation électorale ; \all{der Erstwähler} = le primo-votant ; \all{die Tragweite} = la portée.

\textbf{Barème indicatif (20 pts)} : thème 5 $\times$ 2 pts = 10 (1 pt pour le sens, 1 pt pour la correction grammaticale) ; version 10 pts (fidélité 5 pts, qualité du français 3 pts, vocabulaire 2 pts).
\end{corrige}

\begin{corrige}[Exercice 11 — Production écrite guidée (type Bac)]
\textbf{Barème indicatif (20 pts)} : respect du genre et des conventions (en-tête, appel, clôture / accroche, appel au vote) 4 pts ; respect de la consigne et des contenus exigés 4 pts ; correction grammaticale (place du verbe, cas, temps) 6 pts ; richesse du vocabulaire thématique 3 pts ; cohérence et registre formel 3 pts.

\textbf{Proposition de rédaction modèle — Sujet 1 (lettre formelle) :}

\all{Moussa Njoya, Klassensprecher der Klasse 1\textsuperscript{ère} A, Lycée de Yaoundé, Yaoundé, Kamerun}

\all{An den Bürgermeister der Stadt Bonn, Rathaus, Bonn, Deutschland}

\all{Yaoundé, den 12. März 2025}

\all{Betreff: Partnerschaft für einen Schüleraustausch}

\all{Sehr geehrter Herr Bürgermeister,}

\all{als Klassensprecher der Klasse 1\textsuperscript{ère} A des Lycée de Yaoundé schreibe ich Ihnen, um eine Partnerschaft zwischen unserer Schule und einem Gymnasium in Bonn vorzuschlagen. Unsere Klasse besteht aus 45 Schülern, die seit drei Jahren Deutsch lernen und sich sehr für die deutsche Kultur und das demokratische Leben in Ihrem Land interessieren.}

\all{Ein Austausch von zwei Wochen im nächsten Schuljahr mit etwa zwanzig Schülern würde beiden Seiten viele Vorteile bringen: Die Jugendlichen könnten ihre Sprachkenntnisse verbessern, eine andere Kultur entdecken und als junge Bürger Verantwortung übernehmen. Wir könnten zum Beispiel gemeinsam ein Projekt zum Thema Umwelt oder Jugendengagement organisieren.}

\all{Über eine positive Antwort würde ich mich sehr freuen.}

\all{Mit freundlichen Grüßen}

\all{Moussa Njoya}

\textbf{Proposition de rédaction modèle — Sujet 2 (petit discours) :}

\all{Meine sehr geehrten Damen und Herren, liebe Mitschülerinnen und Mitschüler,}

\all{ich heiße Amina Bekono und ich möchte eure Schulsprecherin werden. Warum ich? Weil ich mich seit zwei Jahren im Sportverein und in der Umweltgruppe engagiere und weil ich eure Ideen ernst nehme.}

\all{Wenn ihr mich wählt, werde ich drei konkrete Dinge tun. Erstens will ich mit der Schulleitung über eine bessere Kantine sprechen. Zweitens werde ich neue Clubs gründen, zum Beispiel einen Deutsch-Club und einen Fußball-Club für Mädchen. Drittens möchte ich jeden Monat eine Versammlung organisieren, damit jeder Schüler und jede Schülerin gehört wird.}

\all{Politik ist nicht nur etwas für Erwachsene — sie beginnt hier, in unserer Schule. Stimmt für mich, und wir werden gemeinsam unsere Schule verändern!}

\all{Ich danke euch für eure Aufmerksamkeit.}

\textbf{Points de vigilance} :
\begin{itemize}
    \item Lettre : virgule après l'appel puis minuscule ; \all{Mit freundlichen Grüßen} sans virgule ; \all{Betreff:} sans article ; date avec \all{den} + date ordinale (\all{den 12. März 2025}).
    \item Discours : alterner \all{ich werde / ich will / ich möchte} ; annoncer le plan avec \all{erstens, zweitens, drittens} ; finir par un appel au vote clair.
    \item Grammaire : verbe conjugué en 2\textsuperscript{e} position, infinitif à la fin avec \all{würde/werden}, majuscules à tous les noms, \all{sich engagieren für + Akk.}, \all{sich freuen über/auf + Akk.}
    \item Relire en comptant les mots (120--150) et en vérifiant chaque verbe.
\end{itemize}
\end{corrige}

\newpage

\coursec{Fiche bilan}

\begin{popfigure}
\begin{tikzpicture}[
  mindmap,
  grow cyclic,
  every node/.style={concept, font=\small, align=center},
  root concept/.append style={concept color=popblue, font=\bfseries\small, text=white, minimum size=3.2cm},
  level 1 concept/.append style={level distance=3.6cm, sibling angle=60, font=\footnotesize},
]
\node[root concept, align=center]{Citoyenneté\\ Traduire \&\\ produire}
  child[concept color=poporange]{ node[align=center]{Lexique\\ der Bürger\\ die Wahl\\ die Pflicht} }
  child[concept color=popgreen]{ node[align=center]{L'interview\\ Fragen stellen\\ W-Fragen} }
  child[concept color=poppurple]{ node[align=center]{La lettre\\ Sehr geehrte...\\ Mit freundlichen\\ Grüßen} }
  child[concept color=poppink]{ node[align=center]{Le discours\\ Liebe...\\ wir müssen\\ wir können} }
  child[concept color=popgold]{ node[align=center]{Grammaire\\ verbe en 2\textsuperscript{e} position\\ subordonnées\\ Perfekt/Präteritum} }
  child[concept color=popblue!60]{ node[align=center]{Thème \& version\\ français $\rightarrow$ allemand\\ allemand $\rightarrow$ français} };
\end{tikzpicture}
\legende{Carte mentale de la leçon : produire et traduire des textes sur la citoyenneté}
\end{popfigure}

\begin{bilan}
\textbf{1. Lexique clé de la citoyenneté (à connaître par cœur)}
\begin{itemize}
  \item \all{der Bürger, - / die Bürgerin, -nen} : le citoyen / la citoyenne
  \item \all{die Staatsangehörigkeit, -en} : la nationalité
  \item \all{das Recht, -e / die Pflicht, -en} : le droit / le devoir
  \item \all{die Wahl, -en / wählen} : l'élection / élire, voter
  \item \all{die Freiheit, -en / die Verantwortung, -en} : la liberté / la responsabilité
  \item \all{sich engagieren (für + Akk.)} : s'engager (pour)
  \item \all{der Freiwillige, -n / die Freiwilligenarbeit} : le bénévole / le bénévolat
  \item \all{die Gemeinde, -n / die Demokratie, -n} : la commune / la démocratie
  \item \all{respektieren / helfen (dat.)} : respecter / aider
  \item \all{das Interview, -s / der Journalist, -en} : l'interview / le journaliste
\end{itemize}

\textbf{2. Règles de grammaire à connaître par cœur}
\begin{center}
\begin{tabularx}{\linewidth}{|X|X|}
\hline
\rowcolor{popblueL} \textbf{Règle} & \textbf{Exemple} \\
\hline
Le verbe conjugué est toujours en 2\textsuperscript{e} position dans la phrase déclarative. & \all{Heute engagieren sich viele Jugendliche.} \\
\hline
Dans une subordonnée (\all{weil, dass, wenn, obwohl}), le verbe conjugué va à la fin. & \all{..., weil sie der Gemeinde helfen wollen.} \\
\hline
\all{W-Fragen} (questions ouvertes) : mot interrogatif + verbe + sujet. & \all{Warum engagieren Sie sich?} \\
\hline
Questions fermées : verbe en 1\textsuperscript{re} position. & \all{Haben Sie Zeit für ein Interview?} \\
\hline
Perfekt : \all{haben/sein} + participe II ; verbes de mouvement avec \all{sein}. & \all{Ich habe geholfen. / Wir sind gegangen.} \\
\hline
Verbes à particule séparable : la particule va en fin de phrase au Präsens. & \all{Der Bürgermeister kommt um 10 Uhr an.} \\
\hline
Majuscule à tous les noms ; \all{Sie} (vouvoiement) avec majuscule. & \all{die Wahl, der Bürger, Sie} \\
\hline
\end{tabularx}
\end{center}

\textbf{3. Modèles de production (structures fixes)}
\begin{itemize}
  \item \textbf{Interview :} salutation + présentation + 4 à 6 questions (\all{Wie, Warum, Seit wann, Was}) + remerciement. Ex. : \all{Guten Tag, Herr Nkoulou! Darf ich Ihnen einige Fragen stellen?}
  \item \textbf{Lettre :} lieu et date à droite (\all{Yaoundé, den 12. Mai}) ; \all{Sehr geehrte Frau... / Lieber...} ; corps en paragraphes ; \all{Mit freundlichen Grüßen} (formel) ou \all{Viele Grüße} (amical).
  \item \textbf{Petit discours :} \all{Liebe Mitbürger!} + annonce du sujet + arguments avec \all{wir müssen, wir können, wir sollen} + appel à l'action + remerciement : \all{Ich danke Ihnen für Ihre Aufmerksamkeit.}
\end{itemize}

\textbf{4. Méthodes express de traduction}
\begin{itemize}
  \item \textbf{Version (allemand $\rightarrow$ français) :} 1) lire tout le texte ; 2) repérer sujet, verbe, compléments ; 3) traduire les structures avant les mots ; 4) rédiger un français naturel ; 5) relire.
  \item \textbf{Thème (français $\rightarrow$ allemand) :} 1) identifier le temps ; 2) placer le verbe en 2\textsuperscript{e} position (ou à la fin en subordonnée) ; 3) vérifier genre, cas, majuscules, particules séparables ; 4) relire chaque phrase.
\end{itemize}
\end{bilan}

\begin{aretenir}[Checklist avant le Bac]
Avant de rendre ma copie, je vérifie :
\begin{itemize}
  \item $\square$ Le verbe conjugué est en 2\textsuperscript{e} position dans chaque phrase déclarative.
  \item $\square$ Dans les subordonnées (\all{weil, dass, wenn}), le verbe est à la fin.
  \item $\square$ Tous les noms allemands prennent une majuscule.
  \item $\square$ Chaque nom est accompagné du bon article (\all{der, die, das}) et du bon cas.
  \item $\square$ Les particules séparables sont bien détachées (\all{Ich stehe um 6 Uhr auf.}).
  \item $\square$ Le Perfekt est formé avec le bon auxiliaire (\all{haben} ou \all{sein}).
  \item $\square$ Les Umlaute (\all{ä, ö, ü}) et le \all{ß} sont correctement écrits.
  \item $\square$ Ma lettre contient : date, formule d'appel, corps, formule de politesse, signature.
  \item $\square$ Mon interview contient des questions variées (\all{W-Fragen} et questions fermées).
  \item $\square$ Mon discours utilise \all{wir müssen / wir können} et s'adresse au public.
  \item $\square$ J'ai traduit les idées, pas mot à mot, et ma traduction française est naturelle.
  \item $\square$ J'ai relu ma copie au moins une fois pour corriger l'orthographe et l'ordre des mots.
\end{itemize}
\end{aretenir}

\findecours

\end{document}
